\documentclass[11pt,reqno]{amsart}

% ================================================================
%  Packages
% ================================================================
\usepackage{amsmath,amssymb,amsthm}
\usepackage{mathrsfs}
\usepackage{enumerate}
\usepackage[shortlabels]{enumitem}
\usepackage{booktabs}
\usepackage{array}
\usepackage{microtype}
\usepackage[dvipsnames]{xcolor}
\usepackage[colorlinks=true,linkcolor=blue!60!black,citecolor=green!40!black,urlcolor=blue!60!black]{hyperref}
\usepackage[all]{xy}

% ================================================================
%  Theorem environments
% ================================================================
\newtheorem{theorem}{Theorem}[section]
\newtheorem{proposition}[theorem]{Proposition}
\newtheorem{lemma}[theorem]{Lemma}
\newtheorem{corollary}[theorem]{Corollary}
\newtheorem{conjecture}[theorem]{Conjecture}
\theoremstyle{definition}
\newtheorem{definition}[theorem]{Definition}
\newtheorem{construction}[theorem]{Construction}
\newtheorem{example}[theorem]{Example}
\newtheorem{principle}[theorem]{Principle}
\theoremstyle{remark}
\newtheorem{remark}[theorem]{Remark}

% ================================================================
%  Macros (providecommand: safe against collision with any input)
% ================================================================
\providecommand{\cA}{\mathcal{A}}
\providecommand{\cB}{\mathcal{B}}
\providecommand{\cC}{\mathcal{C}}
\providecommand{\cD}{\mathcal{D}}
\providecommand{\cE}{\mathcal{E}}
\providecommand{\cF}{\mathcal{F}}
\providecommand{\cG}{\mathcal{G}}
\providecommand{\cH}{\mathcal{H}}
\providecommand{\cI}{\mathcal{I}}
\providecommand{\cJ}{\mathcal{J}}
\providecommand{\cK}{\mathcal{K}}
\providecommand{\cL}{\mathcal{L}}
\providecommand{\cM}{\mathcal{M}}
\providecommand{\cN}{\mathcal{N}}
\providecommand{\cO}{\mathcal{O}}
\providecommand{\cP}{\mathcal{P}}
\providecommand{\cQ}{\mathcal{Q}}
\providecommand{\cR}{\mathcal{R}}
\providecommand{\cS}{\mathcal{S}}
\providecommand{\cT}{\mathcal{T}}
\providecommand{\cU}{\mathcal{U}}
\providecommand{\cV}{\mathcal{V}}
\providecommand{\cW}{\mathcal{W}}
\providecommand{\cX}{\mathcal{X}}
\providecommand{\cY}{\mathcal{Y}}
\providecommand{\cZ}{\mathcal{Z}}
\providecommand{\barB}{\bar{B}}
\providecommand{\barC}{\overline{C}}
\providecommand{\Defcyc}{\mathrm{Def}_{\mathrm{cyc}}}
\providecommand{\Ran}{\mathrm{Ran}}
\providecommand{\MC}{\mathrm{MC}}
\providecommand{\Sym}{\mathrm{Sym}}
\providecommand{\Hom}{\mathrm{Hom}}
\providecommand{\Ext}{\mathrm{Ext}}
\providecommand{\Res}{\mathrm{Res}}
\providecommand{\Aut}{\mathrm{Aut}}
\providecommand{\FM}{\overline{C}}
\providecommand{\fg}{\mathfrak{g}}
\providecommand{\fh}{\mathfrak{h}}
\providecommand{\gmod}{\mathfrak{g}^{\mathrm{mod}}_\cA}
\providecommand{\gAmod}{\mathfrak{g}^{\mathrm{mod}}_\cA}
\providecommand{\Sh}{\mathrm{Sh}}
\providecommand{\OPE}{\mathrm{OPE}}
\providecommand{\Vir}{\mathrm{Vir}}
\providecommand{\Ahat}{\widehat{A}}
\providecommand{\hol}{\mathrm{hol}}
\providecommand{\bC}{\mathbb{C}}
\providecommand{\bP}{\mathbb{P}}
\providecommand{\bA}{\mathbb{A}}
\providecommand{\bZ}{\mathbb{Z}}
\providecommand{\bR}{\mathbb{R}}
\providecommand{\bQ}{\mathbb{Q}}
\providecommand{\Bbbk}{\mathbb{k}}
\providecommand{\coll}{\mathrm{coll}}
\providecommand{\FP}{\mathrm{FP}}
\providecommand{\cross}{\mathrm{cross}}
% Operad/algebra macros
\providecommand{\Ass}{\mathrm{Ass}}
\providecommand{\Com}{\mathrm{Com}}
\providecommand{\SCchtop}{\mathrm{SC}^{\mathrm{ch,top}}}
\providecommand{\Convinf}{\mathrm{Conv}_\infty}
\providecommand{\Convstr}{\mathrm{Conv}_{\mathrm{str}}}
\providecommand{\Ainf}{A_\infty}
\providecommand{\ChirHoch}{\mathrm{ChirHoch}}
\providecommand{\Defcyc}{\mathrm{Def}_{\mathrm{cyc}}}
\providecommand{\orline}[1]{\mathfrak{o}_{#1}}
\providecommand{\fsl}{\mathfrak{sl}}
\providecommand{\HH}{\mathrm{HH}}
\providecommand{\Fun}{\mathrm{Fun}}
\providecommand{\Op}{\mathrm{Op}}
\providecommand{\Rep}{\mathrm{Rep}}
\providecommand{\End}{\mathrm{End}}
\providecommand{\Conf}{\mathrm{Conf}}
% Subscript/superscript decorators
\providecommand{\coll}{\mathrm{coll}}
\providecommand{\dg}{\mathrm{dg}}
\providecommand{\FP}{\mathrm{FP}}
\providecommand{\KZ}{\mathrm{KZ}}
\providecommand{\conn}{\mathrm{conn}}
\providecommand{\bF}{\mathbb{F}}
\providecommand{\der}{\mathrm{der}}
\providecommand{\ch}{\mathrm{ch}}
\providecommand{\op}{\mathrm{op}}
\providecommand{\oc}{\mathrm{oc}}
\providecommand{\SC}{\mathrm{SC}}
\providecommand{\Id}{\mathrm{Id}}
\providecommand{\id}{\mathrm{id}}
\providecommand{\Reg}{\mathrm{Reg}}
\providecommand{\GM}{\mathrm{GM}}
\providecommand{\co}{\mathrm{co}}
\providecommand{\contact}{\mathrm{contact}}
\providecommand{\cross}{\mathrm{cross}}
\providecommand{\grav}{\mathrm{grav}}
\providecommand{\sep}{\mathrm{sep}}
\providecommand{\nsep}{\mathrm{nsep}}
\providecommand{\DS}{\mathrm{DS}}
\providecommand{\BV}{\mathrm{BV}}
\providecommand{\cyc}{\mathrm{cyc}}
\providecommand{\mod}{\mathrm{mod}}
\providecommand{\fib}{\mathrm{fib}}
\providecommand{\Fl}{\mathrm{Fl}}
\providecommand{\sgn}{\mathrm{sgn}}
\providecommand{\Coder}{\mathrm{Coder}}
\providecommand{\Ob}{\mathrm{Ob}}
\providecommand{\Obs}{\mathrm{Obs}}
\providecommand{\clutch}{\mathrm{clutch}}
\providecommand{\Ch}{\mathrm{Ch}}
\providecommand{\red}{\mathrm{red}}
\providecommand{\gh}{\mathrm{gh}}
\providecommand{\fo}{\mathfrak{o}}
\providecommand{\CompCl}{\mathrm{CompCl}}
\providecommand{\Tr}{\mathrm{Tr}}
\providecommand{\Conv}{\mathrm{Conv}}

\DeclareMathOperator{\gr}{gr}
\DeclareMathOperator{\rk}{rk}
\DeclareMathOperator{\depth}{depth}
\DeclareMathOperator{\Spec}{Spec}

\numberwithin{equation}{section}

% ================================================================
% Graceful fallback for cross-references to theorems/propositions
% that live in the full monograph but are cited in this standalone
% survey.  Undefined labels resolve to "full monograph" rather than
% the default LaTeX "??" placeholder.
% ================================================================
\makeatletter
\let\svy@oldref\ref
\renewcommand{\ref}[1]{%
  \@ifundefined{r@#1}{\textit{full monograph}}{\svy@oldref{#1}}%
}
\makeatother

% ================================================================
\begin{document}

\title[Modular Koszul Duality: A Survey]
{Modular Koszul Duality and the Shadow Obstruction Tower:\\
A Survey}

\author{Raeez Lorgat}
\address{Perimeter Institute for Theoretical Physics,
  Waterloo, ON N2L 2Y5, Canada}
\email{rlorgat@perimeterinstitute.ca}

\date{\today}

\begin{abstract}
This survey charts the programme of modular Koszul duality for chiral
algebras on algebraic curves.  The central object is the ordered bar
complex $\barB^{\mathrm{ord}}(\cA)=T^c(s^{-1}\bar\cA)$ with
deconcatenation coproduct, an $E_1$-coalgebra built on ordered
configurations; the symmetric bar $\barB^{\Sigma}(\cA)$ is its
$\Sigma_n$-coinvariant shadow.  The bar differential on the
Fulton--MacPherson compactification squares to zero by Arnold plus
Borcherds at genus zero, and acquires curvature
$\kappa(\cA)\cdot\omega_g$ from the Hodge bundle at higher genus.  A
single universal Maurer--Cartan element $\Theta_{\cA}$ controls the
genus tower.  Five theorems (A--D, H) organise the invariants:
bar--cobar inversion, the Verdier intertwining, complementarity, the
modular characteristic, and the polynomial bound on chiral Hochschild
cohomology.  The shadow obstruction tower partitions the standard
landscape into four depth classes (Gaussian, Lie, contact, mixed),
each realised by an explicit family.  We survey the $E_1/E_\infty$
hierarchy, the seven faces of the collision residue, the arithmetic
of the shadow $L$-function, the open/closed bridge to 3D
holomorphic--topological quantum field theory, modular PVA
quantization, holographic reconstruction, and the outstanding
frontier conjectures.
\end{abstract}

\subjclass[2020]{17B69, 14H10, 18M70, 81T40}

\keywords{Chiral algebras, bar construction, Koszul duality,
moduli of curves, vertex algebras, Maurer--Cartan,
W-algebras, factorization algebras, Yangian, shadow tower}

\maketitle

\tableofcontents

\bigskip

% ====================================================================
% TRACK A: sections 1--6 (bar construction, curvature, modular homotopy,
% universal Maurer--Cartan, tangent complex, shadow obstruction tower)
% ====================================================================

\section*{1.\quad The ordered bar complex}
% ====================================================================

\medskip

\noindent
The primitive object of this programme is the \emph{ordered bar
complex} $\barB^{\mathrm{ord}}(\cA)=T^c(s^{-1}\bar\cA)$, an
$E_1$-coalgebra with deconcatenation coproduct, built on ordered
configurations of points on an algebraic curve.  The symmetric bar
$\barB^{\Sigma}(\cA)$, and with it every statement about the modular
characteristic~$\kappa$ and the shadow obstruction tower, is the
$\Sigma_n$-coinvariant shadow obtained by averaging.  The ordered bar
carries the full $R$-matrix, the Yangian, the braiding; averaging
collapses this profile to a scalar.  What survives averaging is the
subject of the five theorems; what is lost is the subject of the
$E_1$ frontier.

The exposition constructs the ordered bar first, computes its
chain-level structure on explicit examples, extracts the
$R$-matrix as the arity-$2$ collision residue, and only then
derives the symmetric bar by projection.  No symmetric-bar result
is stated without first exhibiting the $E_1$ object it descends from.

\subsection*{1.1.\enspace The $1$-form}

Let $X$ be a smooth algebraic curve over~$\mathbf C$.  On the
configuration space $C_2(X)=X\times X\setminus\Delta$, define the
logarithmic $1$-form
\begin{equation}\label{eq:preface-eta}
\eta_{12}\;=\;d\log(z_1-z_2)\;=\;\frac{dz_1-dz_2}{z_1-z_2}.
\end{equation}
This form has a simple pole along~$\Delta$ with residue~$1$.  The
residue depends on no coordinate, no metric, no level.  A double
pole~$dz/(z-w)^2$ transforms under coordinate change and has no
coordinate-independent residue; a regular $1$-form extracts nothing
from the collision.  The logarithmic derivative is the unique $1$-form
on $C_2(X)$ with a first-order pole along~$\Delta$ and conformally
invariant residue.

On $C_3(X)$, the three pullbacks $\eta_{12}, \eta_{23}, \eta_{31}$
satisfy the \emph{Arnold relation}
\begin{equation}\label{eq:preface-arnold}
\eta_{12}\wedge\eta_{23}+\eta_{23}\wedge\eta_{31}+\eta_{31}\wedge\eta_{12}\;=\;0,
\end{equation}
the differential-form image of the partial-fractions identity
$\tfrac{1}{(z_1-z_2)(z_2-z_3)}+\tfrac{1}{(z_2-z_3)(z_3-z_1)}+\tfrac{1}{(z_3-z_1)(z_1-z_2)}=0$.
Arnold proved that~\eqref{eq:preface-arnold}, replicated across all
triples $\{i,j,k\}\subset\{1,\ldots,n\}$, generates the complete ideal
of relations in $H^*(\mathrm{Conf}_n(\mathbf C),\mathbf C)$.

Under $z\mapsto w(z)$ the propagator transforms as
$\eta_{ij}\mapsto\eta_{ij}+d\log\frac{w(z_i)-w(z_j)}{z_i-z_j}$;
near the diagonal this reduces to $d\log w'(z_i)$, the standard
$\mathrm{Diff}(X)$ cocycle.  The propagator is a BRST ghost for
diffeomorphisms.

\subsection*{1.2.\enspace The ordered bar complex on~$X$}

A \emph{chiral algebra}~$\cA$ on a smooth algebraic curve~$X$ is a
$\mathcal D_X$-module equipped with a chiral bracket
$\mu^{\mathrm{ch}}\colon j_*j^*(\cA\boxtimes\cA)\to\Delta_!(\cA)$
extracting the singular part of sections with poles along the diagonal.
In local coordinates this encodes the operator product expansion
$a(z)\,b(w)\sim\sum_{n\ge 0}a_{(n)}b(w)/(z-w)^{n+1}$,
governed by the Borcherds identity coupling all pole orders.

Place sections $a_1,\dots,a_n\in\cA$ at distinct ordered points
$z_1<\cdots<z_n$ on~$X$ and assign to each consecutive pair
the form $\eta_{i,i+1}=d\log(z_i-z_{i+1})$.  The
Fulton--MacPherson compactification
$\overline{\mathrm{FM}}_n(X)$ resolves $\mathrm{Conf}_n(X)$ by
iterated real oriented blowup along all diagonal strata (a blowup,
not a complement); boundary codimension-one faces are indexed by
subsets $S\subseteq\{1,\ldots,n\}$ with $|S|\ge 2$.

The \emph{ordered chiral bar complex}:
\begin{equation}\label{eq:ordered-bar}
\barB^{\mathrm{ord}}_X(\cA)
\;=\;
\bigoplus_{n\ge 1}(s^{-1}\bar\cA)^{\otimes n}\otimes
\Gamma\bigl(\overline{\mathrm{FM}}_n(X),\Omega^{n-1}_{\log}\bigr).
\end{equation}
Here $\bar\cA=\ker(\varepsilon\colon\cA\to\omega_X)$ is the
augmentation ideal, $T^c(V)=\bigoplus_{n\ge 0}V^{\otimes n}$ is the
cofree tensor coalgebra, and $s^{-1}$ is desuspension with
$|s^{-1}a|=|a|-1$ in cohomological convention ($|d|=+1$).

This is the primitive object.  It carries the \emph{deconcatenation
coproduct}
\[
\Delta[s^{-1}a_1|\cdots|s^{-1}a_n]
=\sum_{i=0}^n[s^{-1}a_1|\cdots|s^{-1}a_i]\otimes[s^{-1}a_{i+1}|\cdots|s^{-1}a_n],
\]
splitting an ordered sequence at every position: $n+1$ terms for an
ordered $n$-tuple.  This makes $\barB^{\mathrm{ord}}(\cA)$ an
$E_1$-coalgebra: coassociative, not cocommutative.  The linear
ordering of the points is retained; there is no $\Sigma_n$-quotient.
Its $E_1$-Koszul dual $\cA^!_{\mathrm{line}}$ governs the category
of line operators: for affine Kac--Moody,
$\cA^!_{\mathrm{line}}$ is the Yangian $Y(\fg)$.

The differential decomposes as
$d_{\barB}=d_{\mathrm{int}}+d_{\mathrm{res}}+d_{\mathrm{dR}}$:
the internal differential of~$\cA$, the OPE residue extraction
along codimension-one boundary strata of $\overline{\mathrm{FM}}_n(X)$,
and the de Rham differential on logarithmic forms.

\subsection*{1.3.\enspace $d^2=0$: Arnold kills the form, Borcherds kills the algebra}

Expand $(d_{\mathrm{int}}+d_{\mathrm{res}}+d_{\mathrm{dR}})^2$.  The
diagonal terms: $d_{\mathrm{int}}^2=0$, $d_{\mathrm{dR}}^2=0$; the
cross-terms $[d_{\mathrm{int}},d_{\mathrm{dR}}]=0$ by Leibniz.  The
remaining terms pair into three cancellation mechanisms:

\emph{Disjoint supports}: when $\{i,j\}\cap\{k,l\}=\emptyset$ the
boundary faces are transverse, and the residue operators commute.

\emph{Shared index}: when $\{i,j\}\cap\{k,l\}=\{j=k\}$ the triple
collision $z_i=z_j=z_l$ lies in a codimension-two stratum.  The
form contribution $\eta_{ij}\wedge\eta_{jl}$ paired with its cyclic
partners vanishes by the Arnold relation~\eqref{eq:preface-arnold};
the corresponding algebraic contractions vanish by Borcherds.
Arnold kills the form; Borcherds kills the algebra; both are consumed.

\emph{Same pair}: $d_{\mathrm{res}}^{(ij)}\,d_{\mathrm{res}}^{(ij)}=0$
by the desuspension sign rule.

Result: $d_{\barB}^2=0$ iff the $n$-th products satisfy the Borcherds
identity.  The ordered bar construction accepts exactly chiral algebras.

All of this is genus~$0$.  At genus $g\ge 1$ the propagator acquires
quasi-periodicity from the prime form, and $d_{\mathrm{fib}}^2\neq 0$:
the curvature $d_{\mathrm{fib}}^2=\kappa(\cA)\cdot\omega_g$ is a
single scalar times the Hodge class, governing the entire genus
expansion.

\subsection*{1.4.\enspace The collision residue and the $R$-matrix}

The ordered bar differential extracts OPE data at each collision
stratum.  At tensor degree~$2$ the extraction is the \emph{collision
residue}: the residue of the OPE integrated against the propagator
$\eta_{12}=d\log(z_1-z_2)$.  Since $d\log$ has a simple pole, it
absorbs one pole order from the OPE:
\begin{equation}\label{eq:pole-absorption}
\text{pole order of }r(z) \;=\; \text{pole order of OPE} - 1.
\end{equation}
The result is the \emph{classical $r$-matrix}
\[
r(z)\;=\;\operatorname{Res}^{\mathrm{coll}}_{0,2}
\bigl(\Theta^{E_1}_{\cA}\bigr),
\]
a matrix-valued rational function encoding the full binary OPE data
of the ordered bar.  The Yang--Baxter equation
$[r_{12},r_{13}]+[r_{12},r_{23}]+[r_{13},r_{23}]=0$ is the
Arnold relation~\eqref{eq:preface-arnold} evaluated on the
Casimir: three boundary divisors of $\overline{\mathrm{FM}}_3(\bC)$,
three terms.

The $r$-matrix is an $E_1$ invariant.  The symmetric bar sees only
its $\Sigma_2$-coinvariant trace $\kappa(\cA)=\mathrm{av}(r(z))$,
a scalar.  This lossy projection is the subject of~\S1.7.

\subsection*{1.5.\enspace Chain-level: the Heisenberg ordered bar}

The Heisenberg algebra $\cH_k$: one field~$\alpha$ of conformal
weight~$1$, $\alpha(z)\alpha(w)\sim k/(z-w)^2$.  Only
$\alpha_{(1)}\alpha=k\cdot\mathbf 1$ is nonzero.

\medskip
\noindent\textbf{Degree~$2$.}\enspace
The ordered bar at tensor degree~$2$ is
$\barB^{\mathrm{ord}}_2=\bC\cdot[s^{-1}\alpha|s^{-1}\alpha]
\otimes\Gamma(\overline{\mathrm{FM}}_2,\Omega^1_{\log})$.
The residue differential acts by
\[
d_{\mathrm{res}}
\bigl(s^{-1}\alpha\otimes s^{-1}\alpha\otimes\eta_{12}\bigr)
\;=\;k\cdot\mathbf{1}.
\]
The OPE has pole order~$2$; the $d\log$ kernel absorbs one order,
producing a simple-pole residue.  This is the complete arity-$2$
ordered bar data.

\medskip
\noindent\textbf{Degree~$3$.}\enspace
At tensor degree~$3$,
$\barB^{\mathrm{ord}}_3=(s^{-1}\alpha)^{\otimes 3}
\otimes\Omega^2(\overline{\mathrm{FM}}_3)$.
Every collision stratum extracts
$\alpha_{(1)}\alpha=k\cdot\mathbf{1}$, leaving $n-2$ copies
of~$\alpha$.  Since $\mathbf{1}_{(p)}\alpha=0$ for all $p\ge 0$,
no further contraction occurs: the degree-$3$ differential
vanishes identically.  The Arnold relation kills the form-level
contribution independently.  All higher degrees vanish by the same
mechanism.

\medskip
\noindent\textbf{Bar cohomology.}\enspace
The bar complex is quadratic: all structure is captured at tensor
degree~$2$.  Bar cohomology concentrates in degree~$1$:
$H^1(\barB(\cH_k))=\bC\cdot\alpha^*$, $H^n=0$ for $n\ge 2$.
The conformal-weight-graded bar Hilbert series is rational.
The partition-graded generating function is transcendental:
$\prod_{m\ge 1}(1-x^m)^{-1}$, reflecting the Fock-space grading
of the single generator.  $\cH_k$ is chirally Koszul.

\medskip
\noindent\textbf{$R$-matrix and averaging.}\enspace
The ordered $r$-matrix: $r^{\cH}(z)=k/z$, a scalar (the OPE double
pole $k/(z-w)^2$ drops to a simple pole under $d\log$ absorption).
The spectral $R$-matrix $R(z)=\exp(k/z)$ is already
$\Sigma_2$-invariant: averaging loses nothing.
\begin{equation}\label{eq:av-heisenberg}
\mathrm{av}\bigl(r^{\cH}(z)\bigr)=k=\kappa(\cH_k).
\end{equation}
Koszul dual:
$\cH_k^!=(\operatorname{Sym}^{\mathrm{ch}}(V^*),m_0=-k\,\omega)$,
a curved symmetric chiral algebra with $m_1^2(a)=[m_0,a]$.
Complementarity (free-field branch, sum $=0$; contrast Virasoro where sum $=13$): $\kappa(\cH_k)+\kappa(\cH_k^!)=k+(-k)=0$.

\subsection*{1.6.\enspace Chain-level: the Kac--Moody ordered bar}

The affine Kac--Moody algebra $\widehat\fg_k$ for simple~$\fg$
of rank~$\ell$ and dual Coxeter number~$h^\vee$:
\[
J^a(z)J^b(w)\sim\frac{f^{ab}{}_{c}J^c(w)}{z-w}
+\frac{k\kappa^{ab}}{(z-w)^2}.
\]
Two pole orders.  The first carries the Lie bracket
$J^a_{(0)}J^b=f^{ab}{}_cJ^c$; the second carries the invariant
form $J^a_{(1)}J^b=k\kappa^{ab}\cdot\mathbf 1$.

\medskip
\noindent\textbf{Degree~$2$: two components.}\enspace
The ordered bar differential at tensor degree~$2$ decomposes as
\[
d_{\barB}\bigl([s^{-1}J^a|s^{-1}J^b]\otimes\eta_{12}\bigr)
\;=\;
\underbrace{f^{ab}{}_{c}\,[s^{-1}J^c]}_{\text{Lie bracket (simple pole)}}
\;+\;
\underbrace{k\kappa^{ab}\cdot\mathbf{1}}_{\text{curvature (double pole)}}.
\]
The Lie bracket contributes a Chevalley--Eilenberg-type
differential; the inner product contributes the curvature.
The coexistence of both poles, absent in the pure double-pole
Heisenberg, makes $\widehat\fg_k$ the first nontrivial example.

\medskip
\noindent\textbf{Degree~$3$: Serre relations.}\enspace
For $\fg=\mathfrak{sl}_3$ with simple roots $\alpha_1,\alpha_2$ and
generators $E_1,E_2$, the degree-$3$ ordered bar complex
$\barB^{\mathrm{ord}}_3=\fg^{\otimes 3}\otimes\Omega^2(\overline{\mathrm{FM}}_3)$
has dimension $8^3\times 2=1024$.
The Serre element
\[
S=[E_1|E_1|E_2]-2[E_1|E_2|E_1]+[E_2|E_1|E_1]
\]
is a bar cycle: $d_{\barB}(S\otimes\eta_{12}\wedge\eta_{13})=0$.
This is the Serre relation $(\mathrm{ad}\,E_1)^2(E_2)=0$ manifesting
as a degree-$3$ bar cocycle.  The two independent Serre relations
(for $\{i,j\}=\{1,2\}$) generate
$H^3(\barB(\widehat{\mathfrak{sl}}_{3,k}))$ at generic level.

\medskip
\noindent\textbf{Bar cohomology across Kac--Moody.}\enspace
The bar cohomology of affine Kac--Moody algebras exhibits the full
range of algebraic complexity:
\begin{center}
\small
\renewcommand{\arraystretch}{1.15}
\begin{tabular}{llll}
\textbf{Algebra} & $\dim H^n(\barB)$ & \textbf{Growth}
  & \textbf{Generating function}\\ \hline
$\widehat{\mathfrak{sl}}_2$ & $2n+1$ & linear
  & rational: $x(3-x)/(1-x)^{2}$\\
$\widehat{\mathfrak{sl}}_3$ & $\{8,36,204,\ldots\}$ & exponential
  & rational: $4x(2{-}13x{-}2x^2)/((1{-}8x)(1{-}3x{-}x^2))$\\
$\widehat\fg_k$ (general) & Garland--Lepowsky & polynomial ($\fg$ simple)
  & D-finite
\end{tabular}
\end{center}
For $\widehat{\mathfrak{sl}}_2$, the linear growth
$\dim H^n=2n+1$ is the Garland--Lepowsky concentration: bar
cohomology is generated in degree~$1$ by the $3$-dimensional
adjoint representation.  For $\widehat{\mathfrak{sl}}_3$, the
rational generating function with denominator $(1-8x)(1-3x-x^2)$
reflects the Serre relations at degree~$3$; the PBW spectral sequence
collapses at~$E_2$.

\medskip
\noindent\textbf{$R$-matrix and averaging.}\enspace
The ordered $r$-matrix is the level-prefixed Casimir:
\[
r^{\mathrm{KM}}(z)\;=\;k\,\Omega/z,
\qquad\Omega=\sum_a J^a\otimes J_a
\]
(at $k=0$ the $r$-matrix vanishes; at $k=-h^\vee$ it is critical).
Averaging collapses the Casimir tensor to its trace, with a Sugawara
shift from simple-pole self-contraction:
\begin{equation}\label{eq:av-km}
\mathrm{av}(k\,\Omega/z)
\;=\;\frac{k\dim\fg}{2h^\vee}
\;=\;\kappa_{\mathrm{dp}},
\qquad
\kappa(\widehat\fg_k)
\;=\;\kappa_{\mathrm{dp}}+\frac{\dim\fg}{2}
\;=\;\frac{(k+h^\vee)\dim\fg}{2h^\vee}.
\end{equation}
The $\dim(\fg)/2$ term is the Sugawara shift: the
simple-pole self-contraction through the adjoint Casimir eigenvalue
$2h^\vee$.  For $\widehat{\mathfrak{sl}}_2$:
$\kappa=3(k+2)/4$; for $\widehat{\mathfrak{sl}}_3$:
$\kappa=4(k+3)/3$.

The Casimir tensor structure $\Omega=\sum_a J^a\otimes J_a$, invisible
to~$\kappa$, is the data that builds the Yangian $Y(\fg)$.
Koszul dual: $\widehat\fg_k^!=\widehat\fg_{-k-2h^\vee}$
(Feigin--Frenkel duality).  Complementarity:
$\kappa(\widehat\fg_k)+\kappa(\widehat\fg_k^!)=0$
(KM branch: sum $=0$; Virasoro branch: sum $=13$).

\subsection*{1.7.\enspace The averaging map: from ordered to symmetric}

The symmetric bar
$\barB^{\Sigma}(\cA)=\operatorname{Sym}^c(s^{-1}\bar\cA)$
is the $\Sigma_n$-coinvariant quotient of the ordered bar, with the
coshuffle coproduct: $2^n$ bipartitions of an unordered set rather
than $n+1$ consecutive splittings.  This is an $E_\infty$-coalgebra,
the Beilinson--Drinfeld factorization coalgebra on $\mathrm{Ran}(X)$.
Its Koszul dual $\cA^!_{\mathrm{ch}}$ is the Francis--Gaitsgory
chiral Koszul dual via Verdier duality on $\operatorname{Ran}(X)$.

The averaging map
\[
\mathrm{av}\colon\barB^{\mathrm{ord}}(\cA)\longrightarrow\barB^{\Sigma}(\cA)
\]
takes $\Sigma_n$-coinvariants.  For a chiral algebra with OPE poles,
this descent uses the spectral $R$-matrix as twisting datum.
Applying $\mathrm{av}$ at each arity:
\begin{equation}\label{eq:av-organizing}
\begin{aligned}
\text{arity }2\colon&\quad\mathrm{av}(r(z))+\kappa_{\mathrm{sp}}=\kappa(\cA)
  &&\text{(Theorem D; $\kappa_{\mathrm{sp}}=0$ for abelian, $\dim\fg/2$ for KM)},\\
\text{arity }3\colon&\quad\mathrm{av}(\text{KZ associator})
  =\mathfrak{C}
  &&\text{(cubic shadow)},\\
\text{arity }4\colon&\quad\mathrm{av}(\text{Yangian coproduct})
  =\mathfrak{Q}
  &&\text{(quartic resonance class)}.
\end{aligned}
\end{equation}
Each arity-$n$ component of the ordered Maurer--Cartan element
$\Theta^{E_1}_\cA$ projects under $\mathrm{av}$ to the arity-$n$
shadow.  The full obstruction tower is the coinvariant image of
the ordered tower.

\medskip
\noindent\textbf{What averaging loses.}\enspace
For Heisenberg, $r(z)=k/z$ is already symmetric: averaging loses
nothing.  For Kac--Moody, $r(z)=k\,\Omega/z$ carries the Casimir
tensor $\Omega=\sum_a J^a\otimes J_a$; averaging collapses
$\Omega$ to $\dim(\fg)/(2h^\vee)$ and discards the tensor structure
that generates the Yangian $Y(\fg)$, the braided monoidal structure
on line operators, and the RTT relations.  For Virasoro,
$r(z)=(c/2)/z^3+2T/z$ carries the stress-tensor coupling at
arity~$2$; averaging collapses it to $\kappa=c/2$ and discards the
higher-spin data.

In each case: the ordered bar carries the $R$-matrix, the Yangian
lineage, the braiding.  The symmetric bar carries~$\kappa$, the
shadow tower, the genus expansion.  The first is the full algebra;
the second is its scalar fingerprint.

\subsection*{1.8.\enspace The five theorems as averaged invariants}

The main theorems are the invariants of
$\Theta^{E_1}_{\cA}$ that survive $\Sigma_n$-coinvariant projection.

\medskip
\noindent\textbf{Theorem A} (bar-cobar adjunction and Verdier
intertwining).  The cobar functor $\Omega(\barB(\cA))$ inverts
the bar construction on the Koszul locus:
$\Omega(\barB(\cA))\xrightarrow{\sim}\cA$.  Verdier duality on
$\mathrm{Ran}(X)$ sends the factorization coalgebra
$\barB_X^\Sigma(\cA)$ to the homotopy Koszul
dual~$\cA^!_\infty$, a factorization algebra whose cohomology
recovers $\cA^!$.  This identification intertwines bar-cobar
inversion with Verdier duality.

\medskip
\noindent\textbf{Theorem B} (bar-cobar inversion).  On the Koszul
locus, the counit $\varepsilon\colon\Omega(\barB(\cA))\to\cA$ is a
quasi-isomorphism at every genus.  Koszulness is characterised by
twelve equivalent conditions (\S9).

\medskip
\noindent\textbf{Theorem C} (complementarity).  The ordered
partition functions $Q^{E_1}_{g,n}(\cA)$ and $Q^{E_1}_{g,n}(\cA^!)$
sum to a topological invariant.  The coinvariant projection
yields the scalar complementarity $\kappa(\cA)+\kappa(\cA^!)=K(\cA)$,
where $K=0$ for Kac--Moody and free fields, $K=13$ for Virasoro.

\medskip
\noindent\textbf{Theorem D} (modular characteristic).  For uniform
weight, $\mathrm{obs}_g(\cA)=\kappa(\cA)\cdot\lambda_g$
(UNIFORM-WEIGHT).  Multi-weight $g\ge 2$ requires the cross-channel
correction $\delta F_g^{\mathrm{cross}}$.

\medskip
\noindent\textbf{Theorem H} (chiral Hochschild).
$\mathrm{ChirHoch}^*(\cA)$ has polynomial Hilbert series,
concentrated in cohomological degrees $\{0,1,2\}$.

\begin{remark}[The zero-dimensional shadow]\label{rem:preface-zero-dim}
When $X=\operatorname{Spec}\Bbbk$ is a point, configuration spaces
collapse, logarithmic forms vanish, and the chiral bar complex reduces
to the classical bar construction of an associative algebra.
Classical Koszul duality (Priddy, Beilinson--Ginzburg--Soergel,
Loday--Vallette) embeds on $X=\bA^1$: the affine line
deformation-retracts to a point (but the retraction is additional
data, not a canonical identification; on~$\bP^1$ the global topology
has no counterpart over a bare point).
\end{remark}

Everything above takes place at genus~$0$: the propagator is globally
defined on $\bP^1$, and the Arnold relation holds without correction.
On a curve of genus $g\ge 1$, $z_1-z_2$ has no global meaning; the
Arnold relation acquires a defect, the bar differential no longer
squares to zero, and the defect is controlled by a single scalar.


% ====================================================================
\section*{2.\quad Curvature and the genus tower}
% ====================================================================

\subsection*{2.1.\enspace The genus-$g$ propagator}

The genus-$0$ propagator $\eta_{12}=d\log(z_1-z_2)$ relied on
$z_1-z_2$ being a globally defined function on $\bP^1$.  On a compact
curve of genus $g\ge 1$, any meromorphic function with a simple zero
must also have a pole: $z_1-z_2$ has no global meaning.  The
replacement is the \emph{prime form}
$E(z_1,z_2)$: a section of $K^{-1/2}\boxtimes K^{-1/2}$
vanishing to first order along the diagonal with no other zeros.  The
bundle is $K^{-1/2}\boxtimes K^{-1/2}$; the prime form is a
$(-\tfrac12,-\tfrac12)$-differential.  In a local coordinate with
diagonal $z_1=z_2$,
\[
E(z_1,z_2)=(z_1-z_2)\bigl(1+O(z_1-z_2)\bigr)\cdot(dz_1)^{-1/2}(dz_2)^{-1/2}.
\]
Analytic continuation around the $B$-cycle $B_j$ in~$z_1$ gives
\[
E(z_1+B_j,z_2)=-\exp\bigl(-\pi i\Omega_{jj}-2\pi i\textstyle\int_{z_2}^{z_1}\omega_j\bigr)\cdot E(z_1,z_2),
\]
where $\Omega=(\Omega_{jk})$ is the period matrix and
$\omega_1,\ldots,\omega_g$ the holomorphic differentials normalised by
$\oint_{A_j}\omega_k=\delta_{jk}$.  The $A$-cycle monodromy is trivial.

The logarithmic derivative
$\omega(z_1,z_2):=d_{z_1}\log E(z_1,z_2)$ is the fundamental
meromorphic differential of the second kind.  The \emph{genus-$g$
propagator}:
\[
\eta^{(g)}_{12}=d\log E(z_1,z_2)+\pi\sum_{j,k=1}^g(\operatorname{Im}\Omega)^{-1}_{jk}\operatorname{Im}\!\Bigl(\int_{z_2}^{z_1}\omega_j\Bigr)\omega_k(z_1).
\]
The period correction $(\operatorname{Im}\Omega)^{-1}_{jk}$ compensates
$B$-cycle monodromy, rendering $\eta^{(g)}_{12}$ single-valued; this
is a full $g\times g$ matrix at genus $g\ge 2$, never collapsed to a
scalar.  The price:
$\eta^{(g)}_{12}$ is no longer holomorphic.  The propagator couples to
the Hodge bundle $\mathbb E\to\overline{\mathcal M}_g$ through
$\Omega$.

\subsection*{2.2.\enspace Curvature}

The Arnold relation at genus~$g$ acquires a defect.  The period
corrections produce a residual $(1,1)$-form, and the fibrewise bar
differential satisfies
\begin{equation}\label{eq:pref-curvature}
d_{\mathrm{fib}}^{2}\;=\;\kappa(\cA)\cdot\omega_g,
\end{equation}
where $\omega_g$ is the Arakelov $(1,1)$-form on the fibre of the
universal curve $\pi\colon\mathcal C_g\to\overline{\mathcal M}_g$, and
$\kappa(\cA)\in\Bbbk$ is the \emph{modular characteristic}, determined
entirely by the leading OPE singularity.

For Heisenberg, every residue extracts $\alpha_{(1)}\alpha=k$ and
$\kappa(\mathcal H_k)=k$.  For Kac--Moody, the curvature receives the
trace over the adjoint representation:
\[
\kappa(\widehat{\mathfrak g}_k)=\frac{(k+h^\vee)\dim\mathfrak g}{2h^\vee}.
\]
At the critical level $k=-h^\vee$ the shifted level vanishes and the
bar complex is uncurved.  The Sugawara construction is undefined (the
central charge $c=k\dim\mathfrak g/(k+h^\vee)$ has a pole) but the
curvature parameter vanishes.

For Virasoro at central charge~$c$, the OPE is
$T(z)T(w)\sim\tfrac{c/2}{(z-w)^4}+\tfrac{2T(w)}{(z-w)^2}+\tfrac{\partial T(w)}{z-w}$.
The bar propagator absorbs one pole order, so the collision
residue is $r(z)=(c/2)/z^3+2T/z$ (cubic, not quartic).
The curvature computation gives
\[
\kappa(\mathrm{Vir}_c)=c/2.
\]
The Koszul dual is $\mathrm{Vir}_{26-c}$, with
$\kappa(\mathrm{Vir}_{26-c})=(26-c)/2$; the two sum to~$13$, not zero.
At $c=13$, Virasoro is self-dual under Koszul duality.

\subsection*{2.3.\enspace Period integrals restore nilpotence}

The curvature~\eqref{eq:pref-curvature} obstructs $d_{\mathrm{fib}}^2=0$
but not the total $D_g^2=0$.  Integration against $A$-cycle periods
defines a connection on the family of bar complexes parametrised by
$\overline{\mathcal M}_g$.  The corrected total differential
\[
D_g=d_{\mathrm{fib}}+\nabla^{\mathrm{GM}}
\]
incorporates the Gauss--Manin connection on the Hodge bundle, and
$D_g^2=0$: the $(1,1)$-curvature of the fibrewise differential is
absorbed by the $(0,1)$-part of the base connection.  The fibrewise
class $\kappa\cdot\omega_g\in H^{1,1}(\mathcal C_g)$ is identified
under the period map $\overline{\mathcal M}_g\to\mathcal A_g$ with
the $(0,1)$-curvature of $\nabla^{\mathrm{GM}}$ on $R^1\pi_*\mathbb C$,
and the two curvatures cancel.

\subsection*{2.4.\enspace The genus tower (uniform weight)}

The bar complex traces a family of curved cochain complexes over
$\overline{\mathcal M}_g$.  For \emph{uniform-weight} algebras
(single-generator, or multi-generator with all generators of the same
conformal weight), the obstruction to flatness is the characteristic
class
\begin{equation}\label{eq:obs-uniform}
\operatorname{obs}_g(\cA)=\kappa(\cA)\cdot\lambda_g\in H^*(\overline{\mathcal M}_g),\qquad\lambda_g=c_g(\mathbb E),
\qquad\text{(UNIFORM-WEIGHT)}
\end{equation}
where $\mathbb E=\pi_*\omega_\pi$ is the Hodge bundle: the vector
bundle on $\overline{\mathcal M}_g$ whose fibre over $[C]$ is the
$g$-dimensional space $H^0(C,\omega_C)$.  At genus~$1$ the
factorization holds unconditionally for all families.  For
\emph{multi-weight} families at genus~$g\ge 2$ the scalar formula
\emph{fails}: a cross-channel correction from mixed-propagator
boundary graphs appears,
\begin{equation}\label{eq:obs-multiweight}
\operatorname{obs}_g(\cA)=\kappa(\cA)\cdot\lambda_g+\delta F_g^{\mathrm{cross}},\qquad\text{(ALL-WEIGHT)}
\end{equation}
(multi-weight genus expansion theorem, full monograph).  Every
appearance of $\operatorname{obs}_g$, $F_g$, or $\lambda_g$ in this
monograph carries explicit tagging: \textsc{(uniform-weight)} or
\textsc{(all-weight, with cross-channel correction)}.

Grothendieck--Riemann--Roch applied to the universal curve
$\pi\colon\mathcal C_g\to\overline{\mathcal M}_g$ computes the total
obstruction: the Chern character of $R\pi_*\omega_\pi$ involves the
Todd class of the relative tangent, the $\hat A$-genus.  The
generating function in the uniform-weight lane:
\[
\sum_{g\ge 1}\operatorname{obs}_g(\cA)\cdot x^{2g}=\kappa(\cA)\cdot\bigl(\hat A(ix)-1\bigr).
\]
Since $\hat A(ix)=(x/2)/\sin(x/2)=1+x^2/24+7x^4/5760+\cdots$, the first
Faber--Pandharipande coefficients are
\[
F_1=\tfrac{\kappa}{24},\qquad F_2=\tfrac{7\kappa}{5760},\qquad F_3=\tfrac{31\kappa}{967680},
\qquad\text{(UNIFORM-WEIGHT)}
\]
with the general term
$\kappa\cdot\tfrac{2^{2g-1}-1}{2^{2g-1}}\cdot\tfrac{|B_{2g}|}{(2g)!}$.
Each of these is the scalar-lane projection; for multi-weight
families at $g\ge 2$, add $\delta F_g^{\mathrm{cross}}$.  Verify
$F_1=\kappa/24$ as sanity check (Cauchy normalization
$1/(2\pi i)$, not $1/(2\pi)$).

The scalar $\kappa(\cA)$ is only the leading term.  The bar complex
carries higher-arity data: a universal Maurer--Cartan element
$\Theta_\cA$ whose projection to arity~$r$ yields a shadow invariant
$S_r(\cA)$.  Write $\alpha:=S_3$ for the cubic shadow.
On each primary deformation line~$L$ (a one-parameter subfamily
at fixed spin content), the shadow
generating function $H(t)=\sum_{r\ge 2}rS_rt^r$ satisfies the
algebraic relation
\begin{equation}\label{eq:preface-riccati-hook}
H(t)^2=t^4\,Q_L(t),\qquad Q_L(t)=4\kappa^2+12\kappa S_3\,t+(9S_3^2+16\kappa S_4)\,t^2,
\end{equation}
a quadratic polynomial in three genus-$0$ invariants
(Riccati algebraicity theorem, full monograph).  The infinite tower is
algebraic of degree~$2$, governed by the discriminant
$\Delta=8\kappa S_4$.

The shadow obstruction tower is not merely formal.  Its projections
yield: the Pixton ideal of the modular curve, the integrable
hierarchies of Gelfand--Dickey type, the Hitchin spectral curve of
the shadow connection, the Baxter $Q$-operator, and the
Bekenstein--Hawking entropy formula of three-dimensional gravity
(twisted holography chapter, full monograph).  These are
projections of one object, not separate results.

\subsection*{2.5.\enspace Complementarity}

Verdier duality on the Ran space splits the center local system into
two halves, one controlled by $\cA$ and one by $\cA^!$, carrying a
shifted-symplectic structure:
\[
R\Gamma(\overline{\mathcal M}_g,\mathcal Z_\cA)\simeq\mathbf Q_g(\cA)\oplus\mathbf Q_g(\cA^!).
\]
What $\cA$ sees as deformation, $\cA^!$ sees as obstruction.  The
natural pairing
\[
\mathbf Q_g(\cA)\otimes\mathbf Q_g(\cA^!)\to R\Gamma(\overline{\mathcal M}_g,\omega_{\overline{\mathcal M}_g})\simeq\Bbbk[-(3g-3)]
\]
is a $(-(3g-3))$-shifted symplectic form, and the two summands are
Lagrangian: isotropic of half the total dimension.  The scalar sum
rule is family-dependent (never universal):
\[
\kappa(\cA)+\kappa(\cA^!)=\begin{cases}0&\text{(Kac--Moody, free field, lattice)},\\ 13&\text{(Virasoro)},\\ 250/3&(\mathcal W_3),\\ K_N\cdot(H_N-1)&(\mathcal W_N),\end{cases}
\]
where $K_N=c(\mathcal W_N)+c(\mathcal W_N^!)$ is the central-charge
conductor and $H_N=\sum_{j=1}^N 1/j$ the harmonic number.
The scalar $\kappa+\kappa'$ is the kappa-level Koszul conductor;
it is zero for KM and free fields, nonzero for $\mathcal W$-algebras.
Self-dual at $c=13$ for Virasoro, $c=50$ for $\mathcal W_3$; qualify
``self-dual at $c=13$ under Koszul
duality,'' never unqualified.

\subsection*{2.6.\enspace Feigin--Frenkel duality}

The Sugawara field
$T(z)=\tfrac{1}{2(k+h^\vee)}\sum_a{:}J^a(z)J^a(z){:}$
has central charge $c=k\dim\mathfrak g/(k+h^\vee)$.  At the critical
level $k=-h^\vee$ the denominator vanishes: $T(z)$ is undefined, not
``$c$ diverges.''  The chiral algebra acquires the Feigin--Frenkel
center $\mathfrak z(\widehat{\mathfrak g}_{-h^\vee})$, isomorphic to
functions on the space of ${}^L\mathfrak g$-opers on the disc.

The Feigin--Frenkel involution $k\leftrightarrow-k-2h^\vee$ is Koszul
duality on configuration spaces: it interchanges $\cA$ and $\cA^!$
while preserving the bar-cobar adjunction.  At the critical level, the
Koszul dual is the Langlands-dual critical-level
algebra~$\widehat{{}^L\mathfrak g}_{-{}^Lh^\vee}$, not a renormalised
copy of the original.

\subsection*{2.7.\enspace The total space}

The universal configuration fibration
$\pi_n\colon\overline{\mathcal C}_n(\mathcal C_g)\to\overline{\mathcal M}_g$
packages the bar complex into a single global object.  Its boundary
divisors carry two distinct geometric roles.  \emph{Collision
divisors} (vertical): the fibrewise boundary of
$\overline{\mathrm{FM}}_n(X_s)$, producing the bar differential
$d_{\mathrm{res}}$ (OPE residues).  \emph{Degeneration divisors}
(horizontal): the boundary of $\overline{\mathcal M}_g$ itself (nodal
curves), producing the period corrections and Gauss--Manin terms.

Algebraic boundary (Borcherds) meets geometric boundary (period map
and clutching).  Their interaction is the curvature~$\kappa$: the
fibrewise bar differential fails to be nilpotent precisely because
the vertical and horizontal boundaries do not commute.


% ====================================================================
\section*{3.\quad Modular homotopy theory}
% ====================================================================

The genus tower demands an organising framework.  The fibrewise
differential, the Gauss--Manin correction, and the clutching morphisms
at boundary strata satisfy compatibility relations dictated by the
stratification of $\overline{\cM}_{g,n}$.  The governing algebraic
structure is the \emph{modular operad}: an operad whose operations are
indexed by connected stable graphs of arbitrary genus.

The \emph{modular homotopy type} of $\cA$ is the formal moduli problem
classifying the inverse system of bar complexes,
\[
\mathcal M^{\mathrm{mod}}_\cA(R):=\operatorname{MC}(\mathfrak g^{\mathrm{mod}}_\cA\widehat\otimes_{\mathfrak m}R)/{\sim}
\]
for local Artin rings~$R$.  Here $\mathfrak g^{\mathrm{mod}}_\cA$ is
the \emph{modular convolution algebra}: a complete filtered cyclic
$L_\infty$-algebra.  Three inputs enter: a homotopy chiral algebra on
the curve, a cooperadic chain complex on logarithmic configuration
spaces, and a convolution $L_\infty$-structure mediating between them.
The bar complex, the Koszul dual, the genus tower, the
complementarity splitting, and the shadow obstruction tower are all
representations of $\mathcal M^{\mathrm{mod}}_\cA$.

The primitive story lives on ordered configurations.  The $E_1$
convolution algebra $\mathfrak g^{\mathrm{mod},E_1}_\cA$, built on
ribbon graphs from $F\!\Ass$, is the primary object; the modular
convolution algebra $\gAmod$ is its image under the av map from
ordered to symmetric.  Every construction in this section begins with
the $E_1$ structure and descends.

\subsection*{3.1.\enspace The modular operad}

At genus~$0$, the combinatorics is governed by trees; at higher
genus, by stable graphs with genus labels at their vertices.  A
\emph{modular operad} (Getzler--Kapranov~\cite{GeK98}) has operations
indexed by connected stable graphs.  A \emph{connected stable graph}
$\Gamma$ has a label $g(v)\in\mathbb Z_{\ge 0}$ at each vertex
satisfying $2g(v)-2+\operatorname{val}(v)>0$ and total genus
$g(\Gamma)=b_1(\Gamma)+\sum_v g(v)$; trees are the genus-$0$ case.
The structure maps are \emph{separating gluing}
$\xi_{\mathrm{sep}}\colon\cO(g_1,n_1{+}1)\otimes\cO(g_2,n_2{+}1)\to\cO(g_1{+}g_2,n_1{+}n_2)$
and \emph{non-separating gluing}
$\xi_{\mathrm{ns}}\colon\cO(g,n{+}2)\to\cO(g{+}1,n)$.

A stable graph records a degeneration of a smooth curve into a nodal
curve: vertices are irreducible components, edges are nodes, genus
labels are component genera, half-edges are marked points.  The set
$\mathsf{Gr}^{\mathrm{st}}_{g,n}$ parametrises the boundary
stratification of $\overline{\cM}_{g,n}$.  Every codimension-one
boundary divisor is a one-edge degeneration; every codimension-two
stratum is the intersection of exactly two codimension-one divisors,
appearing with opposite orientations.

The \emph{Feynman transform} replaces trees by stable graphs,
\begin{equation}\label{eq:preface-feynman-sum}
F\cO(g,n)=\bigoplus_{\Gamma\in\mathsf{Gr}^{\mathrm{st}}_{g,n}}\frac{\hbar^{g(\Gamma)}}{|\operatorname{Aut}(\Gamma)|}\bigotimes_{v\in V(\Gamma)}\cO\bigl(g(v),\operatorname{In}(v)\bigr).
\end{equation}

\subsection*{3.2.\enspace The Feynman transform differential}

The differential contracts one edge at a time,
\[
d_{F\cO}(\gamma_\Gamma)=\sum_{e\in E(\Gamma)}(-1)^{\epsilon(e)}\xi_e(\gamma_{\Gamma/e}),
\]
with the sign from the position of $e$ in the ordering on $E(\Gamma)$
recorded by $\det(E(\Gamma))$.  The identity $d^2=0$ is the
codimension-$2$ cancellation: each pair of edges contracts in either
order, and the two terms cancel by associativity of~$\cO$.  At
genus~$0$, $F\cO$ reduces to the operadic bar construction.

When $\cO=\Com$, each vertex factor in~\eqref{eq:preface-feynman-sum}
is one-dimensional, and the Feynman transform becomes the stable
graph complex
\[
F\Com(g,n)=\bigoplus_{\Gamma\in\mathsf{Gr}^{\mathrm{st}}_{g,n}}\frac{\hbar^{g(\Gamma)}}{|\operatorname{Aut}(\Gamma)|}\det(E(\Gamma)).
\]
When $\cO=\Ass$, the Feynman transform $F\!\Ass$ sums over ribbon
graphs (fatgraphs); this is the $E_1$-primitive input for the ordered
convolution algebra.

\subsection*{3.3.\enspace The convolution ladder}

At each level of generality, the bar-type construction is a
convolution dg Lie algebra whose Maurer--Cartan elements classify
algebraic structures:
\begin{equation}\label{eq:preface-convolution-ladder}
\begin{array}{r@{\;:\quad}c@{\qquad}l}
\text{algebra} & \Hom(B(\Ass),\End_A) & A_\infty\text{-structures on }A\\[3pt]
\text{operad}  & \Hom_\Sigma(B(\cP),\End_V) & \cP_\infty\text{-structures on }V\\[3pt]
\text{modular} & \Hom_\Sigma(F\cO,\End_{\barB(\cA)}) & \text{consistent genus towers}
\end{array}
\end{equation}
The first line is Hochschild; the second is the Ginzburg--Kapranov
deformation complex; the third governs this book.  The Lie bracket is
the convolution bracket
\[
[f,g]=\mu\circ(f\otimes g)\circ\Delta-(-1)^{|f||g|}\mu\circ(g\otimes f)\circ\Delta,
\]
with $\Delta$ the cooperadic decomposition and $\mu$ the operadic
composition.

The \emph{primitive $E_1$ convolution}
\[
\mathfrak g^{\mathrm{mod},E_1}_\cA:=\Hom_\Sigma\bigl(F\!\Ass,\End_{\barB^{E_1}(\cA)}\bigr)
\]
is the ordered, ribbon-graph input: the algebraic form of the 't~Hooft
$1/N$ expansion.  Setting $\hbar=1/N^2$, the MC element
$\Theta^{E_1}_\cA\in\operatorname{MC}(\mathfrak g^{\mathrm{mod},E_1}_\cA)$
is the genus expansion $\sum_g N^{2-2g}F_g(\lambda)$ summed over
ribbon graphs weighted by the genus of their embedding surface.
Planar ($g=0$) diagrams dominate at large~$N$; the shadow obstruction
tower $\Theta^{E_1,\le r}$ is the truncation to $r$-point 't~Hooft
interactions.

The \emph{modular convolution algebra} is the av-image,
\begin{equation}\label{eq:preface-gmod}
\gAmod:=\Hom_\Sigma\bigl(F\Com,\End_{\barB(\cA)}\bigr).
\end{equation}
The coinvariant map $F\!\Ass\to F\Com$ projects
$\Theta^{E_1}\mapsto\Theta$: colour-ordered amplitudes to full
amplitudes.  On each primary line, the shadow tower is algebraic of
degree~$2$: the entire infinite sequence is the Taylor expansion of
$\sqrt{Q_L}$ for a single quadratic $Q_L$, the shadow metric
(\S6 below).

\medskip
\noindent\textbf{Remark} (Homotopy types in three stages).  The three
levels of the convolution ladder mirror three levels of homotopy
theory.  \emph{Over a field} (the bar complex of a graded algebra):
$B(A)=(T^c(s^{-1}\bar A),d_B)$ determines
the homotopy type of an associative algebra; formality corresponds to
intrinsic formality; Massey products measure the obstruction.
\emph{Over a manifold}: the Sullivan minimal model $(\Lambda V,d)$
determines the rational homotopy type; K\"ahler manifolds are formal.
The Heisenberg algebra is the vertex-algebraic analogue: its shadow
obstruction tower terminates at arity~$2$, the Gaussian archetype.
\emph{Over a curve at all genera}: the $A_\infty$-products $m_k$ of
the first stage become the modular $L_\infty$-brackets $\ell_k^{(g)}$;
formality becomes the Gaussian archetype; the shadow obstruction tower
is the modular Massey product tower.  The shadow depth classification
(\S6) makes this analogy precise: the four classes G, L, C, M
correspond to four levels of modular formality.

\medskip
\noindent\textbf{Koszulness as central question.}\enspace
Within this formality classification, the Koszulness property is the
central organizing question of the programme.  A chiral algebra is
\emph{chirally Koszul} iff its bar cohomology concentrates in
diagonal degrees and the cobar-bar counit
$\Omega(\barB(\cA))\to\cA$ is a quasi-isomorphism.  Twelve equivalent
conditions (ten unconditional, one conditional on perfectness, one
one-directional) characterise this locus.  The Koszulness programme
(\S9) proves the core equivalences from PBW concentration, bar-cobar
inversion, Hochschild polynomial growth, and Koszul-complex
acyclicity.  All standard families are Koszul; shadow depth classifies
the complexity of Koszul algebras, not Koszulness itself.

\subsection*{3.4.\enspace Homotopy chiral algebras}

The convolution algebra requires a cochain-level algebraic input, not
merely the cohomology-level chiral bracket.  The \v Cech totalisation
provides this: it replaces the strict chiral algebra with a
homotopy-coherent model carrying all derived OPE data.

Fix an affine cover $\mathfrak U=\{U_\alpha\}$ of~$X$.  The \v Cech
totalisation
\[
A^{\mathrm{ch}}_\infty:=\check C^\bullet(\mathfrak U;\cA)=\bigoplus_{p\ge 0}\prod_{\alpha_0<\cdots<\alpha_p}\cA(U_{\alpha_0}\cap\cdots\cap U_{\alpha_p})
\]
carries a $\mathrm{Ch}_\infty$-algebra structure (Malikov--Schechtman):
the chiral bracket restricted to overlaps satisfies Jacobi up to a
coherent system of higher homotopies, the \emph{secondary Borcherds
operations}
\[
F_n\colon\operatorname{Lie}(n)\otimes\cY(n)\to P_n(\{A^{\mathrm{ch}}_\infty\},A^{\mathrm{ch}}_\infty),\qquad n\ge 2,
\]
encoding derived OPE data at all pole orders.  $F_2$ is the chiral
bracket; $F_3$ is the Jacobiator homotopy arising from
$\overline{\mathcal M}_{0,4}\cong\mathbb P^1$, whose boundary has
three points corresponding to OPE compositions
$(a_{(m)}b)_{(n)}c$ and cyclic permutations.  Higher operations $F_n$
are determined inductively by $C_*(\overline{\mathcal M}_{0,n+1})$;
the transferred operations are the associahedral $A_\infty$-corrections.

Different affine covers yield quasi-isomorphic $\mathrm{Ch}_\infty$-algebras.
The strict chiral algebra $\cA$ is the $H^0$ of this structure: the
primitive local object is the homotopy chiral algebra, and the strict
algebra appears after rectification.  Vallette rectification: the
homotopy category of $\mathrm{Ch}_\infty$-algebras is equivalent to
the homotopy category of strict chiral algebras, via the bar-cobar
adjunction.

\subsection*{3.5.\enspace Logarithmic Fulton--MacPherson spaces}

The classical FM compactification records collisions on a smooth
variety.  Its logarithmic extension (Mok) records collisions along a
divisor and provides the normal-crossings boundary structure on which
$D^2=0$ depends.

For a simple-normal-crossing pair $(X,D)$, Mok constructs
$\mathrm{FM}_n(X|D)$: a smooth variety with corners compactifying
$\mathrm{Conf}_n(X\setminus D)$, with boundary stratification indexed
by \emph{rigid planted-forest types}.  A planted-forest type $\rho$
consists of a partition of $\{1,\ldots,n\}$ into blocks (the roots),
a rooted tree per block recording the hierarchy of collisions, and a
grid-depth assignment $w_v\in\mathbb Z_{\ge 0}$ at each vertex
recording the order of tangency with~$D$.  The ordinary FM is the
case $D=\emptyset$.

The \emph{degeneration formula}: for a semistable $W\to B$ with
singular fibre $W_0=\bigcup Y_v$ and a rigid type $\rho$,
\[
\mathrm{FM}_n(W/B)(\rho)\longrightarrow\prod_{v\in V(S_\rho)}\mathrm{FM}_{I_v}(Y_v|D_v)
\]
is a proper birational correspondence factoring the global FM
geometry into a fibre product of local pieces.  Tropicalisation:
$\operatorname{Trop}(\mathrm{FM}_n(\mathbb C|D))=G_{\mathrm{pf}}$, the
planted-forest poset with grafting as partial order.  The planted
forest poset is the combinatorial skeleton of the geometric input,
just as the stable graph complex is the skeleton of the Feynman
transform.

The \emph{Mok codimension formula}: for a type with grid depths
$(w_1,\ldots,w_r)$ and tree vertex counts $(|V(T_j)|)_{j=1,\ldots,s}$,
\[
\operatorname{codim}(W_\rho)=\sum_{i=1}^r w_i+\sum_{j=1}^s(|V(T_j)|-1).
\]
For $\mathrm{FM}_3$: four codimension-one strata ($\{12\}$, $\{23\}$,
$\{13\}$, $\{123\}$), three codimension-two strata
($\{12{<}123\}$, $\{23{<}123\}$, $\{13{<}123\}$).

The chain complexes $\cC^{\log\mathrm{FM}}_{X|D}(n):=C_\bullet(\mathrm{FM}_n(X|D))$
carry cooperadic structure: the inclusion of boundary strata defines
cocomposition maps.  For a rigid stable graph~$\Gamma$,
\[
\Delta^{\log}_\Gamma:=(\nu_\Gamma)_*\circ\operatorname{Res}_{D^{\log}_\Gamma}\colon\cC^{\log\mathrm{FM}}_{\mathrm{mod}}(g,n)\to\bigotimes_{v\in V(\Gamma)}\cC^{\log\mathrm{FM}}_{\mathrm{mod}}\bigl(g(v),I_v\sqcup H(v)\bigr)[-|E_{\mathrm{int}}(\Gamma)|],
\]
where $\nu_\Gamma$ is the Mok birational correspondence and
$\operatorname{Res}_{D^{\log}_\Gamma}$ is the residue along the
logarithmic boundary divisor.

\subsection*{3.6.\enspace The convolution $L_\infty$-algebra}

The cooperadic and operadic inputs of \S\S3.4--3.5 must be coupled.
The coupling datum is a twisting morphism: a linear map satisfying a
Maurer--Cartan equation in the convolution complex.

A twisting morphism $\alpha\colon\cC\to\cP$ (satisfying
$\partial\alpha+\alpha\star\alpha=0$) determines a convolution
$sL_\infty$-algebra $\operatorname{hom}_\alpha(\cC,\cP)$
(Robert-Nicoud--Wierstra, after Vallette) with underlying space
$\prod_{n\ge 1}\Hom_{\Sigma_n}(\cC(n),\cP(n))$.  The brackets
\[
\ell_k(f_1,\ldots,f_k)(x)=\sum\gamma_\cP\bigl((\alpha\circ 1)(f_1\otimes\cdots\otimes f_k)\Delta^{(k)}_\cC(x)\bigr)
\]
come from iterated cooperadic decomposition.  $\ell_1$ is the
convolution differential; $\ell_2$ is the Lie bracket of
coderivations; higher $\ell_k$, $k\ge 3$, arise whenever $\cC$ is not
concentrated in a single arity.

The $sL_\infty$-algebra extends to $\infty$-morphisms in either slot
separately but not both simultaneously: the convolution is
functorial in each slot, but does not assemble into a bifunctor.  The
MC3 categorical lift proceeds one slot at a time.  The bar-cobar
adjunction is a Quillen equivalence (Vallette); different
$\mathrm{Ch}_\infty$ presentations yield quasi-isomorphic
deformation theories.

\subsection*{3.7.\enspace Assembly}

The modular convolution $L_\infty$-algebra:
\begin{equation}\label{eq:pref-convolution}
\mathfrak g^{\mathrm{mod}}_\cA:=\Convinf\!\Bigl(\cC^{\log\mathrm{FM}}_{\mathrm{mod}},\operatorname{End}_{\mathrm{Ch}_\infty}(A^{\mathrm{ch}}_\infty)\Bigr).
\end{equation}
The cooperadic input is the log-FM chain cooperad
$\cC^{\log\mathrm{FM}}_{\mathrm{mod}}$; the operadic input is the
$\mathrm{Ch}_\infty$-endomorphism operad of
$A^{\mathrm{ch}}_\infty$.  The $E_1$ cousin
$\mathfrak g^{\mathrm{mod},E_1}_\cA$ uses the ribbon-graph-indexed
cooperad and the ordered bar endomorphism operad.

One-slot convention: fix the log-FM (resp.\ ribbon) cooperad as the
``geometric'' slot; vary the chiral algebra input as the ``algebraic''
slot by $\infty$-morphisms.

The \emph{strict chart}: the coderivation dg Lie algebra
\[
\mathfrak g^{\mathrm{mod},\log}_\cA(\mathfrak U):=\Convstr\!\bigl(\cC^{\log\mathrm{FM}}_{\mathrm{mod}},\operatorname{End}_{\mathrm{Ch}_\infty}(A^{\mathrm{ch}}_\infty)\bigr)
\]
has classical commutator bracket, no higher $L_\infty$ brackets, but
the same Maurer--Cartan moduli.  The strict chart suffices for
constructing $\Theta_\cA$ and computing shadows; the full $L_\infty$
structure is needed for transfer, formality, and gauge equivalence.

\subsection*{3.8.\enspace The modular bar functor}

In log-FM coordinates,
\[
\mathrm B^{\log}_{\mathrm{mod}}(\cA)(g,n)=\bigoplus_{\Gamma\in\mathsf{Gr}^{\mathrm{st}}_{g,n}}\left(\bigotimes_{v\in V(\Gamma)}\bar B^{\mathrm{ch}}(A^{\mathrm{ch}}_\infty;\operatorname{In}(v))\otimes C_\bullet(\mathrm{FM}^{\log}_\Gamma)\otimes\orline{E_{\mathrm{int}}(\Gamma)}\right)_{\!\operatorname{Aut}(\Gamma)}.
\]
The ingredients: at each vertex the transferred
$\mathrm{Ch}_\infty$-bar contribution on its incoming half-edges;
at each internal edge the propagator $P_\cA=H_\cA^{-1}$; the geometric
factor $C_\bullet(\mathrm{FM}^{\log}_\Gamma)$; the orientation line
$\orline{E_{\mathrm{int}}(\Gamma)}=\det(\Bbbk E(\Gamma))\otimes\det(H_1(\Gamma;\Bbbk))^{-1}$;
the automorphism quotient.

\subsection*{3.9.\enspace The six-component differential}

The total differential has six components, whose mutual cancellations
constitute $D^2=0$:
\[
D^{\log}_{\mathrm{mod}}=\underbrace{d_{\check C}}_{\text{\v Cech}}+\underbrace{d_{\mathrm{Ch}_\infty}}_{\text{transferred chiral}}+\underbrace{d_{\mathrm{coll}}}_{\text{collision}}+\underbrace{d_{\mathrm{sew}}}_{\text{separating}}+\underbrace{d_{\mathrm{pf}}}_{\text{planted-forest}}+\underbrace{\hbar\Delta}_{\text{genus-raising}}.
\]
$d_{\check C}$ is the \v Cech differential of the cover, acting on
\v Cech indices.  $d_{\mathrm{Ch}_\infty}$ is the transferred chiral
differential on each vertex, from secondary Borcherds operations.
$d_{\mathrm{coll}}$ contracts an internal edge by extracting the OPE
residue at the corresponding codimension-one collision stratum of
$\overline{\mathrm{FM}}_n(X)$.  The three positive-genus boundary
operators act on a pure tensor $a\otimes\eta$ indexed by $\Gamma$:
\begin{align*}
d_{\mathrm{sew}}(a\otimes\eta)&=\sum_{e\in E_{\mathrm{sep}}(\Gamma)}\pm\mu_e(a)\otimes(\xi_e)_*\operatorname{Res}_{D_e^{\log}}(\eta),\\
d_{\mathrm{pf}}(a\otimes\eta)&=\sum_{F\in\mathsf{PF}^{\mathrm{rig}}(\Gamma)}\pm\mu_F(a)\otimes(\xi_F)_*\operatorname{Res}_{D_F^{\log}}(\eta),\\
\hbar\Delta(a\otimes\eta)&=\hbar\sum_{e\in E_{\mathrm{ns}}(\Gamma)}\pm\operatorname{Tr}_e(a)\otimes(\xi_e)_*\operatorname{Res}_{D_e^{\log}}(\eta).
\end{align*}
Separating edges give the clutching morphism (gluing two lower-genus
curves); non-separating edges give the genus-raising self-gluing
(BV operator); planted-forest strata give the genuinely logarithmic
corrections from Mok.  Explicitly,
\[
d_{\mathrm{pf}}=\sum_{\rho\in\mathsf{PF}^{\mathrm{rig}}}\epsilon_\rho(\kappa_\rho)_*\mathrm{pr}_\rho^*\otimes\mu_\rho:
\]
Mok's degeneration formula internalised as a summand of the total
differential.

\subsection*{3.10.\enspace $D^2=0$}

The nilpotence $(D^{\log}_{\mathrm{mod}})^2=0$ holds by three
cancellation mechanisms.  \emph{Vertex labels}: $d_{\check C}^{\,2}=0$,
$d_{\mathrm{Ch}_\infty}^{\,2}=0$, and
$\{d_{\check C},d_{\mathrm{Ch}_\infty}\}=0$ by the
$\mathrm{Ch}_\infty$-algebra axioms on $A^{\mathrm{ch}}_\infty$.
\emph{One-edge expansions}: anticommutators of a vertex differential
with an edge-contraction differential vanish by the Leibniz rule for
coderivations.  \emph{Codimension-two cancellation}: the sum of all
products of two distinct edge-contraction operators vanishes, because
the codimension-two boundary of Mok's log-FM compactification is a
normal-crossings divisor (Mok, Theorem~3.3.1) and each codimension-two
face appears as the intersection of exactly two codimension-one faces
with opposite orientations.

The third mechanism is the nontrivial one.  The normal-crossings
property is essential: without it, codimension-two faces could have
excess intersection, and the cancellation would fail.


% ====================================================================
\section*{4.\quad The universal Maurer--Cartan element}
% ====================================================================

The modular convolution algebra $\gAmod$ is a complete filtered
$L_\infty$-algebra (\S3.7).  Its total differential
$D_\cA$ already squares to zero.  Decomposing $D_\cA=d_0+\Theta_\cA$
into a tree-level local part and its complement produces a universal
element $\Theta_\cA\in\operatorname{MC}(\gAmod)$ satisfying the
Maurer--Cartan equation as a formal consequence of $D_\cA^2=0$.  This
element requires no equation to solve and no convergence hypothesis.
Every invariant in this survey is a projection of~$\Theta_\cA$: the
five theorems A--D and H are its five $\Sigma_n$-coinvariant shadows.

\subsection*{4.1.\enspace $\Gamma$-amplitudes and Taylor coefficients}

Each stable graph produces a multilinear operation on the convolution
algebra, built from chiral operations at vertices, propagators along
edges, and geometric factors from log-FM strata.  For
$\Gamma\in\mathsf{Gr}^{\mathrm{st}}$ with $|V(\Gamma)|=k$ vertices and
elements $f_1,\ldots,f_k\in\gAmod$, the \emph{$\Gamma$-amplitude}
assembles
\[
\ell_\Gamma(f_1,\ldots,f_k)=\mu_\Gamma\circ(\alpha^{\mathrm{mod}}_\Gamma\otimes f_1\otimes\cdots\otimes f_k)\circ\Delta^{\log}_\Gamma,
\]
a $k$-linear map of cohomological degree $2-2g(\Gamma)-k$.  Here
$\mu_\Gamma$ is the graph composition (transferred chiral operations
at vertices, propagators along edges), $\alpha^{\mathrm{mod}}_\Gamma$
is the modular twisting morphism component, and $\Delta^{\log}_\Gamma$
is the cooperadic cocomposition.

The Taylor coefficients of the modular $L_\infty$ structure sum over
stable graphs of fixed genus and vertex count:
\[
\ell_k^{(g)}=\sum_{\substack{\Gamma\in\mathsf{Gr}^{\mathrm{st}},\;|V(\Gamma)|=k\\b_1(\Gamma)+\sum_v g(v)=g}}\frac{1}{|\operatorname{Aut}(\Gamma)|}\ell_\Gamma.
\]
The homotopy Jacobi identity is encoded by $D^2=0$ on the total
convolution algebra.

\medskip
\noindent\textbf{Low-order Taylor coefficients.}\enspace
At tree level, the binary bracket is the transferred Lie bracket:
$\ell_2^{(0)}(\alpha,\beta)=\pi[m_2(\iota(\alpha),\iota(\beta))]$.
The ternary bracket sums over the three binary trees in
$\mathsf{Trees}_3$ (channels $s$, $t$, $u$), each inserting a homotopy
$h$ at the internal edge.  The quaternary bracket sums over the
fifteen labeled trees in $\mathsf{Trees}_4$ (two unlabeled topologies:
three balanced $((ab)(cd))$ and twelve caterpillar $(((ab)c)d)$),
each with two propagator insertions; this is the classical homological
perturbation lemma in its $L_\infty$ incarnation.

The first genus-raised bracket is unary, the BV operator on the single
graph with one vertex and one self-loop:
\[
\ell_1^{(1)}(\alpha)=\operatorname{tr}(P_\cA\circ\alpha\circ\delta^{\mathrm{ns}})=\sum_i\langle e^i,\alpha(e_i)\rangle_{\mathrm{Sh}},
\]
trace over the propagator contracting two half-edges at the
non-separating node of $\overline{\cM}_{1,1}$.  The genus-one binary
bracket
\[
\ell_2^{(1)}(\alpha,\beta)=\operatorname{tr}\bigl(P_\cA\circ m_2(\iota\alpha,-)\circ\delta^{\mathrm{ns}}\circ\iota\beta\bigr)+(\alpha\leftrightarrow\beta)
\]
accounts for the separating node of $\overline{\cM}_{1,2}$.  These
low-order brackets are the only data needed for shadow computations
through arity~$4$.

\subsection*{4.2.\enspace The bar-intrinsic construction}

$\Theta_\cA$ requires no equation to solve: it is extracted from a
differential that already squares to zero.

Decompose the total bar differential
$D_\cA\colon\mathrm B^{\mathrm{mod}}(\cA)\to\mathrm B^{\mathrm{mod}}(\cA)$
as $D_\cA=d_0+\Theta_\cA$, where
$d_0=d_{\mathrm{int}}+[\tau,-]$ is the tree-level local part.  The
\emph{universal Maurer--Cartan element} is the remainder,
\begin{equation}\label{eq:pref-theta}
\Theta_\cA:=D_\cA-d_0=\sum_{\Gamma\in\mathsf{Gr}^{\mathrm{st}}}\frac{W^{\log\mathrm{FM}}_\Gamma}{|\operatorname{Aut}(\Gamma)|}\otimes\Phi^\cA_\Gamma\in\operatorname{MC}(\gAmod),
\end{equation}
with $W^{\log\mathrm{FM}}_\Gamma\in H^*(\overline{\cM}^{\log}_{g(\Gamma),n(\Gamma)})$
the log-FM fundamental class of the stratum and
$\Phi^\cA_\Gamma\in\Hom(\bar\cA^{\otimes n(\Gamma)},\Bbbk)$ the chiral
amplitude (propagators on internal edges, operations at vertices,
cyclic trace at the output).

The MC property is immediate.  $D_\cA^2=0$ by the normal-crossings
property of $\overline{\cM}_{g,n}^{\log}$; $d_0^2=0$ by construction;
expanding $(d_0+\Theta_\cA)^2=0$ gives
\[
[d_0,\Theta_\cA]+\tfrac12[\Theta_\cA,\Theta_\cA]=0.
\]
The MC property is a formal consequence of $\partial^2=0$ on the
moduli space of curves, requiring no simple-Lie restriction, no
convergence hypothesis, and no choice of basis.  This is MC2
(conjecture-turned-theorem) in its cleanest form.

The $E_1$ analogue $\Theta^{E_1}_\cA$ is constructed identically on
the ordered, ribbon-graph-indexed convolution algebra, and
$\Theta_\cA=\mathrm{av}(\Theta^{E_1}_\cA)$ under the coinvariant map
$F\!\Ass\to F\Com$.

\subsection*{4.3.\enspace $\Theta_\cA$ as universal modular twisting morphism}

The Maurer--Cartan space of the modular convolution $L_\infty$
algebra is identified with the space of modular twisting morphisms,
\[
\operatorname{MC}_\bullet(\gAmod)\simeq\operatorname{Tw}_{\alpha^{\mathrm{mod}}}\Bigl(\cC^{\log\mathrm{FM}}_{\mathrm{mod}},\operatorname{End}_{\mathrm{Ch}_\infty}(A^{\mathrm{ch}}_\infty)\Bigr).
\]
Under this identification $\Theta_\cA$ is the \emph{universal modular
twisting morphism}, the operadic analogue of the canonical flat
connection on a principal bundle.  A twisting morphism
$\alpha\colon\cC\to\cP$ determines a twisted composition product
$\cC\circ_\alpha\cP$ whose differential squares to zero iff $\alpha$
satisfies the MC equation.  $\Theta_\cA$ is the unique twisting
morphism that simultaneously twists the modular bar complex and
intertwines with the log-FM boundary structure.

The principal-bundle analogy is exact.  The modular bar bundle
$\mathrm B^{\log}_{\mathrm{mod}}(\cA)\to\overline{\cM}_{g,n}$ is the
principal bundle; $\Theta_\cA$ is the connection; the MC equation is
flatness; the shadow tower (\S6) is the Chern--Weil theory.  The
structure group is the pro-unipotent MC gauge group
$\exp(\mathfrak g^{\mathrm{mod},0}_\cA)$.

\subsection*{4.4.\enspace Tridegree, depth, and convergence}

Each coefficient of $\Theta_\cA$ carries a natural tridegree
\[
(g,n,d)=(\text{loop genus},\text{external arity},\text{log boundary depth}),
\]
with $d(\rho)=\#E_{\mathrm{sep}}(\Gamma_\rho)+\#E_{\mathrm{ns}}(\Gamma_\rho)+\operatorname{depth}_{\mathrm{pf}}(\Gamma_\rho)$
the codimension of the log-FM stratum.  The depth filtration
$F^d\gAmod=\bigoplus_{d'\ge d}(\gAmod)_{*,*,d'}$ makes the convolution
algebra complete:
\[
\gAmod=\varprojlim_d\gAmod/F^d\gAmod.
\]
The inverse limit $\Theta_\cA=\varprojlim_r\Theta_\cA^{\le r}$
converges unconditionally: no analytic input, no growth estimate, no
restriction on the algebra.  Convergence is purely algebraic,
guaranteed by completeness of the depth filtration.  This is MC4
(completion tower): proved.  Generating depth (the arity at
which higher operations are determined by lower ones) is distinct
from algebraic depth (whether the tower terminates).

At fixed $(g,n)$, only finitely many stable graph types contribute.
The depth filtration refines this further: at depth $d$, the stratum
has codimension $d$ in $\overline{\cM}_{g,n}^{\log}$, and the
amplitude has internal-edge count bounded by $d$.

\subsection*{4.5.\enspace The weight filtration}

Decompose the MC equation by internal-edge count,
\[
\Theta_\cA=\sum_{r\ge 1}\Theta^{[r]}_\cA,\qquad\Theta^{[r]}_\cA=\sum_{\substack{\Gamma\in\mathsf{Gr}^{\mathrm{st}}\\|E_{\mathrm{int}}(\Gamma)|=r}}\frac{W^{\log\mathrm{FM}}_\Gamma}{|\operatorname{Aut}(\Gamma)|}\otimes\Phi^\cA_\Gamma.
\]
Weight~$1$: $[D_{\mathrm{loc}},\Theta^{[1]}_\cA]=0$.  Weight~$N$:
$[D_{\mathrm{loc}},\Theta^{[N]}_\cA]+\tfrac12\sum_{i+j=N}[\Theta^{[i]}_\cA,\Theta^{[j]}_\cA]=0$.
The obstruction to extending from weight $N-1$ to weight $N$ lies in
$H^2(\gAmod/F^{N+1},D_{\mathrm{loc}})$.  The bar-intrinsic
construction solves all stages simultaneously.

\medskip
\noindent\textbf{Weight~$1$: boundary divisors.}\enspace
In degree $(0,4)$: the three boundary points of
$\overline{\cM}_{0,4}\cong\mathbb P^1$ correspond to the three
channels $(12|34)$, $(13|24)$, $(14|23)$,
\[
\Theta^{[1]}_\cA\big|_{(0,4)}=[\delta_s]\otimes\Phi_s+[\delta_t]\otimes\Phi_t+[\delta_u]\otimes\Phi_u.
\]
The weight-one equation $\Phi_s+\Phi_t+\Phi_u=F_3$ identifies the
Jacobiator homotopy with a boundary in the bar complex: the first
Borcherds identity.  In degree $(1,1)$: the single boundary divisor
$\delta_{\mathrm{irr}}$ gives
$\Theta^{[1]}_\cA|_{(1,1)}=[\delta_{\mathrm{irr}}]\otimes\Delta$ with
$\Delta=\sum_i e_i\otimes e^i$ the Casimir.  In degree $(0,5)$: the
ten boundary divisors of $\overline{\cM}_{0,5}$ (Petersen graph)
recover the pentagon identity.  The entire $A_\infty$ hierarchy is
the tree-level projection of a single MC element.

\medskip
\noindent\textbf{Weight~$2$: the first planted-forest layer.}\enspace
At weight~$2$, codimension-two strata enter: classical products of
two boundary divisors plus genuinely logarithmic planted-forest
strata with no DM counterpart.  The weight-two equation is
$[D_{\mathrm{loc}},\Theta^{[2]}_\cA]=-\tfrac12[\Theta^{[1]},\Theta^{[1]}]$,
and the planted-forest correction $\Theta^{[2]}_{\mathrm{pf}}$ records
how collision limits in $\mathrm{FM}_n(X|D)$ correct the naive square.
For the $\mathcal W_3$ central-parameter line, cubic gauge triviality
kills the weight-$2$ obstruction; the first genuine invariant appears
at weight~$3$.

\subsection*{4.6.\enspace The all-genus recursion}

Projecting the MC equation to fixed $(g,n)$,
\[
d\Theta^{(g,n)}_\cA+\tfrac12\!\!\!\sum_{\substack{g_1+g_2=g\\n_1+n_2=n+2}}\!\!\!\ell_2^{(0)}(\Theta^{(g_1,n_1)}_\cA,\Theta^{(g_2,n_2)}_\cA)+\Delta_{\mathrm{ns}}\Theta^{(g-1,n+2)}_\cA+\sum_{k\ge 3}\frac{1}{k!}\ell_k^{(\mathrm{tree})}(\Theta_\cA,\ldots,\Theta_\cA)^{(g,n)}=0.
\]
Four terms: the internal differential, separating clutching,
non-separating BV, and planted-forest corrections.  This recursion
determines $\Theta^{(g,n)}_\cA$ inductively from lower $(g',n')$ in
lexicographic order.  The base cases are
$\Theta^{(0,3)}_\cA$ (the chiral bracket) and $\Theta^{(0,4)}_\cA$
(the Jacobiator data).  The recursion terminates at each $(g,n)$.

\medskip
\noindent\textbf{The clutching identity.}\enspace
On each rigid boundary stratum~$\rho$,
\[
\Theta_\cA\big|_\rho=(\kappa_\rho)_*\!\left(\bigotimes_{v\in V(S_\rho)}\Theta_{\cA,v}\right).
\]
The amplitude on the degenerate curve equals the product of amplitudes
on the irreducible components, sewn by propagators along the nodes.
This is the MC-level sewing axiom, and the starting point for the
clutching law of the quartic shadow (\S6).

\bigskip
\noindent
Every invariant in this monograph is a projection of $\Theta_\cA$:
the curvature $\kappa$, the spectral discriminant $\Delta_\cA$, the
shadow metric $Q_L$, the sewing amplitudes, the $\hat A$-genus, the
Fredholm determinant, and the quartic contact invariant
$Q^{\mathrm{contact}}$.  The five main theorems are five projections.

\medskip
\begin{center}
\small
\renewcommand{\arraystretch}{1.15}
\begin{tabular}{lll}
\textbf{Projection of $\Theta_\cA$} & \textbf{Invariant} & \textbf{Section}\\
\hline
Scalar ($r=2$, genus tower) & $\kappa(\cA)$, $F_g$, $\hat A$-genus & \S2\\[1pt]
Binary genus-$0$ residue & classical $r$-matrix, Yangian seed & \S7.7\\[1pt]
Binary Verdier dual & complementarity $Q_g\oplus Q_g^!$ & \S2\\[1pt]
Bar-cobar counit & inversion $\Omega(\barB(\cA))\xrightarrow{\sim}\cA$ & \S1\\[1pt]
Chiral Hochschild & $\operatorname{ChirHoch}^*(\cA)$, deformation ring & \S9\\[1pt]
Spectral discriminant & $\Delta_\cA(x)$, genus-$1$ Hessian & \S5\\[1pt]
Cubic shadow ($r=3$) & gauge-trivial obstruction & \S6\\[1pt]
Quartic resonance ($r=4$) & $\mathfrak Q^{\mathrm{contact}}$, clutching law & \S6\\[1pt]
MC tautological descent & shadow CohFT, $R^*(\overline{\cM}_{g,n})$ & \S5\\[1pt]
Primitive kernel & shell equations, branch BV & \S5\\[1pt]
Dirichlet--sewing lift & $S_\cA(u)$, Euler--Koszul class & \S8\\[1pt]
Swiss-cheese colour decomposition & $\alpha_T$, six projections (Vol.~II) & \S10
\end{tabular}
\end{center}
\medskip

\noindent
These five theorems are five projections of $\Theta_\cA$.  Theorem~A
constructs the arena in which $\Theta_\cA$ lives.  Theorem~B
recovers $\cA$ from $\Theta_\cA$ on the Koszul locus.  Theorem~C
splits the genus tower into Lagrangian summands (C1 unconditional, C2
\textsc{(uniform-weight)}).  Theorem~D extracts the scalar $\kappa$
with $\hat A$-genus generating function (uniform-weight; multi-weight
adds $\delta F_g^{\mathrm{cross}}$).  Theorem~H identifies the
deformation ring via $\operatorname{ChirHoch}^*$ concentration.

The finite-order projections form an obstruction tower
$\Theta_\cA^{\le 2}\to\Theta_\cA^{\le 3}\to\Theta_\cA^{\le 4}\to\cdots$
whose generating function is algebraic of degree~$2$.  On each
primary line,
$H(t)=t^2\sqrt{Q_L(t)}$ is determined by three seed invariants
$(\kappa,\alpha,S_4)$.  The inverse limit exists and is automatically
Maurer--Cartan because $D_\cA^2=0$ (MC4).

In three dimensions, $\Theta_\cA$ extends to an open/closed
Maurer--Cartan element
$\Theta^{\mathrm{oc}}_\cA=\Theta_\cA+\sum_j\mu^{M_j}$
packaging the holomorphic-topological QFT on
$\mathbb C_z\times\mathbb R_t$: the bar complex encodes the closed
(holomorphic) colour, the coproduct encodes the open (topological)
colour, and the chiral derived center $\cZ^{\mathrm{der}}_{\mathrm{ch}}(\cA)$
is the universal bulk (Volume~II).

The bar complex thus organises five mathematical structures as
projections of one object: the genus expansion
$F_g=\kappa\cdot\lambda_g^{\mathrm{FP}}$ \textsc{(uniform-weight)} is
a GUE matrix model at $N^2=\kappa$; the ordered bar
$\barB^{\mathrm{ord}}(\cA)$ produces $U_q(\mathfrak g)$ via FRT
reconstruction; the quantum-group $R$-matrix yields knot invariants;
the shadow metric $Q_L$ defines a Bridgeland stability condition with
universal wall phase $\pi/4$; and the Stokes constants
$S_n/S_1=(-1)^{n+1}n$ encode the resurgent structure.

\medskip
\noindent\textbf{Corollary} (MC replaces the sum over topologies).
On the Koszul locus, the tree-level complex carries a contracting
homotopy (from the PBW degeneration), so $\Theta^{(g)}$ is determined
up to gauge-equivalent MC elements by induction from $\Theta^{(0)}$.
No free parameters appear at any genus; distinct choices of
contracting homotopy yield gauge-equivalent MC elements; no choice of
bulk topologies is required.  The genus expansion is computed
algebraically from the genus-$0$ chiral bracket by the clutching
recursion of \S4.6, not approximated perturbatively.  The
Maloney--Witten problem (negative Fourier coefficients in the
Poincar\'e series) does not arise: the MC equation determines each
$F_g$ uniquely by algebraic induction, without summing over bulk
topologies at fixed genus.  The non-perturbative sum over distinct
three-manifold topologies lies outside the MC framework.


% ====================================================================
\section*{5.\quad The modular tangent complex and Chern--Weil theory}
% ====================================================================

The MC element $\Theta_\cA$ exists unconditionally and satisfies the
MC equation.  Its deformation theory is controlled by the modular
tangent complex, obtained by twisting the convolution algebra by
$\Theta_\cA$ itself.

\subsection*{5.1.\enspace The $L_\infty$-twisted differential}

The MC element twists the convolution differential.  For $x\in\gAmod$,
\[
d_{\Theta_\cA}(x):=\sum_{n\ge 0}\sum_{g\ge 0}\frac{\hbar^g}{n!}\ell_{n+1}^{(g)}(\Theta_\cA^{\otimes n},x)=d(x)+[\Theta_\cA,x]+\tfrac{1}{2!}\ell_3^{(0)}(\Theta_\cA,\Theta_\cA,x)+\hbar\ell_1^{(1)}(x)+\cdots
\]
In the strict chart this collapses to $d_{\Theta_\cA}(x)=d(x)+[\Theta_\cA,x]$.
The identity $d_{\Theta_\cA}^2=0$ is the MC equation for $\Theta_\cA$.

\subsection*{5.2.\enspace The modular tangent complex}

The \emph{modular tangent complex at $\Theta_\cA$}:
\[
T^{\mathrm{mod}}_{\Theta_\cA}:=(\gAmod,d_{\Theta_\cA}).
\]
Its cohomology governs deformations: $H^0$ infinitesimal automorphisms
(stabiliser), $H^1$ first-order deformations (tangent space), $H^2$
obstructions.  $H^2=0$ makes the MC moduli smooth at $\Theta_\cA$;
$H^1=0$ makes $\Theta_\cA$ rigid.

For a one-parameter family $\cA_t$ varying in $k(t)$ or $c(t)$, the
tangent complex forms a local system over the parameter space and
the obstruction class traces a path in $H^2$.  The vanishing locus
is the moduli of liftable deformations.

\subsection*{5.3.\enspace Three characteristic projections}

Three successive projections of $\Theta_\cA$ capture progressively
finer structure: $\kappa(\cA)$ (leading scalar curvature),
$\Delta_\cA(x)$ (characteristic polynomial of a finite-rank operator),
$\mathfrak R^{\mathrm{mod}}_4(\cA)$ (first nonlinear characteristic
class).

\medskip
\noindent\textbf{The scalar curvature.}\enspace
\[
\kappa(\cA)=\operatorname{tr}(\Theta_\cA)_{2,0}=\operatorname{tr}(\pi\circ m_2\circ\iota^{\otimes 2}\circ\Delta_{\mathrm{sep}})\in\Bbbk.
\]
This is the first Chern class of the modular bar bundle:
$\operatorname{obs}_g(\cA)=\kappa(\cA)\cdot\lambda_g$
\textsc{(uniform-weight)}; multi-weight at $g\ge 2$ adds
$\delta F_g^{\mathrm{cross}}$ \textsc{(all-weight)}.  The Mumford relation
$\lambda_g^{g+1}=0$ confines genus-$g$ scalar obstructions to
dimension at most~$g$.

Explicit values across the standard landscape:
\begin{center}
\renewcommand{\arraystretch}{1.2}
\begin{tabular}{lccl}
\textbf{Algebra $\cA$} & $\kappa(\cA)$ & $\kappa(\cA^!)$ & \textbf{Sum rule}\\
\hline
Heisenberg $\cH_k$ & $k$ & $-k$ & $0$\\[1pt]
Affine $\widehat{\fg}_k$ & $\tfrac{(k+h^\vee)\dim\fg}{2h^\vee}$ & $-\tfrac{(k+h^\vee)\dim\fg}{2h^\vee}$ & $0$\\[4pt]
Virasoro $\mathrm{Vir}_c$ & $c/2$ & $(26-c)/2$ & $13$\\[1pt]
$\beta\gamma$ system ($\lambda{=}1$) & $1$ & $-1$ & $0$\\[1pt]
Lattice $V_\Lambda$ & $\operatorname{rank}(\Lambda)$ & $-\operatorname{rank}(\Lambda)$ & $0$\\[1pt]
$\mathcal W_N(\widehat{\mathfrak{sl}}_N)$ & $c\cdot(H_N-1)$ & $c'\cdot(H_N-1)$ & $K_N\cdot(H_N-1)$\\[4pt]
Free fermion $\psi$ & $1/4$ & $-1/4$ & $0$
\end{tabular}
\end{center}
$H_N=\sum_{j=1}^N 1/j$ is the harmonic number
($H_N-1\neq H_{N-1}$: at $N=2$, $H_2-1=1/2$ but $H_1=1$).  The sum rule
$\kappa(\cA)+\kappa(\cA^!)=0$ holds for Kac--Moody, free fields,
lattice algebras.  For $\mathcal W$-algebras the sum is a nonzero
constant: $\kappa(\mathrm{Vir}_c)+\kappa(\mathrm{Vir}_{26-c})=13$ and
$\kappa(\mathcal W_{N,c})+\kappa(\mathcal W_{N,c'})=K_N(H_N-1)$ with
$K_N=c+c'$.  Self-duality occurs at $c=13$ for Virasoro and $c=50$
for $\mathcal W_3$.

\medskip
\noindent\textbf{The spectral discriminant.}\enspace
The Koszul dual Hilbert series
\[
h_{\cA^!}(q)=\sum_{n\ge 0}\dim H^n(\barB(\cA))\,q^n
\]
is rational for every standard family: it satisfies a holonomic
recurrence of finite order $d$, and the transfer matrix
$T_{\mathrm{rec}}$ is a $d\times d$ matrix.  Define the
\emph{spectral branch object} as a finite-rank perfect complex
$V^{\mathrm{br}}_\cA$ with \emph{branch transport operator}
$T_{\mathrm{br},\cA}\colon V^{\mathrm{br}}_\cA\to V^{\mathrm{br}}_\cA$
and
\begin{equation}\label{eq:pref-delta-def}
\Delta_\cA(x):=\operatorname{sdet}(1-xT_{\mathrm{br},\cA}).
\end{equation}
The rank of $V^{\mathrm{br}}_\cA$ equals the degree of algebraicity of
$h_{\cA^!}$.  Three properties: $\Delta_\cA$ is self-dual
($\Delta_{\cA^!}=\Delta_\cA$), unlike $\kappa$ which is anti-symmetric
for KM/free; exact factorization functors preserving quadratic OPE
data preserve $T_{\mathrm{br}}$; the spectral radius equals
$\dim\mathfrak g$ for affine KM at generic level.

For $\mathfrak{sl}_3$: the degree-$3$ recurrence has transfer matrix
eigenvalues $8$, $(3\pm\sqrt{13})/2$, and
$\Delta_{\widehat{\mathfrak{sl}}_3}(x)=(1-8x)(1-3x-x^2)$.  For
Heisenberg: $h_{\mathcal H_k^!}(q)=1+q$ and $\Delta_{\mathcal H_k}(x)=1$.
For $\beta\gamma$: $h=\sqrt{(1+q)/(1-3q)}$ and
$\Delta_{\beta\gamma}(x)=(1-3x)(1+x)$.

In the Chern--Weil dictionary, $\Delta_\cA$ plays the role of the
Chern character:
$\operatorname{ch}_r(\cA)=\tfrac{1}{r!}\tfrac{d^r}{dt^r}|_{t=0}[-\log\Delta_\cA(t)]=\tfrac{1}{r}\operatorname{tr}(T_{\mathrm{br}}^r)$.

\medskip
\noindent\textbf{The quartic shadow.}\enspace
\[
\mathfrak R^{\mathrm{mod}}_{4,g,n}(\cA)=\operatorname{pr}_{4,g,n}(\Theta_\cA)\in H^*(\overline{\cM}_{g,n})\otimes(\bar\cA^\vee)^{\otimes n}_{\mathrm{cyc}}.
\]
This is the first \emph{nonlinear} characteristic class: the first
invariant not recoverable from $\kappa$ and $\Delta_\cA$.  Its
genus-zero, four-point specialisation is the quartic contact invariant
\[
Q^{\mathrm{contact}}_{\mathrm{Vir}}=\frac{10}{c(5c+22)},\qquad Q^{\mathrm{contact}}_{\cH}=0,\qquad Q^{\mathrm{contact}}_{\beta\gamma}\neq 0.
\]
(Verify: at $c=1$,
$Q_{\mathrm{Vir}}=10/27$.)  For Heisenberg the
quartic vanishes (no cubic OPE pole).  For $\beta\gamma$ the quartic
is the first gauge-nontrivial shadow (the cubic is gauge-removable).
For Virasoro, both cubic and quartic are nonvanishing.

\subsection*{5.4.\enspace The Chern--Weil dictionary}

The analogy between $\Theta_\cA$ and a connection $1$-form is exact.
The modular bar bundle
$\mathrm B^{\log}_{\mathrm{mod}}(\cA)\to\overline{\cM}_{g,n}$ is a
principal bundle for the gauge group of the convolution algebra;
$\Theta_\cA$ is its flat connection (flatness automatic); the shadow
tower is its Chern--Weil theory.  Holonomy becomes factorisation
homology, gauge equivalence becomes MC gauge, and the Weil
homomorphism becomes the shadow algebra
$\cA^{\mathrm{sh}}=H_\bullet(\Defcyc^{\mathrm{mod}})$.

\begin{center}
\renewcommand{\arraystretch}{1.25}
\begin{tabular}{ll}
\textbf{Classical Chern--Weil} & \textbf{Modular factorization}\\
\hline
Principal $G$-bundle $P\to M$ & Modular bar bundle $\mathrm B^{\log}_{\mathrm{mod}}(\cA)\to\overline{\cM}_{g,n}$\\[1pt]
Connection $\omega\in\Omega^1(P,\fg)$ & MC element $\Theta_\cA\in\operatorname{MC}(\gAmod)$\\[1pt]
Curvature $\Omega=d\omega+\tfrac12[\omega,\omega]$ & $D\Theta_\cA+\tfrac12[\Theta_\cA,\Theta_\cA]=0$\\[1pt]
Flatness $\Omega=0$ & Automatic (bar-intrinsic)\\[1pt]
First Chern $c_1$ & $\kappa(\cA)=\operatorname{tr}(\Theta_\cA)_{2,0}$\\[1pt]
Chern character & Spectral discriminant $\Delta_\cA(t)$\\[1pt]
Higher Chern classes $c_r$ & Shadows $\mathfrak R^{\mathrm{mod}}_{4,g,n}$, $\operatorname{Sh}_r$ ($r\ge 5$)\\[1pt]
Weil homomorphism & Shadow algebra $\cA^{\mathrm{sh}}=H_\bullet(\Defcyc^{\mathrm{mod}})$\\[1pt]
$\operatorname{ad}$-invariant polynomial & Cyclic cocycle in $\Defcyc^{\mathrm{mod}}(\cA)$\\[1pt]
Cyclic pairing $\operatorname{tr}(XY)$ & Verdier duality on $\operatorname{Ran}(X)$\\[1pt]
Group multiplication & Operadic composition in $\cP^{\mathrm{ch}}$\\[1pt]
Coproduct on $\cO(G)$ & Operadic cocomposition / clutching\\[1pt]
Lagrangian splitting $T^*G=\fg\oplus\fg^*$ & Complementarity $Q_g(\cA)\oplus Q_g(\cA^!)$\\[1pt]
Holonomy $\operatorname{Hol}(\gamma)$ & Factorisation homology $\int_\Sigma\cA$\\[1pt]
Gauge equivalence & MC gauge in $\operatorname{MC}_\bullet(\gAmod)$\\[1pt]
Transgression $\operatorname{CS}(\omega)$ & Shadow obstruction tower (\S6)\\[1pt]
Flat sections $\nabla s=0$ & Twisted cochains $d_{\Theta_\cA}(x)=0$
\end{tabular}
\end{center}

\subsection*{5.5.\enspace The genus spectral sequence}

The genus filtration $F^g L_{\mathrm{mod}}:=\bigoplus_{g'\ge g}(\gAmod)_{g',*,*}$
yields a convergent spectral sequence.  Write
$D_0=D_{\mathrm{int}}+D_{\mathrm{sep}}$ for the genus-preserving part
and $D_1=D_{\mathrm{nsep}}$ for the genus-raising part; these satisfy
$D_0^2=0$, $\{D_0,D_1\}=0$, $D_1^2=0$ (a bicomplex).

The spectral sequence: $E_0^{p,q}=\operatorname{gr}_F^p L_{\mathrm{mod}}^{p+q}$
with $d_0=[D_0,-]$; $E_1^{p,*}$ is the associated graded cohomology,
genus by genus; the differential $d_r$ at page $r$ encodes the
obstruction $\operatorname{Ob}_r$; $E_\infty^{p,q}\cong\operatorname{gr}_F^p H^{p+q}(L_{\mathrm{mod}})$.
Convergence follows from completeness of the depth filtration and
pronilpotence of the weight filtration; the $\varprojlim^1$ term
vanishes by Mittag-Leffler.

This spectral sequence is distinct from the PBW spectral sequence
(which decomposes by conformal weight within fixed genus).  PBW
computes \emph{what} bar cohomology is; the genus spectral sequence
determines \emph{whether} the MC element extends across genera.  For
Heisenberg it degenerates at $E_1$.  For Virasoro, $d_1\neq 0$ and
nontrivial differentials persist at every page.

\subsection*{5.6.\enspace Homotopy invariance and functoriality}

The shadow algebra $\cA^{\mathrm{sh}}=H_\bullet(\Defcyc^{\mathrm{mod}}(\cA))$
is a homotopy invariant.  For any $\infty_\alpha$-quasi-isomorphism
$f\colon\cA\xrightarrow{\sim}\cA'$,
\[
f_*\colon\operatorname{MC}_\bullet(\gAmod)\xrightarrow{\;\sim\;}\operatorname{MC}_\bullet(\mathfrak g^{\mathrm{mod}}_{\cA'})
\]
is a weak homotopy equivalence.  The scalar $\kappa$, the spectral
determinant $\Delta_\cA$, and the full shadow obstruction tower
$(\kappa,\Delta,\mathfrak C,\mathfrak Q,\operatorname{Sh}_5,\ldots)$
consist of quasi-isomorphism invariants; $\Theta_\cA$ itself is a
quasi-isomorphism invariant up to gauge.  The assignment
$\cA\mapsto(\gAmod,\Theta_\cA)$ is a functor from chiral algebras with
$\infty_\alpha$-quasi-isos inverted to MC moduli problems.  The
convolution $L_\infty$-algebra is functorial in each slot separately
but not as a bifunctor; functoriality of $\cA\mapsto\gAmod$ proceeds
through the source slot.

\subsection*{5.7.\enspace The shadow CohFT and topological recursion}

Projected to $\overline{\cM}_{g,n}$, the MC equation produces
tautological classes and tautological relations.  The MC element
determines shadow tautological classes
$\tau_{g,n}(\cA)\in R^*(\overline{\cM}_{g,n+1})$ assembling into a
cohomological field theory (Theorem~\ref{thm:shadow-cohft}): the
factorisation axiom holds by the clutching law, the
$\Sigma_n$-equivariance by the symmetry of stable-graph summation,
the flat identity conditionally on the vacuum lying in~$V$.

Projecting the MC equation to the $(g,n)$-component yields a
tautological relation in $R^*(\overline{\cM}_{g,n})$ at each $(g,n)$:
\[
\sum_{\mathrm{sep}/\mathrm{nsep}}\xi_{\mathrm{sep},*}\bigl(\tau_{g_1,|S|}(\cA)\cdot\tau_{g_2,|S^c|}(\cA)\bigr)+\delta_{\mathrm{pf}}^{(g,n)}=0.
\]

\medskip
\noindent\textbf{WDVV and Mumford from MC.}\enspace
At genus~$0$, arity~$3$: the MC tautological relation on
$\overline{\cM}_{0,4}$ reduces to WDVV.  At arity~$2$, genus~$g\ge 2$:
substituting $\tau_{g,2}(\cA)=\kappa\cdot\pi^*\lambda_g$
\textsc{(uniform-weight)} recovers the Mumford $\lambda$-class
clutching relation on $\overline{\cM}_g$.

The complementarity propagator
$P_\cA\colon H^*(\cA)\otimes H^*(\cA^!)\to\mathbb C$ is the Givental
$R$-matrix of the shadow CohFT; when it carries a flat unit (all
standard families with vacuum in $V$), Teleman reconstruction
recovers $\{\tau_{g,n}\}$ at all $(g,n)$ from genus-$0$ data and
$P_\cA$.  Eynard--Orantin topological recursion is the MC equation
applied to the shadow obstruction tower.

\begin{center}
\small
\renewcommand{\arraystretch}{1.15}
\begin{tabular}{ll}
\textbf{MC datum} & \textbf{Tautological output}\\
\hline
$\Theta_\cA\in\operatorname{MC}(\gAmod)$ & shadow CohFT $\Omega_{g,n}^\cA$\\[1pt]
$D^2=0$ on $\mathrm{FM}^{\log}$ & CohFT axioms\\[1pt]
$D\Theta+\tfrac12[\Theta,\Theta]=0$ at $(g,n)$ & tautological relation\\[1pt]
$[\Theta,\Theta]|_{\mathrm{sep}}$ & WDVV / associativity\\[1pt]
$\hbar\Delta(\Theta)|_{\mathrm{nsep}}$ & Mumford $\lambda$-class relations\\[1pt]
$P_\cA$ & Givental $R$-matrix\\[1pt]
MC recursion & EO topological recursion
\end{tabular}
\end{center}

\medskip
\noindent\textbf{The Pixton shadow bridge.}\enspace
For class-M algebras, $\Theta_\cA$ generates tautological classes at
every genus.  At genus~$2$--$3$, the relations
$\tau_{g,n}(\mathrm{Vir}_c)$ for generic $c$ produce elements of the
Pixton ideal in $R^*(\overline{\cM}_{g,n})$: Pixton's tautological
ring relations arise as shadows of the Virasoro MC equation.  On the
semisimple locus, MC-descended relations together with Mumford
relations generate the Pixton ideal
(Theorem~\ref{thm:pixton-from-mc-semisimple}); the non-semisimple
frontier remains open.

\medskip
\noindent\textbf{The B\"ocherer bridge.}\enspace
At genus~$2$, the shadow tautological class $\tau_{2,0}(\mathrm{Vir}_c)$
for lattice vertex algebras computes a central $L$-value via the
B\"ocherer conjecture (Furusawa--Morimoto).  For the Leech lattice,
$\chi_{12}=\operatorname{SK}(f_{22})$, and the MC-derived second
Fourier coefficient $c_2\approx-1.918\times 10^{-6}\neq 0$ is the
first nonzero central $L$-value produced by the shadow obstruction
tower.

\subsection*{5.8.\enspace The primitive kernel and cofree reduction}

$\Theta_\cA$ lives in the cofree conilpotent coalgebra underlying the
modular bar construction.  By Milnor--Moore for modular operads,
$\Theta_\cA$ is determined by its restriction to primitive corollas
and rigid cutting strata, the \emph{primitive logarithmic modular
kernel}
\[
\mathfrak K_\cA=\sum_{g\ge 0,n\ge 2}K_{g,n}^\cA+\sum_{\rho\in\mathrm{Rig}}K_\rho^\cA,
\]
and the equivalence (Theorem~\ref{thm:primitive-to-global-reconstruction})
\[
(Q_\cA^{\mathrm{mod}})^2=0\iff d\mathfrak K_\cA+\mathfrak K_\cA\star\mathfrak K_\cA+\hbar\Delta_{\mathrm{ns}}(\mathfrak K_\cA)=0,
\]
with $\star$ the primitive pre-Lie product from log-FM residues,
replaces the genus-by-genus obstruction programme by a single equation
on primitive data.

\medskip
\noindent\textbf{Shell equations.}\enspace
Projecting to genus~$2$ and~$3$:
\begin{align*}
R_2(\mathfrak K_\cA)&=\Delta_{\mathrm{ns}}(K_1)+\tfrac12[K_1,K_1]+\sum_{\rho\in\mathrm{Rig}_2}R_\rho^\cA(K_{0,2}^\cA),\\
R_3(\mathfrak K_\cA)&=\Delta_{\mathrm{ns}}(K_2)+[K_1,K_2]+\tfrac{1}{3!}\ell_3(K_1,K_1,K_1)+\sum_\rho R_\rho^\cA(\ldots).
\end{align*}
For Heisenberg, only $\Delta_{\mathrm{ns}}$ survives.  For affine
Kac--Moody, $[K_1,K_1]$ contributes the cubic shadow.  For
$\beta\gamma$, cubic gauge triviality kills the bracket and the rigid
stratum contributes the quartic contact invariant.  For Virasoro and
$\mathcal W_N$, all three terms are active at every genus.

\medskip
\noindent\textbf{The branch BV action.}\enspace
The finite-rank branch-active quotient
$V_\cA^{\mathrm{br}}=V_\cA^+\oplus V_\cA^-$ carries a BV master action
\[
S_\cA^{\mathrm{br}}(x,\xi)=\tfrac12\Omega_\cA(K_\cA^{\mathrm{br}}x,x)+\sum_{m+n\ge 3}\frac{1}{m!n!}\mu_{m,n}^\cA(x^{\otimes m};\xi^{\otimes n}),
\]
satisfying
$\hbar\Delta_\cA^{\mathrm{br}}S_\cA^{\mathrm{br}}+\tfrac12\{S_\cA^{\mathrm{br}},S_\cA^{\mathrm{br}}\}_{\mathrm{BV}}=0$.
The metaplectic half-density $\delta_\cA\in 1+x\Bbbk[[x]]$ satisfies
$\delta_\cA^2=\Delta_\cA=\det(1-xT_\cA^{\mathrm{br,red}})$: the
spectral discriminant is the square of a canonical half-density.

\subsection*{5.9.\enspace Entanglement from the modular characteristic}

The bar coalgebra carries a natural bipartite structure: its coproduct
splits a configuration of points into two halves, one on each side of
a spatial cut.  For a spatial bipartition of an interval of length $L$
with UV cutoff $\varepsilon$, the scalar projection of $\Theta_\cA$
restricted to the bipartite cut yields
\[
S_{\mathrm{EE}}(\cA)=\frac{c(\cA)}{3}\log\frac{L}{\varepsilon}.
\]
For Virasoro $c=2\kappa$; the relationship between $c$ and $\kappa$
is family-dependent.  Complementarity gives the entanglement sum rule
$S_{\mathrm{EE}}(\cA)+S_{\mathrm{EE}}(\cA^!)=\tfrac{c(\cA)+c(\cA^!)}{3}\log(L/\varepsilon)$;
for the Virasoro family $\mathrm{Vir}_c^!=\mathrm{Vir}_{26-c}$ this
gives $c+c'=26$ and the universal sum $26/3\cdot\log(L/\varepsilon)$.

The four shadow-depth classes (G/L/C/M) classify entanglement
complexity: class G (Gaussian, $r_{\max}=2$) exhibits area-law
entanglement only; class L ($r_{\max}=3$) adds subleading logarithmic
corrections from the cubic shadow; class C ($r_{\max}=4$) adds a
contact correction from the quartic invariant; class M
($r_{\max}=\infty$) carries an infinite tower of corrections.

\medskip
\noindent\textbf{The entanglement programme (G11--G16).}\enspace
Six theorem targets extend the scalar result.  G11: Ryu--Takayanagi
from the scalar projection of $\Theta_\cA$.  G12: Koszulness as exact
quantum error correction; the Knill--Laflamme condition is the
Lagrangian isotropy condition from the complementarity pairing;
shadow depth equals the number of redundancy channels; the twelve
Koszulness characterisations translate into a twelve-fold code
dictionary; failure of Koszulness is failure of the code.  G13:
replica trick from the MC equation.  G14: modular flow from the
tangent complex.  G15: QES as the shadow connection Ward identity.
G16: Page curve at $c=13$ from the shadow connection at the self-dual
point; the scrambling time is controlled by $\Delta_\cA$.


% ====================================================================
\section*{6.\quad The shadow obstruction tower}
% ====================================================================

\subsection*{6.1.\enspace The extension tower}

Can $\Theta_\cA$ be approximated by finitely many arities, and what
controls the passage from one finite approximation to the next?  The
total weight filtration answers both.

Define the \emph{total weight} $w(g,r,d):=2g-2+r+d$.  The filtration
$F^N\mathfrak g^{\mathrm{amb}}_\cA:=\bigoplus_{w\ge N}(\mathfrak g^{\mathrm{amb}}_\cA)_{g,r,d}$
is exhaustive, separated, complete.  $\mathfrak g^{\mathrm{amb}}_\cA$
is the \emph{ambient} (planted-forest-extended) modular convolution
algebra.  The \emph{modular extension tower}:
\[
\cE_\cA(N):=\operatorname{MC}(\mathfrak g^{\mathrm{amb}}_\cA/F^{N+1}),\qquad N\ge 1,
\]
with restriction maps $\cE_\cA(N+1)\to\cE_\cA(N)$.  A full lift is a
compatible point of the inverse system; the bar-intrinsic construction
provides one unconditionally.  The shadow obstruction tower studies
$\cE_\cA(1)\leftarrow\cE_\cA(2)\leftarrow\cE_\cA(3)\leftarrow\cdots$.

\subsection*{6.2.\enspace Obstruction classes and the all-arity master equation}

Define $\Theta_\cA^{\le r}$ by projecting to total weight $\le r$.  The
\emph{shadow obstruction tower} is the sequence
$\Theta_\cA^{\le 2},\Theta_\cA^{\le 3},\Theta_\cA^{\le 4},\ldots$  At
each truncation the MC equation fails by a remainder,
\[
D\Theta^{\le r}_\cA+\tfrac12[\Theta^{\le r}_\cA,\Theta^{\le r}_\cA]=-o_{r+1}(\cA),
\]
with obstruction class $o_{r+1}(\cA)\in Z^2(F^{r+1}/F^{r+2},d_{\le r})$.
Extension is possible iff $[o_{r+1}]=0$ in
$H^2(F^{r+1}/F^{r+2},d_{\le r})$.  The bar-intrinsic construction
guarantees extension, so the obstruction class is always exact; the
shadow obstruction tower records the gauge choices at each stage.

The all-arity master equation at arity~$r$:
\[
\nabla_H(\operatorname{Sh}_r(\cA))+o^{(r)}(\cA)=0,\qquad r\ge 3,
\]
with $\nabla_H$ the Hodge covariant derivative.  This equation
determines $\operatorname{Sh}_r$ uniquely modulo gauge if
$H^1(F^r/F^{r+1},d_{\le r-1})=0$.

\subsection*{6.3.\enspace Cubic gauge triviality and the canonical quartic}

\begin{quote}
\textbf{Cubic gauge triviality.}\enspace
When $H^1(F^3/F^4,d_2)=0$, the cubic MC term $\Theta^{[3]}_\cA$ is
gauge-trivial: there exists $\exp(\xi_3)$ with $\xi_3\in F^3$ such that
$\exp(\xi_3)\cdot\Theta^{\le 3}_\cA$ has vanishing weight-$3$
component.
\end{quote}

After this gauge transformation, the quartic obstruction class
$[\mathfrak Q(\cA)]\in H^2(F^4/F^5,d_2)$ is canonical: independent of
homotopy retraction data, cubic gauge, and MC gauge orbit
representative.  This is the first characteristic class not polynomial
in $\kappa$.  The mechanism is standard Postnikov: trivial $H^1$ at
level $n$ makes the $H^2$ class at level $n+1$ canonical.

The clutching law for the quartic follows from Mok's log FM
degeneration formula.  For $C\leadsto C_0=C_1\cup_p C_2$,
\[
\mathfrak Q(\cA)|_{C_0}=\sum_{a+b=4}\mathfrak Q^{(a)}(\cA)|_{C_1}\cdot\mathfrak Q^{(b)}(\cA)|_{C_2}+\text{planted-forest correction}.
\]
The correction is genuinely logarithmic, with no DM counterpart.

The quartic contact invariant
$Q^{\mathrm{contact}}_{\mathrm{Vir}}=10/[c(5c+22)]\in H^2(F^4/F^5,d_2)$.
The denominator has physical meaning: $c=0$ is the logarithmic
threshold; $c=-22/5$ is the minimal model accumulation point
$\mathcal M(2,5)$.  The numerator $10=\binom{5}{2}$ counts arity-four
Catalan channels.

\subsection*{6.4.\enspace The Hamilton--Jacobi generating function}

For a multi-generator algebra ($\mathcal W_3$ with generators $T$
(weight~$2$) and $W$ (weight~$3$)), the deformation space is
two-dimensional.  The Riccati ODE lifts to a Hamilton--Jacobi PDE.
Write $U=\sum_{s\ge 3}\operatorname{Sh}_s(x_T,x_W)$; the master equation
becomes
\[
2E(U)+\tfrac12\lVert\nabla U\rVert^2_H=R(x_T,x_W),
\]
with Euler operator $E=x_T\partial_T+x_W\partial_W$, kinetic energy
$\lVert\nabla U\rVert^2_H=(2/c)(\partial_T U)^2+(3/c)(\partial_W U)^2$,
and $R$ a polynomial of degree $\le 4$ in
$(\kappa,\mathfrak C,\mathfrak Q)$.  The gradient $dU\colon X\to T^*X$
embeds the deformation space as a Lagrangian submanifold of $T^*X$.
Whether the genus expansion admits a quantisation with $U$ as leading
semiclassical term is open: it would require a quantum Hamiltonian
whose WKB expansion has leading term $U$ and whose quantum corrections
reproduce the genus-$g$ shadow amplitudes.

On the $T$-line the conformal weight forces the transverse momentum to
vanish and the PDE collapses to the Riccati ODE.  On the $W$-line the
cubic shadow $\operatorname{Sh}_3=2x_T^3+3x_Tx_W^2$ has nonzero normal
derivative $(\partial_T\operatorname{Sh}_3)|_{x_T=0}=3x_W^2$, and the
first inter-channel coupling correction enters at arity~$6$:
\[
\delta_6^{W\text{-line}}=\frac{576(5c+38)}{c^3(5c+22)^3}.
\]
The $T$-line is autonomous; the $W$-line is not.  The asymmetry is
intrinsic to the weight difference between $T$ (weight~$2$) and $W$
(weight~$3$).

\subsection*{6.5.\enspace The four depth classes}

The \emph{shadow depth} $r_{\max}(\cA)$ is the smallest $r$ such that
$o_{r+1}(\cA)=0$ for all subsequent levels, or $\infty$ if the tower
never terminates.  The standard landscape partitions into four depth
classes; no other values occur on any one-dimensional primary slice.
The single-line dichotomy: the shadow metric $Q_L(t)$ is quadratic in
$(\kappa,\alpha,S_4)$, and its factorisation type (perfect square or
irreducible) determines whether the tower terminates.

\begin{center}
\renewcommand{\arraystretch}{1.3}
\begin{tabular}{lcccll}
\textbf{Class} & \textbf{Symbol} & $r_{\max}$ & $\ell_3^{(0)},\ell_4^{(0)}$ & \textbf{Potential} & \textbf{Examples}\\
\hline
Gaussian & G & $2$ & $0,0$ & quadratic & Heisenberg $\cH_k$, free fermion $\psi$\\
Lie/tree & L & $3$ & $\neq 0,0$ & cubic & Affine $\widehat{\mathfrak{sl}}_N$, currents\\
Contact & C & $4$ & $0,\neq 0$ & quartic & $\beta\gamma$, $bc$ ghosts\\
Mixed & M & $\infty$ & $\neq 0,\neq 0$ & non-polynomial & Virasoro, $\mathcal W_N$, $\mathcal W_\infty$
\end{tabular}
\end{center}

\noindent
\textbf{Gaussian} ($r_{\max}=2$): transferred brackets $\ell_k^{(0)}$
vanish for $k\ge 3$ on bar cohomology, so the tower terminates at
$\kappa$.  Two mechanisms: Heisenberg has no simple pole; the free
fermion has a simple pole but fermionic antisymmetry forces $S_3=0$.
Sewing is by Fredholm determinant (Heisenberg) or Pfaffian (fermion).
\textbf{Lie/tree} ($r_{\max}=3$): the OPE has a simple pole generating
a Lie bracket, but $\ell_4^{(0)}$ vanishes on bar cohomology.
\textbf{Contact} ($r_{\max}=4$): the simple-pole bracket vanishes on
bar cohomology, but the quaternary bracket survives; the first
nontrivial shadow is $\mathfrak Q(\cA)$.  \textbf{Mixed}
($r_{\max}=\infty$): both cubic and quartic shadows are nontrivial,
and their Poisson bracket $\{\mathfrak C,\mathfrak Q\}_H$ generates a
nonvanishing quintic, then sextic, and so on.

\medskip
\noindent\textbf{Shadow depth is not Koszulness.}\enspace
All four depth classes contain chirally Koszul algebras.  Heisenberg
(depth~$2$), affine at generic level (depth~$3$), $\beta\gamma$
(depth~$4$), and Virasoro (depth~$\infty$) are all chirally Koszul.
Shadow depth classifies the complexity of Koszul algebras, not
Koszulness itself.

The \emph{operadic complexity theorem}
(Theorem~\ref{thm:operadic-complexity-detailed}):
\[
r_{\max}(\cA)=A_\infty\text{-depth}(\cA)=L_\infty\text{-formality level}(\cA).
\]
$A_\infty$-depth is the largest $r$ with $m_r\neq 0$ in the minimal
$A_\infty$ model; $L_\infty$-formality level is the largest $r$ with a
nontrivial $r$-ary bracket on bar cohomology not removable by
$L_\infty$-isomorphism.

\subsection*{6.6.\enspace The genus--arity bigrading}

The MC equation decomposes along two orthogonal axes.  The \emph{genus
axis} extends $\Theta_\cA$ from genus~$0$ to~$1$ to~$2$, controlled by
the genus spectral sequence.  The \emph{arity axis} extends
$\Theta_\cA^{\le 2}\to\Theta_\cA^{\le 3}\to\Theta_\cA^{\le 4}$,
controlled by cyclic obstruction classes.  At fixed arity~$2$ the
genus tower is governed by $\kappa$ alone \textsc{(uniform-weight)};
at fixed genus~$0$ each new shadow coefficient introduces independent
data.  The axes couple through the visibility formula
$g_{\min}(S_r)=\lfloor r/2\rfloor+1$: $S_r$ first contributes at
genus $\lfloor r/2\rfloor+1$ and two new shadow coefficients enter at
each genus.

\subsection*{6.7.\enspace Genus-two shells}

Genus~$2$ is the first level where all three boundary strata interact:
\[
\Theta^{(2)}_\cA=\underbrace{\Theta^{(2)}_{\mathrm{loop}\circ\mathrm{loop}}}_{\text{iterated BV}}+\underbrace{\Theta^{(2)}_{\mathrm{sep}\circ\mathrm{loop}}}_{\text{clutching $\times$ BV}}+\underbrace{\Theta^{(2)}_{\mathrm{pf}}}_{\text{planted-forest}}.
\]
The iterated-BV term involves the figure-eight graph:
\[
\Theta^{(2)}_{\mathrm{loop}^2}=\tfrac12\operatorname{tr}_1\operatorname{tr}_2\bigl(P_\cA^{\otimes 2}\circ m^{(h)}_{0,4}\circ\Delta^{\otimes 2}\bigr).
\]
The separating-times-loop term:
\[
\Theta^{(2)}_{\mathrm{sep}\circ\mathrm{loop}}=\operatorname{tr}_{\mathrm{ns}}\bigl(P_\cA\circ\ell_2^{(0)}(\Theta^{(1,\bullet)},\Theta^{(1,\bullet)})\bigr).
\]
The planted-forest term is the genuinely logarithmic contribution,
present only when cubic or quartic shadows are nontrivial.

\medskip
\noindent
Family-by-family structure:

\begin{center}
\renewcommand{\arraystretch}{1.2}
\begin{tabular}{lcccc}
& $\Theta^{(2)}_{\mathrm{loop}^2}$ & $\Theta^{(2)}_{\mathrm{sep}\circ\mathrm{loop}}$ & $\Theta^{(2)}_{\mathrm{pf}}$ & \textbf{Depth}\\
\hline
Heisenberg $\cH_k$ & $\checkmark$ & --- & --- & G\\
Affine $\widehat{\fg}_k$ & $\checkmark$ & $\checkmark$ & $\checkmark$ & L\\
$\beta\gamma$ system & $\checkmark$ & --- & --- & C\\
Virasoro $\mathrm{Vir}_c$ & $\checkmark$ & $\checkmark$ & $\checkmark$ & M\\
$\mathcal W_N$ & $\checkmark$ & $\checkmark$ & $\checkmark$ & M
\end{tabular}
\end{center}

Heisenberg: only iterated BV.  Affine: adds separating-times-loop.
Virasoro and $\mathcal W_N$: all three active, with planted-forest
carrying the quartic resonance.  $\beta\gamma$: only iterated BV
survives; the separating term vanishes because the Lie bracket is
trivial on bar cohomology, and the planted-forest vanishes because
the planted-forest coupling $\mu^{\mathrm{pf}}_{\beta\gamma}=0$
(the planted-forest stratum requires coupling between cubic and quartic
shadows, and the cubic $\mathfrak{C}_{\beta\gamma}=0$).  Contact termination via the iterated BV shell, not
Gaussian collapse.

\subsection*{6.8.\enspace The genus loop operator}

The \emph{genus loop operator}
$\Lambda_P\colon\operatorname{Sym}^r(\bar\cA^\vee)_{\mathrm{cyc}}\to\operatorname{Sym}^{r-2}(\bar\cA^\vee)_{\mathrm{cyc}}$
contracts two legs with the propagator:
\[
\Lambda_P(\phi)(a_1,\ldots,a_{r-2})=\sum_i\phi(a_1,\ldots,a_{r-2},e_i,e^i).
\]
This raises genus by~$1$ and lowers arity by~$2$: the shadow-level
incarnation of $\Delta_{\mathrm{ns}}$.

\medskip
\noindent\textbf{Explicit Virasoro computation.}\enspace
On the Virasoro primary line, the conformal-weight-$2$ bar cohomology
is one-dimensional, spanned by the class $x$ dual to $T$.  The Killing
form is $H=(c/2)x^2$, the propagator $P=2/c$, and the quartic contact
invariant is $Q^{\mathrm{contact}}_{\mathrm{Vir}}\cdot x^4$ with
$Q^{\mathrm{contact}}_{\mathrm{Vir}}=10/[c(5c+22)]$.  Applying
$\Lambda_P$:
\[
\Lambda_P\!\left(\frac{10}{c(5c+22)}\cdot x^4\right)=\binom{4}{2}\cdot\frac{2}{c}\cdot\frac{10}{c(5c+22)}\cdot x^2=\frac{120}{c^2(5c+22)}x^2,
\]
the genus-one Hessian correction.  The loop ratio
\[
\rho^{(1)}_{\mathrm{Vir}}=\frac{\Lambda_P(Q\cdot x^4)}{\kappa\cdot x^2}=\frac{120/[c^2(5c+22)]}{c/2}=\frac{240}{c^3(5c+22)}
\]
is not determined by $\kappa=c/2$: it detects the quartic structure
through its genus-one projection.  Asymptotically
$\rho^{(1)}\sim 48/c^4$ for $c\to\infty$.

\subsection*{6.9.\enspace Quintic forcing and infinite towers}

The quintic obstruction determines whether the tower terminates at
arity~$4$.

\medskip
\noindent\textbf{Virasoro.}\enspace
Both $\mathfrak C_{\mathrm{Vir}}$ and $Q^{\mathrm{contact}}_{\mathrm{Vir}}$
are nonzero.  Their Poisson bracket
\[
o^{(5)}_{\mathrm{Vir}}=\{\mathfrak C_{\mathrm{Vir}},Q^{\mathrm{contact}}_{\mathrm{Vir}}\}_H=\frac{480}{c^2(5c+22)}x^5\neq 0.
\]
For all $c\neq 0,-22/5$, the Virasoro shadow tower does not terminate;
$r_{\max}(\mathrm{Vir}_c)=\infty$.  Same for every $\mathcal W_N$.

\medskip
\noindent\textbf{Heisenberg.}\enspace
$\mathfrak C_{\cH}=0$, $Q^{\mathrm{contact}}_{\cH}=0$, hence $o^{(5)}=0$
and all higher obstructions vanish.  $r_{\max}=2$.

\medskip
\noindent\textbf{Affine.}\enspace
$\mathfrak C_{\widehat\fg}\neq 0$ but $Q^{\mathrm{contact}}_{\widehat\fg}=0$
(quaternary bracket vanishes on bar cohomology by the PBW theorem).
$o^{(5)}=0$; $r_{\max}=3$.

\medskip
\noindent\textbf{$\beta\gamma$.}\enspace
$\mathfrak C_{\beta\gamma}=0$ (Lie bracket trivial on bar cohomology)
but $Q^{\mathrm{contact}}_{\beta\gamma}\neq 0$.  $o^{(5)}=0$;
$r_{\max}=4$.

\begin{center}
\renewcommand{\arraystretch}{1.15}
\begin{tabular}{lcccc}
& $\mathfrak C$ & $\mathfrak Q$ & $o^{(5)}=\{\mathfrak C,\mathfrak Q\}_H$ & $r_{\max}$\\
\hline
Heisenberg & $0$ & $0$ & $0$ & $2$\\
Affine & $\neq 0$ & $0$ & $0$ & $3$\\
$\beta\gamma$ & $0$ & $\neq 0$ & $0$ & $4$\\
Virasoro & $\neq 0$ & $\neq 0$ & $\neq 0$ & $\infty$
\end{tabular}
\end{center}

Finite termination requires $\mathfrak C=0$ or $\mathfrak Q=0$; when
both are nonzero, their bracket generates the quintic, and the
induction continues without bound.

\medskip
\noindent\textbf{The shadow growth rate.}\enspace
For class~M algebras, the Taylor coefficients $S_r$ of $\sqrt{Q_L}$
grow as $S_r\sim A\cdot\rho(\cA)^r\cdot r^{-5/2}\cos(r\theta+\varphi)$
with shadow growth rate $\rho(\cA)=\sqrt{9\alpha^2+2\Delta}/(2|\kappa|)$,
a continuous invariant.  Self-dual Virasoro ($c=13$) has
$\rho\approx 0.467$.  The critical cubic $5c^3+22c^2-180c-872=0$ has
root $c^*\approx 6.125$: the tower converges for $c>c^*$ and diverges
for $c<c^*$.

\medskip
\noindent\textbf{The shadow partition function.}\enspace
The double sum $Z^{\mathrm{sh}}_\cA=\sum_{g,r}\hbar^{2g-2}Z_g^{(r)}$
converges absolutely.  Genus decay at rate $1/(2\pi)^{2g}$
(Faber--Pandharipande); arity decay at rate $\rho^r r^{-5/2}$
(Darboux analysis).  The string partition function diverges as
$(2g-2)!$.  For Virasoro at $c=26$, $\rho\approx 0.234$ and
$Z^{\mathrm{sh}}$ has positive radius of convergence in both $\hbar$
and arity.

\subsection*{6.10.\enspace Independent sum factorisation}

For a direct sum $L=L_1\oplus L_2$ with vanishing mixed OPE, all
shadow invariants separate:
\begin{align*}
\kappa(L_1\oplus L_2)&=\kappa(L_1)+\kappa(L_2),\\
T^{\mathrm{br}}(L_1\oplus L_2)&=T^{\mathrm{br}}(L_1)\oplus T^{\mathrm{br}}(L_2),\\
\Delta_{L_1\oplus L_2}(t)&=\Delta_{L_1}(t)\cdot\Delta_{L_2}(t),\\
R_4^{\mathrm{mod}}(L_1\oplus L_2)&=R_4^{\mathrm{mod}}(L_1)+R_4^{\mathrm{mod}}(L_2).
\end{align*}
The full MC element factorises:
$\Theta_{L_1\oplus L_2}=\Theta_{L_1}\hat\otimes 1+1\hat\otimes\Theta_{L_2}$,
and $Z_{L_1\oplus L_2}=Z_{L_1}\cdot Z_{L_2}$.

Every algebra in the standard landscape is either primitive
(Heisenberg, Virasoro, $\beta\gamma$) or built from primitive pieces
(affine = current algebra over simple Lie, lattice = exponential
extension of Heisenberg, $\mathcal W_N$ = quantum Drinfeld--Sokolov
reduction of affine).  Drinfeld--Sokolov reduction is the
\emph{shadow-depth escalator}
(Principle~\ref{princ:shadow-depth-escalator}): it sends current
algebras (class L, $r_{\max}=3$, finite towers) to $\mathcal W$-algebras
(class M, $r_{\max}=\infty$, infinite towers).  The nonlinearity of
the Dirac bracket on the Slodowy slice creates the quartic and all
higher shadows.

\subsection*{6.11.\enspace The representation-theoretic perspective}

The categorical logarithm $\eta_{ij}=d\log(z_i-z_j)$ of \S1 is the
integral kernel of a tautology: $D_\cA^2=0$ on the modular bar
complex.  Three structure theorems emerge.  \emph{Algebraicity}: on
each primary slice $L$, $H(t)^2=t^4Q_L(t)$ with $Q_L$ quadratic in
$\kappa$, $\alpha$, $S_4$.  \emph{Formality}: the genus-$0$ shadow
obstruction tower is the $L_\infty$ formality obstruction tower; the
four depth classes G/L/C/M are the corresponding formality types,
while positive-genus corrections form a separate quantum layer.
\emph{Complementarity}: for a Koszul pair $(\cA,\cA^!)$, the summands
$\mathbf Q_g(\cA)\oplus\mathbf Q_g(\cA^!)$ are Lagrangian (Theorem~C).

The discriminant $\Delta=8\kappa S_4$ governs the single-line
dichotomy: $\Delta=0$ forces tower termination
($L_\infty$-formality); $\Delta\neq 0$ forces an infinite tower
(intrinsic non-formality).  The contact class C ($\beta\gamma$)
escapes via stratum separation.  No invariant of the classical bar
complex of a graded algebra sees~$\Delta$: it is visible only on
curves of positive genus, where the Arnold relation acquires a defect.

Classical Koszul duality (Priddy, Beilinson--Ginzburg--Soergel)
embeds into the genus~$0$, arity~$2$, $\Delta=0$ stratum via
formal-disk restriction (the formal, class-G case); the
$E_1/E_\infty$ distinction and Fulton--MacPherson geometry are
already absent over a bare point.
The $R$-matrix of quantum groups arises as the genus-$0$ shadow
$\operatorname{Sh}_{0,2}(\Theta_{\widehat\fg_k})$ (\S7.7).  At the
critical level $k=-h^\vee$, the scalar class vanishes ($\kappa=0$)
and bar cohomology recovers the Feigin--Frenkel center
$\operatorname{Fun}(\operatorname{Op}_{{}^L\fg})$; the Langlands
programme lives on this orthogonal axis, not in the MC element but
in the cohomology that emerges on the uncurved scalar locus.


% ====================================================================

% ====================================================================
% TRACK B: sections 7--8 (standard landscape, arithmetic of the tower)
% ====================================================================

\section*{7.\quad The standard landscape: $E_1$ chain-level census}
% ====================================================================

The bar machine of Sections 1--6 accepts any chiral algebra as input.
For each standard family, we construct the ordered bar
$\barB^{\mathrm{ord}}(\cA)$ at the chain level, compute the
differential explicitly at low tensor degrees, extract the
$r$-matrix and bar cohomology, and then average to recover
$\kappa$ and the symmetric shadow data.  The $E_1$ ordered
construction is primary; every symmetric result below is derived
from it by $\Sigma_n$-coinvariant projection.

\subsection*{7.0.\quad Three-tier $R$-matrix classification}

The genus-$0$ binary shadow $r(z) = \Res^{\coll}_{0,2}(\Theta_\cA)$
sorts chiral algebras into three tiers:
\begin{itemize}[nosep]
\item \emph{Tier (a): scalar $R$-matrix.}  The Heisenberg
algebra $\cH_k$ and lattice extensions.  The spectral $R$-matrix is
scalar, $R(z) = \exp(k/z)$, and the Yang--Baxter equation is trivial.
Shadow class $\mathbf G$.
\item \emph{Tier (b): matrix $R$-matrix from OPE poles.}  Affine
Kac--Moody $\widehat\fg_k$, Virasoro, principal $\cW_N$, $\beta\gamma$,
non-principal $\cW$-algebras.  The classical $r$-matrix is
$k\Omega/z$ (affine) or $(c/2)/z^3 + 2T/z$ (Virasoro) or the
DS-reduced analogue.  The level prefix is load-bearing:
at $k = 0$, the affine $r$-matrix vanishes.
Shadow class $\mathbf L$, $\mathbf C$, or $\mathbf M$ depending
on the depth of the OPE.
\item \emph{Tier (c): genuinely $E_1$.}  The dg-shifted Yangian
$Y^{\dg}_\hbar(\fg)$, Etingof--Kazhdan quantum vertex algebras,
line-operator atoms of 3d HT theories.  The $R$-matrix carries
genuine topological direction data: nonsymmetric braiding,
meromorphic poles of higher order, RTT relations that do not
factor through a commutative vertex algebra.  Every standard VOA
is $E_\infty$ (local); $E_1$ primacy requires passing
to the tier (c) atoms.
\end{itemize}
Tier (a) is the free-field archetype, tier (b) supplies the
standard landscape, tier (c) is the genuine $E_1$ frontier
occupied in Volumes II and III.

The seven families below exhaust the tier (a)/(b) landscape.
Each is presented in the $E_1$ order: ordered bar differential,
bar cohomology, $R$-matrix, then averaged invariants
($\kappa$, complementarity, shadow class).  No two families share
all invariants; taken together they realise every depth class
$\mathbf G / \mathbf L / \mathbf C / \mathbf M$ and every
genus-$2$ shell profile.

\medskip
\noindent\textbf{Master table of $E_1$ bar cohomology.}\enspace
\begin{center}
\small
\renewcommand{\arraystretch}{1.2}
\begin{tabular}{lllllll}
\textbf{Family} & $r(z)$ & $\dim H^n(\barB)$
  & \textbf{Growth} & \textbf{GF type}
  & $\kappa$ & \textbf{Class}\\ \hline
$\cH_k$ & $k/z$ & $\delta_{n,1}$
  & $0$ & rational & $k$ & $\mathbf G$\\
$\widehat{\mathfrak{sl}}_2$ & $k\Omega/z$ & $2n{+}1$
  & linear & rational & $\frac{3(k{+}2)}{4}$ & $\mathbf L$\\
$\widehat{\mathfrak{sl}}_3$ & $k\Omega/z$ & $\{8,36,204,\ldots\}$
  & exponential & rational & $\frac{4(k{+}3)}{3}$ & $\mathbf L$\\
$\Vir_c$ & $\frac{c/2}{z^3}{+}\frac{2T}{z}$ & $\{1,2,5,\ldots\}$
  & $3^n$ & algebraic & $c/2$ & $\mathbf M$\\
$\cW_{3,c}$ & $r_c^{W_3}(z)$ & $\{2,5,16,\ldots\}$
  & $3^n$ & algebraic & $5c/6$ & $\mathbf M$\\
$\beta\gamma$ & $r_\lambda(z)$ & $\{2,4,10,\ldots\}$
  & $3^n$ & algebraic & $6\lambda^2{-}6\lambda{+}1$ & $\mathbf C$\\
$V_\Lambda$ & $N/z$ & $\{N,\ldots\}$
  & sub-exp & transcendental & $N$ & $\mathbf G$
\end{tabular}
\end{center}

% ----- 7.1 Heisenberg -----
\subsection*{7.1.\quad Heisenberg}

The Gaussian archetype: shadow depth $2$, tower terminates at
arity $2$, every higher-genus invariant is determined by $\kappa$
alone.  One bosonic field $\alpha$ of conformal weight~$1$,
$\alpha(z)\alpha(w)\sim k/(z-w)^2$, $c=1$.

\medskip\noindent\textbf{Ordered bar differential.}\enspace
The ordered bar complex
$\barB^{\mathrm{ord}}(\cH_k)=T^c(s^{-1}\overline{\cH_k})$
has a single generator $s^{-1}\alpha$ in the augmentation ideal.

\emph{Degree~$2$}: the residue differential extracts the unique
OPE coefficient:
\[
d_{\mathrm{res}}\bigl([s^{-1}\alpha|s^{-1}\alpha]
\otimes\eta_{12}\bigr)
\;=\;k\cdot\mathbf{1}.
\]
The double-pole OPE $k/(z-w)^2$ drops by one order under $d\log$
absorption, producing a simple-pole residue.

\emph{Degree~$3$}: the collision at any boundary stratum extracts
$\alpha_{(1)}\alpha=k\cdot\mathbf{1}$, leaving $n-2$ copies
of~$\alpha$; since $\mathbf{1}_{(p)}\alpha=0$ for $p\ge 0$,
no further contraction occurs.  The Arnold relation kills the
form contribution independently.  The degree-$3$ differential
vanishes identically.  All higher degrees vanish by the same
mechanism.

\medskip\noindent\textbf{Bar cohomology.}\enspace
The complex is quadratic: all structure lives at tensor degree~$2$.
\[
H^n(\barB(\cH_k))\;=\;
\begin{cases}
\bC\cdot\alpha^* & n=1,\\
0 & n\ge 2.
\end{cases}
\]
Conformal-weight-graded bar Hilbert series: rational.
Partition-graded generating function: transcendental,
$\prod_{m\ge 1}(1-x^m)^{-1}$ (Fock-space grading).
The PBW spectral sequence collapses at~$E_1$.  $\cH_k$ is
chirally Koszul.

\medskip\noindent\textbf{$R$-matrix.}\enspace
$r^{\cH}(z)=k/z$, scalar (Tier~(a)).  Spectral $R$-matrix
$R(z)=\exp(k/z)$; Yang--Baxter trivial.

\medskip\noindent\textbf{Averaging.}\enspace
$r(z)=k/z$ is already $\Sigma_2$-invariant: $\mathrm{av}(r(z))=k
=\kappa(\cH_k)$.  Averaging loses nothing.  Koszul dual:
$\cH_k^!=(\Sym^{\mathrm{ch}}(V^*),m_0=-k\omega)$ (curved).
Complementarity (free-field branch, sum $=0$; Virasoro branch
gives $13$):
$\kappa(\cH_k)+\kappa(\cH_k^!)=k+(-k)=0$.
Shadow class $\mathbf{G}$; genus-$2$ shell: loop$\circ$loop only.

\medskip\noindent\textbf{Higher-genus.}\enspace
(UNIFORM-WEIGHT.)  The $\hat A$-genus controls every $F_g$:
\begin{equation}\label{eq:ahat-universality-preface-b}
\textstyle\sum_{g\ge 1} F_g\, x^{2g}
\;=\;
k\bigl(\hat A(ix)-1\bigr)
\;=\;
k\Bigl(\tfrac{x/2}{\sin(x/2)}-1\Bigr),
\end{equation}
so $F_1 = k/24$, $F_2 = 7k/5760$, $F_3 = 31k/967680$.  The series
converges exponentially, $|F_g| \sim 2|k|/(2\pi)^{2g}$.

% ----- 7.2 Kac-Moody -----
\subsection*{7.2.\quad Affine Kac--Moody}

The Lie/tree archetype: shadow depth $3$, cubic shadow is the
Lie bracket, Jacobi kills everything higher.  For simple~$\fg$
of rank~$\ell$ and dual Coxeter~$h^\vee$, currents $J^a$ of
weight~$1$ with
$J^a(z)J^b(w)\sim f^{ab}{}_cJ^c(w)/(z{-}w)+k\kappa^{ab}/(z{-}w)^2$.
Two pole orders: simple pole carries the bracket, double pole
carries the level.

\medskip\noindent\textbf{Ordered bar differential.}\enspace

\emph{Degree~$2$}: the differential decomposes into two components,
one per pole order:
\[
d_{\barB}\bigl([s^{-1}J^a|s^{-1}J^b]\otimes\eta_{12}\bigr)
\;=\;
\underbrace{f^{ab}{}_{c}\,[s^{-1}J^c]}_{\text{CE (simple pole)}}
\;+\;
\underbrace{k\kappa^{ab}\cdot\mathbf{1}}_{\text{curvature (double pole)}}.
\]
The Chevalley--Eilenberg component comes from the Lie bracket
$J^a_{(0)}J^b=f^{ab}{}_cJ^c$; the curvature component from the
invariant form $J^a_{(1)}J^b=k\kappa^{ab}\cdot\mathbf{1}$.

\emph{Degree~$3$ ($\mathfrak{sl}_3$)}: the bar complex
$\barB^{\mathrm{ord}}_3=\fg^{\otimes 3}
\otimes\Omega^2(\overline{\mathrm{FM}}_3)$ has dimension
$8^3\times 2=1024$.  The Serre element
\[
S\;=\;[E_1|E_1|E_2]-2[E_1|E_2|E_1]+[E_2|E_1|E_1]
\]
is a bar $3$-cocycle: $d_{\barB}(S\otimes\eta_{12}\wedge\eta_{13})=0$.
This is the Serre relation $(\mathrm{ad}\,E_1)^2(E_2)=0$
manifesting at degree~$3$.  The two independent Serre
relations (for $\{i,j\}=\{1,2\}$) generate
$H^3(\barB(\widehat{\mathfrak{sl}}_{3,k}))$ at generic level.

\emph{Degree~$4$}: for $\mathfrak{sl}_3$,
$\dim\barB^{\mathrm{ord}}_4=4096\times 6=24576$.
The bracket differential $d\colon\barB_4\to\barB_3$ is surjective
(rank $1024$); degree-$4$ cohomology is conjectured to be
$1352$-dimensional.  Full Borcherds (not Jacobi alone) is required
for $d_{\barB}^2=0$.

\medskip\noindent\textbf{Bar cohomology.}\enspace
\begin{center}
\small
\renewcommand{\arraystretch}{1.15}
\begin{tabular}{lllll}
$\fg$ & $\dim H^n$ & Growth & GF & Source\\ \hline
$\mathfrak{sl}_2$ & $2n+1$ & linear & $x(3{-}x)/(1{-}x)^{2}$
  & Garland--Lepowsky\\
$\mathfrak{sl}_3$ & $\{8,36,204,\ldots\}$ & exponential
  & $4x(2{-}13x{-}2x^2)/((1{-}8x)(1{-}3x{-}x^2))$ & Serre at deg~$3$\\
$\fg$ (general) & polynomial & $\le$ exponential & D-finite
  & PBW collapse at $E_2$
\end{tabular}
\end{center}
For $\widehat{\mathfrak{sl}}_2$, the linear growth $\dim H^n=2n+1$
is the Garland--Lepowsky concentration: bar cohomology is generated
in degree~$1$ by the $3$-dimensional adjoint.  For
$\widehat{\mathfrak{sl}}_3$, the rational generating function with
denominator $(1-8x)(1-3x-x^2)$ reflects the Serre relations; PBW
concentration at all genera is proved.

\medskip\noindent\textbf{$R$-matrix.}\enspace
$r^{\mathrm{KM}}(z)=k\,\Omega/z$ with
$\Omega=\sum_a J^a\otimes J_a$ (Tier~(b));
at $k=0$ the $r$-matrix vanishes, at $k=-h^\vee$ it is critical.
The CYBE is the Arnold relation on $\overline{\mathrm{FM}}_3(\bC)$
evaluated on the Casimir.

\medskip\noindent\textbf{Averaging.}\enspace
Averaging collapses the Casimir tensor to its trace, with
Sugawara shift:
\[
\mathrm{av}(k\,\Omega/z)=\frac{k\dim\fg}{2h^\vee}=\kappa_{\mathrm{dp}},
\qquad
\kappa(\widehat\fg_k)=\kappa_{\mathrm{dp}}+\frac{\dim\fg}{2}
=\frac{(k+h^\vee)\dim\fg}{2h^\vee}.
\]
The $\dim(\fg)/2$ is the simple-pole self-contraction through the
adjoint Casimir eigenvalue $2h^\vee$.  The Casimir tensor structure,
invisible to~$\kappa$, is the data that builds the Yangian $Y(\fg)$.

Koszul dual: $\widehat\fg_k^!=\widehat\fg_{-k-2h^\vee}$
(Feigin--Frenkel duality, $k+h^\vee\mapsto-(k+h^\vee)$).
Complementarity (KM branch, sum $=0$; Virasoro branch gives $13$):
$\kappa(\widehat\fg_k)+\kappa(\widehat\fg_k^!)=0$.
Shadow class $\mathbf L$; genus-$2$ shells: loop$\circ$loop and
sep$\circ$loop active, planted-forest absent.

\medskip\noindent\textbf{Critical level.}\enspace  At
$k=-h^\vee$: $\kappa=0$, the Sugawara construction is undefined,
and the Feigin--Frenkel centre
$\mathfrak z(\widehat\fg_{-h^\vee})\cong\Fun(\Op_{{}^L\fg}(D))$
is a commutative VA of ${}^L\fg$-opers on the formal disk.
Critical-level chiral Koszul duality is bar-level geometric
Langlands.

% ----- 7.3 Virasoro -----
\subsection*{7.3.\quad Virasoro}

The mixed archetype: shadow depth infinity, tower never terminates.
One generator $T$ of weight~$2$,
$T(z)T(w)\sim(c/2)/(z{-}w)^4+2T/(z{-}w)^2+\partial T/(z{-}w)$.
Three pole orders.  The self-coupling $T_{(1)}T=2T$ drives all
nonlinear complexity.

\medskip\noindent\textbf{Ordered bar differential.}\enspace

\emph{Degree~$2$}: the bar differential mixes both pole channels:
\[
d_{\barB}\bigl([s^{-1}T|s^{-1}T]\otimes\eta_{12}\bigr)
\;=\;
\underbrace{\tfrac{c}{2}\cdot\mathbf{1}}_{\substack{\text{curvature}\\
\text{(order-$4$ pole)}}}
\;+\;
\underbrace{2\,[s^{-1}T]}_{\substack{\text{self-coupling}\\
\text{(order-$2$ pole)}}}
\;+\;
\underbrace{[\partial T]}_{\substack{\text{translation}\\
\text{(order-$1$ pole)}}}.
\]
Three contributions from three OPE pole orders, each reduced
by one under $d\log$ absorption: order-$4\to 3$, $2\to 1$, $1\to 0$.
The curvature element $m_0=c/2\in\barB^0=\bC$ is the bar-dual
of the quartic OPE pole.

\emph{Degree~$3$}: the ordered bar
$\barB_3=(s^{-1}T)^{\otimes 3}\otimes\Omega^2(\overline{\mathrm{FM}}_3)$
has differential
\begin{align*}
d_{\barB}\bigl([s^{-1}T|s^{-1}T|s^{-1}T]
&\otimes\eta_{12}\wedge\eta_{13}\bigr)\\
&=\;\tfrac{c}{2}\bigl([s^{-1}T]\otimes(\eta_{12}+\eta_{13}+\eta_{23})\bigr)
\;+\;\text{self-coupling terms}.
\end{align*}
The vacuum term $\frac{c}{2}[s^{-1}T]\otimes(\eta_{12}+\eta_{13}+\eta_{23})$
vanishes by the Arnold relation: three boundary divisors, three
forms, zero sum.  The surviving self-coupling terms
$[\partial T|T]\otimes\eta_{13}+[T|\partial T]\otimes(\eta_{12}+\eta_{23})$
are nontrivial: this is where infinite depth begins.

\medskip\noindent\textbf{Bar cohomology.}\enspace
The bar cohomology of $\Vir_c$ exhibits exponential growth $3^n$,
in contrast to the linear growth of Kac--Moody.
The \emph{Motzkin path model}: $\dim H^n(\barB(\Vir_c))=M(n{+}1)-M(n)$,
where $M(n)$ is the $n$-th Motzkin number
(paths in $\bZ_{\ge 0}$ of length~$n$ with steps
$+1$, $0$, $-1$).  The generating function is algebraic:
\[
\sum_{n\ge 1}\dim H^n\cdot x^n
\;=\;x+2x^2+5x^3+12x^4+30x^5+\cdots\,.
\]
The conformal-weight-graded spectral sequence collapses at~$E_2$;
PBW concentration holds.

\medskip\noindent\textbf{$R$-matrix.}\enspace
$r_c(z)=(c/2)/z^3+2T/z$: the quartic and quadratic OPE poles
drop by one order each under $d\log$ absorption.  Tier~(b),
class~$\mathbf{M}$.

\medskip\noindent\textbf{Averaging.}\enspace
$\mathrm{av}(r_c(z))=c/2=\kappa(\Vir_c)$.  The stress-tensor
coupling $2T/z$, invisible to~$\kappa$, drives the infinite
shadow tower.  Koszul dual: $\Vir_c^!=\Vir_{26-c}$ (self-dual at $c=13$).  Complementarity (Virasoro branch, sum $=13$; KM
branch gives $0$):
\[
\kappa(\Vir_c)+\kappa(\Vir_c^!)\;=\;\tfrac{c}{2}+\tfrac{26-c}{2}
\;=\;13.
\]
Shadow class $\mathbf M$; genus-$2$: all three shells active
including planted-forest (first family to see log-FM boundary
strata of $\overline\cM_{2,n}$).

\medskip\noindent\textbf{Quartic shadow and infinite tower.}\enspace
The quartic contact shadow
\begin{equation}\label{eq:vir-quartic-b}
Q^{\mathrm{contact}}_{\Vir} \;=\; \frac{10}{c(5c+22)},
\end{equation}
has poles at $c=0$ and $c=-22/5$ (the $(2,5)$ minimal model).
The quintic obstruction
$o^{(5)}=\{C_{\Vir},Q^{\mathrm{contact}}\}_H
=480/[c^2(5c+22)]\,x^5$
is nonzero: no finite truncation captures Virasoro.

\medskip\noindent\textbf{Critical dimension $c=26$.}\enspace
The shadow connection
$\nabla^{\hol}_c=d-(c/2)\omega_g
-[10/(c(5c+22))]\delta_4-\cdots$
is flat for every~$c$ by the MC equation.  The critical dimension
arises from the \emph{dual} curvature vanishing:
$\kappa(\Vir_{26-c})=(26-c)/2=0$ forces $c=26$.  At $c=26$ the
Koszul dual $\Vir_0$ is uncurved; the algebra itself retains
$\kappa=13$.

\medskip\noindent\textbf{Polyakov identification.}\enspace
The Polyakov formula
$\log(\det'\Delta_{g_1}/\det'\Delta_{g_0})
=-(1/6\pi)\int(|\nabla\sigma|^2+R\sigma)$
is the arity-$2$ shadow evaluation at $\kappa=1$; the
higher-arity corrections $\delta_4$, $\delta_H$ are
inaccessible to the free-field path integral.

% ----- 7.4 W_3 -----
\subsection*{7.4.\quad $\boldsymbol{\cW_3}$}

First multi-generator family: two generators $T$ (weight~$2$)
and $W$ (weight~$3$).  The multi-weight structure produces
cross-channel mixing that the Virasoro single-generator story
cannot see.

\medskip\noindent\textbf{Ordered bar differential.}\enspace
The $TT$ OPE contributes the same three-pole structure as
Virasoro.  The $WW$ OPE introduces a sixth-order pole:
\begin{align*}
W(z)W(w)
&\sim\frac{c/3}{(z{-}w)^6}+\frac{2T}{(z{-}w)^4}
+\frac{\partial T}{(z{-}w)^3}\\
&\quad+\frac{1}{(z{-}w)^2}\Bigl[\tfrac{16}{22+5c}\Lambda
+\tfrac{3}{10}\partial^2 T\Bigr]+\cdots,
\end{align*}
where $\Lambda=\!:TT:\!-(3/10)\partial^2 T$ is the weight-$4$
quasi-primary composite field.  At degree~$2$ the ordered
bar differential on $[s^{-1}W|s^{-1}W]$ produces five
terms from five pole orders (order $6\to 5$, $4\to 3$,
$3\to 2$, $2\to 1$, $1\to 0$ under $d\log$ absorption),
including the composite field $\Lambda$ at the second-order
pole.  The coefficient $\alpha(c)=16/(22+5c)$ is fixed
uniquely by quasi-primarity ($L_1\Lambda=0$) and the Jacobi
identity.  At $c=-22/5$, $\Lambda$ becomes null and $\cW_3$
is not finitely generated by $T,W$ alone.

\emph{Cross-channel at degree~$3$}: the degree-$3$ bar
differential mixes $TT$, $TW$, and $WW$ channels.  The
$TW$ cross-channel produces terms involving $\Lambda$ that
have no analogue in the single-generator Virasoro;
these drive the multi-weight genus-$2$ correction.

\medskip\noindent\textbf{Bar cohomology.}\enspace
Bar cohomology growth rate $3^n$ (same as Virasoro: the
self-coupling $T_{(1)}T=2T$ generates the same exponential
mechanism).  The generating function is algebraic.  The
two-generator PBW spectral sequence collapses at~$E_2$
with a $2\times 2$ spectral discriminant:
$\Delta_{\cW_3}(x)=(1-\tfrac{c}{2}x)(1-\tfrac{c}{3}x)$,
branch points $2/c$ and $3/c$.

\medskip\noindent\textbf{$R$-matrix.}\enspace
The classical $r$-matrix is the Casimir piece plus $T$- and
$W$-mediated corrections; the per-channel eigenvalues are
$\kappa_T=c/2$ (from $T$) and $\kappa_W=c/3$ (from $W$).
Tier~(b), class~$\mathbf{M}$.

\medskip\noindent\textbf{Averaging.}\enspace
$\mathrm{av}(r_c^{W_3}(z))=c/2+c/3=5c/6=\kappa(\cW_3)$:
the weight-$2$ generator contributes $c/2$, the weight-$3$
generator contributes $c/3$, additively.
Koszul dual: $\cW_3^!=\cW_{3,100-c}$ (self-dual at $c=50$).
\[
\kappa(\cW_{3,c})+\kappa(\cW_{3,100-c})\;=\;\tfrac{250}{3}.
\]
Shadow class $\mathbf M$; genus-$2$: all three shells active,
planted-forest sensitive to $WW\Lambda$ channel.

\medskip\noindent\textbf{Genus-$2$ cross-channel.}\enspace
(ALL-WEIGHT, with cross-channel correction.)
\begin{equation}\label{eq:w3-genus2-b}
F_2(\cW_3)
\;=\;
\underbrace{\tfrac{7c}{6912}}_{\kappa\cdot\lambda_2^{\FP}}
\;+\;
\underbrace{\tfrac{c+204}{16c}}_{\delta F_2^{\mathrm{cross}}},
\end{equation}
with $\delta F_2^{\mathrm{cross}}$ computed graph-by-graph over
the seven genus-$2$ stable graphs.  The scalar UNIFORM-WEIGHT
formula fails at $g\ge 2$ for multi-weight algebras.
A universal closed form for principal $\cW_N$:
\[
\delta F_2^{\mathrm{grav}}(\cW_N,c)
\;=\;
\tfrac{(N-2)(N+3)}{96}
+\tfrac{(N-2)(3N^3+14N^2+22N+33)}{24c}.
\]

\medskip\noindent\textbf{Holographic datum.}
$\cH(\cW_{3,c})=(\cW_{3,c},\cW_{3,100-c},\cC_c^{W_3},
r_c^{W_3}(z),\Theta_c,\nabla^{\hol}_c)$.

% ----- 7.5 beta-gamma -----
\subsection*{7.5.\quad $\boldsymbol{\beta\gamma}$ system}

The contact archetype: shadow depth~$4$, first nonlinearity at
arity~$4$ (not~$3$).  Fields $\beta$ of weight~$\lambda$,
$\gamma$ of weight $1-\lambda$, with
$\beta(z)\gamma(w)\sim 1/(z-w)$ and both self-OPEs vanishing.
Simple pole only.

\medskip\noindent\textbf{Ordered bar differential.}\enspace
The ordered bar has two generators $s^{-1}\beta$, $s^{-1}\gamma$
in the augmentation ideal.  At degree~$2$:
\[
d_{\barB}\bigl([s^{-1}\beta|s^{-1}\gamma]\otimes\eta_{12}\bigr)
\;=\;\mathbf{1},
\qquad
d_{\barB}\bigl([s^{-1}\beta|s^{-1}\beta]\otimes\eta_{12}\bigr)
\;=\;0,
\]
since $\beta_{(0)}\gamma=\mathbf{1}$ (simple pole) and
$\beta_{(0)}\beta=0$ (self-OPE vanishes).  The simple-pole-only
structure means $d\log$ absorption produces a degree-$0$ residue:
the OPE pole is order~$1$, absorbed to order~$0$.

At degree~$3$: the cubic shadow $\ell_3^{(0)}$ vanishes because
$\beta$ and $\gamma$ carry different conformal weights; the
$\mathrm{U}(1)$ ghost-number grading forces all odd-arity
ordered shadows to zero.  At degree~$4$: the quartic
$\ell_4^{(0)}\neq 0$ from $\langle\beta\gamma\beta\gamma\rangle$
cross-contractions, but the quartic contact invariant
$\mu_{\beta\gamma}=\langle\eta,m_3(\eta,\eta,\eta)\rangle=0$
by ghost-number rigidity.  The tower terminates: depth~$4$, not
by Lie cancellation (Jacobi) but by an abelian grading mechanism.

\medskip\noindent\textbf{Bar cohomology.}\enspace
Growth rate $3^n$ (same exponential class as Virasoro).
Generating function: algebraic.  The two-generator PBW spectral
sequence collapses at~$E_2$.  Central charge:
$c_{\beta\gamma}(\lambda)=2(6\lambda^2-6\lambda+1)$; at
$\lambda=1$, $c=2$.  Partner: the fermionic $bc$ system has
$c_{bc}(\lambda)=1-3(2\lambda-1)^2$, and $c_{\beta\gamma}+c_{bc}=0$
(verify: $\lambda=1$ gives $2+(-2)=0$).

\medskip\noindent\textbf{$R$-matrix.}\enspace
Proportional to the identity Casimir on the $\beta\gamma$-state
space; the spectral $R$-matrix factors through the ghost-number
grading.  Tier~(b), class~$\mathbf{C}$.

\medskip\noindent\textbf{Averaging.}\enspace
$\kappa(\beta\gamma)=c_{\beta\gamma}/2=6\lambda^2-6\lambda+1$.
Koszul dual: $(\beta\gamma)^!=bc$ (bosonic/fermionic exchange).
Complementarity (free-field branch, sum $=0$; Virasoro branch
gives $13$): $\kappa(\beta\gamma)+\kappa(bc)=0$.
Shadow class $\mathbf{C}$; genus-$2$: loop$\circ$loop only
(contact termination, distinct from Gaussian collapse).

% ----- 7.6 Lattice -----
\subsection*{7.6.\quad Lattice vertex algebras}

Arithmetic laboratory: the shadow tower inherits the Gaussian
archetype from the Heisenberg factor; the braiding carries the
full arithmetic of the lattice.  Curvature sees only rank;
braiding sees the quadratic form.

\medskip\noindent\textbf{Ordered bar differential.}\enspace
For an even lattice $\Lambda$ of rank~$N$, quadratic form~$q$,
$2$-cocycle~$\varepsilon$:
$V_\Lambda^{N,q}=\cH_1^{\otimes N}\otimes_\bC\bC[\Lambda]_\varepsilon$
(level $k=1$ from the lattice bilinear form).
The ordered bar decomposes as a tensor product of $N$ copies of
the Heisenberg ordered bar, tensored with the group algebra of
the lattice.  At degree~$2$, the differential acts on each
Heisenberg factor independently:
\[
d_{\barB}\bigl([s^{-1}\alpha^i|s^{-1}\alpha^j]
\otimes\eta_{12}\bigr)
\;=\;\delta^{ij}\cdot\mathbf{1},
\]
producing the identity Casimir in the rank-$N$ current space.
The vertex operators $e^{\lambda\cdot\phi}$ contribute only
through the lattice cocycle $\varepsilon(\lambda,\mu)$, which
enters the ordering of tensor factors but not the collision
residue.  All structure is captured at tensor degree~$2$:
the bar complex is quadratic.

\medskip\noindent\textbf{Bar cohomology.}\enspace
$H^1(\barB(V_\Lambda))$ has dimension~$N$ (the $N$ current
generators $\alpha^{i*}$); higher cohomology inherits the
Heisenberg vanishing $H^n=0$ for $n\ge 2$ in each tensor
factor, but the lattice group algebra contributes additional
classes through the vertex operator sector.  The
partition-graded generating function is transcendental:
$\prod_{m\ge 1}(1-x^m)^{-N}$, the $N$-fold Fock-space grading.

\medskip\noindent\textbf{$R$-matrix.}\enspace
$r^{\Lambda}(z)=N/z$ (scalar, rank copies of Heisenberg
$r$-matrix at level $k=1$).  Tier~(a); spectral $R$-matrix
$R(z)=\exp(N/z)$; Yang--Baxter trivial.  The braiding
$\pi_1(\mathrm{Conf}_n(X))\to\mathrm{Aut}(V_\Lambda^{\otimes n})$
from the lattice $2$-cocycle is invisible to the $r$-matrix.

\medskip\noindent\textbf{Averaging.}\enspace
$\kappa(V_\Lambda^{N,q})=\mathrm{rank}(\Lambda)=N$,
independent of $\varepsilon$ and~$q$.  Curvature-braiding
orthogonality: the curvature class
$[\kappa\cdot\omega_g]\in H^{1,1}(\overline\cM_g)$ and the
braiding representation act on complementary factors of the
state space.  Shadow class $\mathbf{G}$; depth~$2$;
genus-$2$: loop$\circ$loop only.
Self-dual lattices ($E_8$: $N=8$; Leech: $N=24$) give
holomorphic VAs with trivial modular functor and theta-function
partition $\Theta_\Lambda/\eta^N$.

% ----- 7.7 Yangians -----
\subsection*{7.7.\quad Yangians and quantum groups}

The line-operator algebra.  Yangian $Y(\fg)$ governs the
topological direction of a 3d HT theory; its $R$-matrix is the
genus-$0$ binary shadow of $\Theta$.  Yang--Baxter is the
arity-$3$ MC equation; RTT encodes all arities.

\medskip\noindent\textbf{$R$-matrix as genus-$0$ binary shadow.}
\[
r(z) \;=\; \Res^{\coll}_{0,2}(\Theta_\cA),
\qquad
R(z) \;=\; 1 + r/z + r^{(2)}/z^2 + \cdots.
\]
The classical Yang--Baxter equation
$[r_{12}, r_{13}] + [r_{12}, r_{23}] + [r_{13}, r_{23}] = 0$
is the Arnold relation on $\FM_3(\bC)$ evaluated on the
Casimir, and the three terms correspond to the three boundary
divisors of $\overline C_3(\bC)$.

\medskip\noindent\textbf{MC3 and all-types CG.}  On the
type-$A$ Yangian surface the post-CG problem reduces to a
single compact/completed extension packet by (a) chromatic
filtration of $K_0(\Rep)$ by $q$-character support,
(b) prefundamental Clebsch--Gordan closure
$V_n \otimes L^- = \bigoplus_{j=0}^n L^-(n-2j)$ at the
character level, and (c) Francis--Gaitsgory pro-nilpotent
completion.  The CG decomposition is extended to all simple
types via the Chari--Moura multiplicity-free $\ell$-weight
property (Theorem~\ref{thm:categorical-cg-all-types},
Corollary~\ref{cor:mc3-all-types}); the remaining post-CG
completion packets are open beyond type~$A$.

\medskip\noindent\textbf{dg-shifted Yangian.}  The modular
dg-shifted Yangian
$Y_T^{\mod} = \varprojlim_r Y_T^{\mod,\le r}$ is the
pro-nilpotent completion of the convolution dg Lie algebra of
the 3d HT chiral algebra for $T$, and the modular $R$-matrix
$R_T^{\mod}(z;\hbar) \in \MC(Y_T^{\mod})$ is an MC element.
The Yang--Baxter equation is the arity-$3$ projection; the
higher RTT relations encode all arities; the inverse limit
converges by Mittag-Leffler stabilisation.

\medskip\noindent\textbf{Drinfeld--Kohno.}  The KZ connection
$\nabla_{\KZ} = d - \hbar\sum_{i<j} \Omega_{ij}\,dz_{ij}/z_{ij}$
on $\Conf_n(\bC)$ has monodromy that agrees with the
quantum-group representation category by the affine
Drinfeld--Kohno theorem; shadow-side this is the statement
$\nabla_{\KZ} = \mathrm{Sh}_{0,n}(\Theta_{\widehat\fg_k})$ on
the affine comparison surface.

\subsection*{7.8.\quad Drinfeld--Sokolov reduction as package morphism}

$\mathrm{DS}\colon \widehat{\mathfrak{sl}}_3{}_k \to \cW_{3,c(k)}$
with $c(k) = 2 - 24(k+2)^2/(k+3)$ extends to a morphism of
holographic data:
\begin{center}
\renewcommand{\arraystretch}{1.15}
\begin{tabular}{lll}
& $\cH(\widehat{\mathfrak{sl}}_3{}_k)$
  & $\mathrm{DS}(\cH)$\\
\hline
$\cA$ & $\widehat{\mathfrak{sl}}_3{}_k$ & $\cW_{3,c(k)}$\\
$\cA^!$ & $\widehat{\mathfrak{sl}}_3{}_{-k-6}$ & $\cW_{3,100-c(k)}$\\
$r(z)$ & $k\,\Omega_{\mathfrak{sl}_3}/z$ & $r_c^{W_3}(z)$\\
$\kappa$ & $4(k+3)/3$ & $5c(k)/6$\\
depth & $3$ & $\infty$\\
$\Theta^{(2)}_{\mathrm{pf}}$ & $0$ & $\neq 0$
\end{tabular}
\end{center}
The Casimir $r$-matrix $k\Omega/z$ specialises to
$r_c^{W_3}(z)$; the curvature shifts; the shadow depth jumps
$3 \to \infty$; the planted-forest shell activates.  At the
critical level $k = -3$: $\kappa \to 0$ on the affine side,
$c \to \infty$ on the $\cW_3$ side (Sugawara pole).

The DS-HPL transfer theorem (Volume II) makes the chain-level
mechanism precise.  Homological perturbation through the BRST
SDR transfers the affine dg-shifted Yangian triple
$(m, r, \Delta)$ to the $\cW$-algebra side: the $A_\infty$
products transfer with nontrivial quartic pole (infinite
arity); the $r$-matrix transfers with the expected pole-order
shift; the coproduct is annihilated by a ghost-number
obstruction, so $\Delta_z^\cW = 0$ at all arities.  Gravity
does not produce particles, only braids them.  Central-charge
additivity: $c_{\mathrm{ghost}} = 6 = 2|\Delta^+(\fsl_3)|$,
$k$-independent.

\subsection*{7.9.\quad Geometric localization principle}

Drinfeld--Sokolov reduction is Hamiltonian reduction from
$J_\infty(\fg^*)$ to $J_\infty(S_f)$, where $S_f = f + \fg^e$
is the Slodowy transverse slice to the nilpotent orbit of $f$.
The Slodowy slice meets every orbit transversally, carries a
Kirillov--Kostant Poisson structure, and its quantisation is
the finite $W$-algebra (Gan--Ginzburg, Premet).  The chiral
upgrade replaces $S_f$ by its arc space, and
$\gr_F \cW^k(\fg, f) \cong \bC[J_\infty(S_f)]$.

In this reading: the PBW--Slodowy collapse is smoothness of
the Slodowy slice; the shadow depth jump $3 \to \infty$
under DS records the nonlinearity of the Dirac bracket on
$J_\infty(S_f)$; Barbasch--Vogan orbit duality
$f \leftrightarrow f^D$ is Springer duality of Slodowy slices,
and curvature complementarity is an algebraic identity under
the Feigin--Frenkel involution; the master commutative square
(Theorem~\ref{thm:master-commutative-square}) organises the
whole picture, with bar complexes on the upper row and
Hamiltonian reduction on arc spaces on the lower row.


% ====================================================================
\section*{8.\quad Arithmetic of the shadow obstruction tower}
% ====================================================================

The shadow obstruction coefficients for the standard families are
rational functions of the central charge.  Their denominators,
zeros, and growth rates carry number-theoretic information
connecting the shadow calculus to $L$-functions.  What arithmetic
does the MC equation force?

\subsection*{The Dirichlet--sewing lift}

The genus-$1$ connected free energy $F_1^{\conn}(\cA;q) =
-\log\det(1 - K_q)$ encodes connected sewing amplitudes
$a_\cA(N)$ as Fourier coefficients.  The Dirichlet--sewing lift
\begin{equation}\label{eq:preface-dirichlet-sewing-b}
S_\cA(u)
\;=\;
\sum_{N \ge 1} a_\cA(N) N^{-u}
\;=\;
\frac{(2\pi)^u}{\Gamma(u)}
\int_0^\infty F_\cA^{\conn}(e^{-2\pi y})\, y^{u-1}\, dy
\end{equation}
is a Dirichlet series whose analytic structure is controlled by
$\Theta_\cA$.  The sewing amplitudes define an inner product
$\langle f, g\rangle_\cA = \sum_N a_\cA(N) f(N) g(N)$ with
reproducing kernel $K_\cA(s,t) = S_\cA(s+t)$; the surface
moment matrix
$\mathscr M_{ij} = K_\cA(\alpha/2 + i, \alpha/2 + j)$ is
positive definite for all $\alpha \ge 2$
(Theorem~\ref{thm:gram-positivity}), a Gram matrix $V^T D V$
with Vandermonde $V$.

Explicit values:
\[
\begin{aligned}
S_\cH(u) &= \zeta(u)\zeta(u+1),\\
S_{V_k(\fg)}(u) &= \dim(\fg)\cdot\zeta(u)\zeta(u+1),\\
S_{\Vir}(u) &= \zeta(u+1)(\zeta(u) - 1),\\
S_{\cW_N}(u) &= \zeta(u+1)\Bigl((N-1)\zeta(u)
- \textstyle\sum_{j=1}^{N-1} H_j(u)\Bigr),
\end{aligned}
\]
where $H_j(u) = \sum_{n=1}^j n^{-u}$ are partial harmonic sums.

\subsection*{The three-tier Euler--Koszul classification}

The \emph{Euler defect} $D_\cA(u) := S_\cA(u)/(|W(\cA)|\,\zeta(u)\,\zeta(u+1))$,
where $|W(\cA)|$ is the number of strong generators
($|W(\cH)|=1$, $|W(\widehat\fg_k)|=\dim\fg$, $|W(\beta\gamma)|=2$),
sorts chiral algebras:
\begin{center}
\small
\renewcommand{\arraystretch}{1.15}
\begin{tabular}{lll}
\textbf{Tier} & \textbf{Condition} & \textbf{Families}\\ \hline
Exact Euler--Koszul & $D_\cA(u) \equiv 1$
  & $\cH$, $V_k(\fg)$, $\beta\gamma$\\
Finitely defective & $1 - D_\cA(u)$ finite harmonic ratio
  & Virasoro, $\cW_N$\\
Non-Euler--Koszul & $S_\cA$ decomposes into Hecke $L$-functions
  & lattice VOAs with $h(D) \ge 2$
\end{tabular}
\end{center}
The Euler--Koszul class $\mathrm{ek}(\cA) = \max_i(w_i - 1)$
and the shadow depth $r_{\max}(\cA)$ are independent invariants
(Theorem~\ref{thm:shadow-euler-independence}): character
arithmetic and OPE interaction are logically independent axes.

\subsection*{The two-variable $L$-object}

Rankin--Selberg unfolding against $E^*(\tau, s)$ gives
\begin{equation}\label{eq:preface-two-variable-L-b}
L_\cA(s, u)
\;=\;
\int_\cF^{\mathrm{reg}} F_\cA^{\conn}(\tau) E^*(\tau, s) y^{u-2}\,dx\,dy.
\end{equation}
The \emph{sewing--Hecke reciprocity theorem} collapses both
variables to one: $L_\cA(s, u) = \Gamma(v)/(2\pi)^v\,S_\cA(v)$
with $v = s + u - 1$.  The scattering poles of $E^*$ at
$s = \rho/2$ (nontrivial zeros $\rho$ of $\zeta$) are
\emph{not} poles of $L_\cA$; the nontrivial zeros appear as
structural obstructions to the Euler product, not as poles of
the $L$-object.

\subsection*{Shadow-moduli resolution}

The shadow-moduli map $t\colon \mathfrak H \to \bC$ solves
$\prod_j(1 - \lambda_j\,t(q))^{c_j} = \det(1 - K_q^{\conn})$
with spectral atoms $\lambda_j$ and multiplicities $c_j$.  The
connected free energy factorises as $F_1^{\conn}(q) =
G_\cA(t(q))$ with
\[
G_\cA(t) \;=\; \int\log(1 - \lambda t)\,d\rho_\cA(\lambda),
\]
determined by the spectral measure alone
(Theorem~\ref{thm:universality-of-G}).  The moments
$\mu_r = \int\lambda^r\,d\rho_\cA$ satisfy the MC recursion
\[
\mu_{r+1} \;=\; (r+1)\,\mathrm{Br}_{r+1}(\mu_2,\ldots,\mu_r;\cA)
\]
(the bar-recursion polynomial determined by the MC structure)
arity by arity.  For two Satake atoms $\alpha, \beta$, the
Newton relations give the power sums
$p_r = e_1 p_{r-1} - e_2 p_{r-2}$, encoding the MC recursion
in spectral coordinates.

\subsection*{Interacting Gram matrix and resonance}

The shadow obstruction tower provides corrections to Gram
positivity.  The interacting weight
$w_\cA(N;\alpha) = a_\cA(N)\cdot\prod_j(1 - \lambda_j N^{-\alpha/2})^{c_j}$
is unconditionally positive when all atoms $\lambda_j$ are
negative.  For Virasoro with $c > 0$:
$\lambda_{\mathrm{eff}} = -6/c < 0$, so positivity holds
unconditionally for $\alpha \ge 2$.  The resonance locus $\cR$
is empty for $c > 0$ and lies in $c < 0$, with the Lee--Yang
point $c = -22/5$ as the first entry.

\subsection*{The Weil analogy}

\begin{center}
\small
\renewcommand{\arraystretch}{1.1}
\begin{tabular}{ll}
\textbf{Weil proof for $C/\bF_q$} & \textbf{MC equation for $\cA$}\\ \hline
Curve $C/\bF_q$ & chirally Koszul $\cA$\\
Frobenius $\mathrm{Frob}_q$ & MC element $\Theta_\cA = D_\cA - d_0$\\
$\mathrm{Frob}^2 = q$ & $D_\cA^2 = 0$\\
\'Etale cohomology $H^i_{\text{\'et}}$ & cyclic deformation $\Defcyc(\cA)$\\
Weight filtration & arity filtration $\Theta_\cA^{\le r}$\\
Eigenvalues $\alpha_i$ on $H^1$ & spectral atoms $\lambda_j$\\
$\|\alpha_i\| = q^{1/2}$ & $|a_{f_j}(p)| \le 2p^{(k-1)/2}$\\
Lefschetz $L$ & bracket with $\kappa(\cA)$
\end{tabular}
\end{center}
For even unimodular lattice VOAs $V_\Lambda$ with
$\Theta_\Lambda = c_E E_{r/2} + \sum_j c_j f_j$, the
Hecke--Newton closure gives: MC equation $\Rightarrow$
prime-local CG closure $\Rightarrow$ CPS automorphy
$\Rightarrow \Sym^{r-1}(f)$ $\Rightarrow$ Ramanujan bound.
Deligne's theorem supplies the bound itself; the MC framework
supplies the systematic Hecke decomposition.

\subsection*{Polyakov--bar-cobar dictionary}

\begin{center}
\small
\renewcommand{\arraystretch}{1.1}
\begin{tabular}{ll}
\textbf{Polyakov} & \textbf{Bar-cobar}\\ \hline
Conformal anomaly $(c/12)R$ & modular characteristic $\kappa$\\
$\det'_\zeta \Delta$ & Fredholm $\det(1 - K)$\\
Critical dimension $c = 26$ & dual curvature $\kappa(\Vir_{26-c}) = 0$\\
Ghost system $bc$ & Koszul dual $\cA^!$ (BRST pairing)\\
Genus expansion $Z = \sum g_s^{2g-2}Z_g$
  & MC5 sewing, $F_g = \kappa\lambda_g^{\FP}$ (uniform)\\
BPZ equations & shadow connection on comparison surface\\
Instanton sectors & Novikov-completed bar $B^\Lambda(\cA)$\\
Wilson line $\langle W(C)\rangle$ & Swiss-cheese boundary module
\end{tabular}
\end{center}
Anomaly, genus expansion, duality, critical dimension, and
non-perturbative completions all descend from
$\Theta_\cA \in \MC(\gAmod)$.  The four depth classes
$\mathbf G/\mathbf L/\mathbf C/\mathbf M$ correspond to
increasing departures from the Gaussian path integral.

\subsection*{Arithmetic packet connection}

Hecke operators $T_p$ act on the space $M_\cA$ of activated
graph amplitudes $\ell_\Gamma^{(g)}$.  Decompose
$M_\cA = \bigoplus_\chi M_\chi$ into Hecke generalised
eigenspaces with eigenvalue $a_\chi(p)$ and nilpotent part
$N_\chi$; attach to each eigenspace its packet function
$\Lambda_\chi(s)$ (completed $L$-function of the eigenform).
The \emph{arithmetic packet connection}
\[
\nabla_\cA^{\mathrm{arith}}
\;=\;
d \;-\; \bigoplus_\chi
\tfrac{\Lambda_\chi'(s)}{\Lambda_\chi(s)}
(\mathrm{id}_{M_\chi} + N_\chi)\,ds
\]
is flat (each block is scalar form times constant endomorphism).
The singular divisor $D_\cA$ is the zero-pole set of the packet
functions, depending only on the arithmetic skeleton
$M_\cA^{\mathrm{ss}}$.

The Eisenstein scattering function on $\cM_{1,1}$ is
$\varphi(s) = \Lambda^*(1-s)/\Lambda^*(s)$ with
$\Lambda^*(s) = \pi^{-s}\Gamma(s)\zeta(2s)$.  The frontier
defect form
$\Omega_\cA = d\log\Lambda_{\mathrm{Eis}}(s) - d\log\varphi(s)$
vanishes iff the Eisenstein block is gauge-equivalent to the
scattering connection.  For lattice VOAs at rank $r$,
$\Omega_{V_\Lambda} \neq 0$ in general because $\zeta(s)$ and
$\zeta(2s)$ have different zero sets.

\subsection*{Finite Miura defect}

The principal $\cW_N$ character is
$\chi_{\cW_N}(q) = \prod_{n \ge 2}
(1 - q^n)^{-\min(n-1, N-1)}$;
the Heisenberg character is
$\chi_\cH(q) = \prod_{n \ge 1}(1 - q^n)^{-1}$.  The ratio is
finite:
\[
\chi_{\cW_N}(q)
\;=\;
\chi_\cH(q)^{N-1}
\cdot \prod_{m=1}^{N-1}(1 - q^m)^{N-m}.
\]
The finite product is the Miura defect; all genus-$1$
arithmetic content in $\cW_N$ beyond Heisenberg is carried by
it.  The prime-side Dirichlet analogue is
$D_{\cW_N}^{\mathrm{prime}}(u) = -\zeta(u+1)
\sum_{m=1}^{N-1}(N - m)m^{-u}$; the packet connection splits
into $(N-1)$ Heisenberg copies plus a defect governing modes
$m = 1, \ldots, N - 1$.

\subsection*{Quartic closure programme}

The quartic resonance class
$\mathfrak{R}_{4,g,n}^{\mod}(\cA)$ satisfies a clutching law
along every separating stratum
\[
\xi^*\mathfrak{R}_{4,g,n}^{\mod}
\;=\;
p_1^*\mathfrak{R}_{4,g_1,n_1+1}^{\mod}
+ p_2^*\mathfrak{R}_{4,g_2,n_2+1}^{\mod}
+ \mathfrak{T}_{3,\xi}
\]
(Theorem~\ref{thm:nms-clutching-law-modular-resonance}).  The
$3 \times 3$ Schur complement of the mixed $(s,u)$-transform
extracts the quartic contact potential $S_4(\cA)$ on the potential side and its
reciprocal on the Gram side.  For Virasoro:
$S_4(\Vir_c) = 10/[c(5c+22)]$.

Three closure conjectures form a hierarchy (each implies its
predecessors): \emph{quartic closure} (intersection of Gram
loci equals $\{\Re\rho = 1/2\}$,
Conjecture~\ref{conj:quartic-closure}),
\emph{Beilinson closure} (global discrepancy
$\cB^{\mathrm{glob}}(\rho)$ vanishes iff $\Re\rho = 1/2$,
Conjecture~\ref{conj:beilinson-closure}), and
\emph{clutching closure} (residue clutching defect
$K_{\cA,\rho,\xi}$ vanishes for all families iff
$\Re\rho = 1/2$, Conjecture~\ref{conj:clutching-closure}).
The chain of constructions
$\Theta_\cA \to w_{c,\rho,u_0} \to$ Gram matrix $\to$ Schur
complement $\to$ quartic resonance divisor $\to$ residue
clutching defect $\to$ closure conjecture is proved up to the
final link, which is conjectural.

\subsection*{Constrained Epstein zeta}

For a lattice VOA $V_\Lambda$ of rank $r$, the constrained
Epstein zeta
$\varepsilon^c_s(V_\Lambda) = \sideset{}{'}\sum_{\ell \in
\Lambda^{\bar}}\|\ell\|^{-2s}$
restricts the Epstein series to the sublattice
$\Lambda^{\bar} \subset \Lambda$ determined by the bar complex.
For the rank-$1$ lattice at self-dual radius $R = 1$, every
nonzero integer contributes:
$\varepsilon^1_s(R{=}1) = 4\zeta(2s)$.  The simplest lattice
VOA's constrained Epstein zeta is the Riemann zeta.

The Hecke decomposition:
$\varepsilon^r_s(V_\Lambda) = \sum_j c_j L(s, \chi_j)$.  Shadow
depth $d$ contributes $d - 1$ critical lines to the
meromorphic continuation, one for each arity stratum beyond
the Gaussian base.  The nontrivial zeros of the component
$L$-functions are the scattering resonances of the hyperbolic
Laplacian on $\cM_{1,1}$: frequencies at which the genus-$1$
sewing amplitude fails to decay.

\medskip\noindent\textbf{Genus-$2$ B\"ocherer bridge.}
For the Leech lattice VOA, the Siegel modular form
$\chi_{12} = \mathrm{SK}(f_{22})$ is the Saito--Kurokawa lift
of the unique weight-$22$ cuspidal eigenform, and the
MC-derived second Fourier coefficient $c_2 \approx
-1.918 \times 10^{-6} \neq 0$ is the first nonzero central
$L$-value produced by the shadow obstruction tower,
connecting $\Theta_\cA$ to genus-$2$ Siegel modular forms via
the Furusawa--Morimoto theorem.


% ====================================================================

% ====================================================================
% TRACK A: section 9 (scalar saturation and the Koszulness programme)
% ====================================================================

\section*{9.\quad Scalar saturation and the Koszulness programme}
% ====================================================================

The bar construction works for all chiral algebras.  It works best,
the cobar inverts the bar, the spectral sequence collapses, the
resolution is linear, precisely on the Koszul locus.  Characterizing
this locus is the central foundational question introduced in~\S3.3
and anticipated throughout sections~4 and~5.  Two complementary
results: scalar saturation (the MC element is determined by a single
scalar for all standard families) and the Koszulness meta-theorem
(twelve characterisations, ten unconditionally equivalent, one
conditional on perfectness, one one-directional).

\subsection*{9.1.\enspace Scalar saturation}

For a one-parameter algebraic family $\cA_k$ with rational OPE
coefficients (structure constants rational functions of~$k$), the
universal MC element satisfies
\[
\Theta_{\cA_k}=\kappa(\cA_k)\cdot\eta_{\mathrm{univ}}\otimes\Lambda_{\mathrm{univ}}+(\text{terms in }F^3\mathfrak g^{\mathrm{mod}}_{\cA_k}).
\]
\emph{Scalar saturation universality}: higher-order terms vanish.  For
every algebraic family with rational OPE coefficients, $\Theta_{\cA_k}$
is determined by $\kappa(\cA_k)$ modulo terms in the arity-$\ge 3$
filtration.

\medskip
\noindent\textbf{Step~1: Whitehead reduction.}\enspace
For a vertex algebra with reductive symmetry at generic level, the
primitive cyclic cohomology
$H^2_{\mathrm{cyc,prim}}(\mathfrak g^{\mathrm{mod}}/F^3,d_2)$ is
computed by a Whitehead-type spectral sequence with
$E_2^{p,q}=H^p(\mathfrak g;\operatorname{Sym}^q V)$.  For reductive
$\mathfrak g$ and generic level, $H^2(\mathfrak g;\operatorname{Sym}^q V)=0$
by the Whitehead lemma, and the arity-$2$ MC class lifts uniquely.

\medskip
\noindent\textbf{Step~2: Algebraic semicontinuity.}\enspace
The vanishing locus
\[
E=\{k\in\mathbb C:H^2_{\mathrm{cyc,prim}}(\cA_k)\neq 0\}
\]
is Zariski-closed by upper semicontinuity.  At $k=\infty$ the algebra
is Heisenberg and $H^2_{\mathrm{cyc,prim}}=0$, so $E$ is a proper
Zariski-closed subset.  Explicit computation verifies $E=\varnothing$
for $\mathfrak{sl}_N$ ($N=2,\ldots,10$): no exceptional level.  The
result covers the entire standard Lie-theoretic landscape at all
non-critical levels, including admissible.

\medskip
\noindent\textbf{Scope.}\enspace
Scalar saturation applies to every family in the standard landscape:
affine Kac--Moody (all types, all non-critical levels), Virasoro (all
$c$), $\mathcal W_N$ (all $N$, all non-degenerate levels),
$\beta\gamma$ and lattice algebras (all parameters).  The residual
conjecture, that $E=\varnothing$ for every positive-energy chiral
algebra with rational OPE coefficients, is open only for
non-algebraic-family constructions.

\subsection*{9.2.\enspace The Koszulness meta-theorem}

The Koszul condition admits twelve formulations, spanning homological
algebra (PBW collapse, Ext vanishing), homotopy theory ($A_\infty$
formality, $E_2$ formality), operadic algebra
(Barr--Beck--Lurie, FM acyclicity), modular geometry (FH
concentration, tropical Koszulness), and symplectic topology
(Lagrangian criterion).  Ten are unconditionally equivalent; one (the
Lagrangian criterion) is conditional on perfectness/nondegeneracy;
one ($\mathcal D$-module purity) is one-directional (forward direction
proved, converse open).

\begin{center}
\emph{For $\cA$ positive-energy and finitely strongly generated, the
following are equivalent:}
\end{center}

\begin{enumerate}[label=\textup{(\roman*)},nosep]
\item $\cA$ is chirally Koszul.
\item The PBW spectral sequence on $\barB^{(g)}(\cA)$ collapses at $E_2$ for every $g$.
\item The minimal $A_\infty$-model of $\barB(\cA)$ (the transferred structure on $H^\bullet$) is formal: $m_n=0$ for $n\ge 3$.
\item $\operatorname{Ext}^{p,q}_\cA(\omega_X,\omega_X)=0$ for $p\neq q$.
\item The bar-cobar counit $\Omega(\barB(\cA))\to\cA$ is a quasi-isomorphism.
\item The Barr--Beck--Lurie comparison for $\barB\dashv\Omega$ is an equivalence.
\item $\int_X\cA$ is concentrated in degree~$0$ for every smooth proper $X$.
\item $\operatorname{ChirHoch}^*(\cA)$ vanishes outside degrees $\{0,1,2\}$.
\item The Kac--Shapovalov determinant $\det G_h\neq 0$ in the bar-relevant range.
\item The restriction $i_S^!\barB_n(\cA)$ to every FM boundary stratum $S$ is acyclic outside degree~$0$.
\end{enumerate}

\noindent
Under the additional hypothesis that the ambient tangent complex is
perfect and nondegenerate:

\begin{enumerate}[label=\textup{(\roman*)},nosep,start=11]
\item $\mathcal M_\cA$ and $\mathcal M_{\cA^!}$ are transverse Lagrangians in the $(-1)$-shifted symplectic deformation space $\mathcal M_{\mathrm{comp}}$.
\end{enumerate}

\noindent
Finally (xii)$\Rightarrow$(x), where:

\begin{enumerate}[label=\textup{(\roman*)},nosep,start=12]
\item Each $\barB_n(\cA)$ is pure of weight~$n$ as a mixed Hodge module, with characteristic variety aligned to FM boundary strata ($\mathcal D$-module purity).
\end{enumerate}

\noindent The converse of (xii) is open.

\medskip
\noindent\textbf{The proof circuit.}\enspace
The core equivalence chain is
(i)$\Leftrightarrow$(ii)$\Leftrightarrow$(iii)$\Leftrightarrow$(v)$\Leftrightarrow$(viii),
proved by PBW concentration $\to$ bar-cobar inversion $\to$ Hochschild
polynomial growth $\to$ Koszul-complex acyclicity, with the $A_\infty$
formality branch via HPL transfer and Keller classicality.  From this
core: (iv) by Ext spectral sequence, (vi) by conservativity with
log-FM totalization, (vii) by $\barB=\int_X$, (ix) by non-degeneracy
forcing PBW injectivity, (x) by stratum-by-stratum PBW concentration.
The Lagrangian criterion (xi) uses PTVV derived intersection theory
on the perfectness hypothesis.  The $\mathcal D$-module purity
implication (xii)$\Rightarrow$(x) uses weight filtration on the bar
complex; the converse would require concentration to force purity.
The cycle closes via (viii)$\Rightarrow$(i): Koszul-complex acyclicity
implies the defining characterisation.

\medskip
\noindent\textbf{Beyond the meta-theorem.}\enspace
Six additional characterisations are proved in the text:

\begin{enumerate}[label=\textup{(\alph*)},nosep]
\item \emph{Shadow-formality at low arity.}  The shadow obstruction tower $\Theta^{\le r}_\cA$ equals the $L_\infty$-formality obstruction tower of $\mathfrak g^{\mathrm{mod}}_\cA$ at arities $r=2,3,4$
(Proposition~\ref{prop:shadow-formality-low-arity}).  Full equivalence to~(iii) requires all arities.

\item \emph{$E_2$-formality of chiral Hochschild cohomology.}  Koszulness implies $\operatorname{ChirHoch}^*(\cA)$ is formal as an $E_2$-algebra (Deligne conjecture applied to the bar-cobar resolution).

\item \emph{Genus-$0$ curve independence.}  Koszulness on $\mathbb P^1$ implies Koszulness on every smooth proper curve (Proposition~\ref{prop:genus0-curve-independence}).  Proved independently by factorization descent.

\item \emph{PBW universality} (sufficient).  Every freely strongly generated vertex algebra is chirally Koszul (Proposition~\ref{prop:pbw-universality}).  Applies to universal algebras: $V_k(\mathfrak g)$ is Koszul at every level including critical and admissible.  For simple quotients $L_k(\fg)$ at admissible levels, status is rank-dependent: $L_k(\mathfrak{sl}_2)$ is Koszul at all admissible levels, but at $\operatorname{rk}(\fg)\ge 2$ with denominator $q\ge 3$, Koszulness fails via a $\operatorname{rk}(\fg)$-dimensional bar obstruction from the abelian Cartan
(Conjecture~\ref{conj:admissible-koszul-rank-obstruction}).  PBW universality is one-directional: not every Koszul algebra is freely strongly generated.

\item \emph{Tropical Koszulness.}  $\cA$ is chirally Koszul iff $\barB^{\mathrm{trop}}(\cA)$ is acyclic.  Proved by the weight spectral sequence on the real blowup.

\item \emph{Bifunctor obstruction decomposition} (one-directional).  Koszulness implies the RNW one-slot obstruction decomposes into independent bar-side and cobar-side components
(Theorem~\ref{thm:bifunctor-obstruction-decomposition}).  The converse is open.
\end{enumerate}

\medskip
\noindent\textbf{Status summary.}
\begin{center}
\renewcommand{\arraystretch}{1.15}
\begin{tabular}{clll}
\multicolumn{4}{l}{\emph{Meta-theorem (Theorem~\ref{thm:koszul-equivalences-meta}):}}\\[1pt]
& (i)--(x) & \textbf{10 unconditional equivalences} & proved\\
& (xi) Lagrangian & conditional on perfectness & proved (unconditional for standard landscape)\\
& (xii) $\mathcal D$-module purity & (xii)$\Rightarrow$(x) only & converse open\\[3pt]
\multicolumn{4}{l}{\emph{Additional characterisations:}}\\[1pt]
& (a) Shadow-formality & equivalence at arities $2,3,4$ & proved\\
& (b) $E_2$-formality & consequence of Koszulness & proved\\
& (c) Curve independence & independent theorem & proved\\
& (d) PBW universality & sufficient condition & proved\\
& (e) Tropical Koszulness & equivalence & proved\\
& (f) Bifunctor decomposition & one-directional & converse open
\end{tabular}
\end{center}

\subsection*{9.3.\enspace Shadow depth and operadic complexity}

The four depth classes G/L/C/M admit refinement: in general
$r_{\max}\in\{2,3,4,5,\ldots\}\cup\{\infty\}$, with class $\mathbf F_r$
(tower terminates at arity~$r$) structurally permitted for all
$r\ge 2$; $\mathbf G=\mathbf F_2$, $\mathbf L=\mathbf F_3$,
$\mathbf C=\mathbf F_4$ name the standard-landscape representatives.

\medskip
\noindent\textbf{The operadic complexity theorem}
(Theorem~\ref{thm:operadic-complexity-detailed}):
\[
r_{\max}(\cA)=A_\infty\text{-depth}(\cA)=L_\infty\text{-formality level}(\mathfrak g^{\mathrm{mod}}_\cA).
\]
$A_\infty$-depth is the minimal $r$ with $m_k=0$ for $k>r$ on
$H^\bullet(\barB(\cA))$; $L_\infty$-formality level is the minimal $r$
such that the $L_\infty$-structure on $\mathfrak g^{\mathrm{mod}}_\cA$
is formal through arity~$r$.  Three consequences: the shadow
obstruction tower \emph{is} the $L_\infty$-formality obstruction
tower; the depth classification becomes a structural invariant of the
operadic type of $\cA$; the tower of a Gaussian algebra is a quadratic
form, of a Lie algebra a Chevalley--Eilenberg complex, of a contact
algebra a cyclic $A_\infty$-structure, of a mixed algebra an infinite
$L_\infty$-tower with no polynomial truncation.

The programme extends to non-principal $\mathcal W$-algebras: the
Bershadsky--Polyakov algebra $\mathrm{BP}_k$ is chirally Koszul at generic level ($k \neq -3$) with
dual $\mathrm{BP}_{-k-6}$ and Koszul conductor
$K_{\mathrm{BP}}=196$ (a global duality invariant, not
a local ghost constant).  Nilpotent transport from hook-type seeds
covers all orbits in $\mathfrak{sl}_2$ through $\mathfrak{sl}_8$.

\subsection*{9.4.\enspace The PBW spectral sequence}

The computational engine for bar cohomology within fixed genus is the
PBW (Poincar\'e--Birkhoff--Witt) spectral sequence.  Filter
$\barB^{(g)}(\cA)$ by \emph{conformal weight}: assign each normally
ordered monomial its total weight (sum of constituent field weights).
The associated graded is the bar complex of a free-field algebra,
reducing to symmetric-function combinatorics.

The $E_1$ page at genus~$g$:
\[
E_1^{p,q}(g)=E_1^{p,q}(g{=}0)\oplus\cE^{p,q}_g,
\]
with the first summand the genus-$0$ bar cohomology determined by
Koszulness of the associated graded and $\cE^{p,q}_g$ the genus-$g$
enrichment from the Hodge bundle, arising from the regular Eisenstein
series coupled to the curvature.  \emph{Concentration}:
$E_\infty^{p,q}(g)=0$ for $q\neq 0$ at all $g\ge 1$, for all standard
families.  PBW degeneration: the spectral sequence collapses at $E_2$
and bar cohomology concentrates on the diagonal.

The PBW spectral sequence is distinct from the genus spectral sequence
of \S5.5.  PBW uses conformal weight within fixed genus to compute bar
cohomology groups; the genus spectral sequence uses the genus
filtration on the deformation algebra to control the MC lift across
genera.  One computes \emph{what}; the other determines \emph{whether}.

\subsection*{9.5.\enspace Chiral Hochschild cohomology (Theorem~H)}

The chiral Hochschild complex $\operatorname{ChirHoch}^\bullet(\cA)$
is the derived endomorphism algebra of the identity functor on
$\cA$-modules, computed by the bar construction with bimodule
coefficients.  At genus~$0$,
\[
\operatorname{ChirHoch}^\bullet(\cA)\simeq R\Hom_{\cA\text{-bimod}}(\cA,\cA)\simeq\operatorname{Ext}^\bullet_{\cA^e}(\cA,\cA),
\]
where $\cA^e=\cA\otimes\cA^{\mathrm{op}}$ is the enveloping chiral
algebra.

For $\cA=\mathcal W^k(\mathfrak g,f_{\mathrm{prin}})$ a principal
$\mathcal W$-algebra at generic level, the chiral Hochschild cohomology
is concentrated in degrees $\{0,1,2\}$:
\[
\operatorname{ChirHoch}^n(\mathcal W^k(\mathfrak g))=0\text{ for }n>2,\qquad P(t)=1+t^2.
\]
The deformation class in $\operatorname{ChirHoch}^2=\mathbb C$ is the
level deformation $k\mapsto k+\epsilon$.  The concentration bound is
the de Rham amplitude on a curve ($\dim X=1$); the continuous
Lie algebra cohomology
$H^*_{\mathrm{cont}}(\operatorname{Lie}(\mathcal W),\operatorname{Lie}(\mathcal W)_0;\mathbb C)=\mathbb C[\Theta_1,\ldots,\Theta_r]$
(Gelfand--Fuchs, $\deg\Theta_i=m_i+1$) is a different, unbounded
invariant.  For $\mathfrak g=\mathfrak{sl}_2$ (Virasoro):
$P(t)=1+t^2$.  For $\mathfrak g=\mathfrak{sl}_3$ ($\mathcal W_3$):
$P(t)=1+t^2$.  Total dimension $2$, never $\mathbb C[\Theta]$.

On the Koszul locus, $\operatorname{ChirHoch}^\bullet(\cA)$ carries an
$E_2$-algebra structure that is formal.  The Connes periodicity
operator $S\colon\operatorname{ChirHoch}^n\to\operatorname{ChirHoch}^{n-2}$
intertwines Koszul duality: $S\circ\operatorname{KD}=\operatorname{KD}\circ S$.
At genus~$0$, the chiral Hochschild complex identifies with the bulk
chiral homology complex of Volume~II.  The bulk-boundary-line triangle
is the genus-$0$ shadow of this identification:
$\cA_{\mathrm{bulk}}\simeq Z_{\mathrm{der}}(\cB_\partial)\simeq\operatorname{ChirHoch}^\bullet(\cA^!)$.

\subsection*{9.6.\enspace The shadow algebra}

The \emph{shadow algebra} $\cA^{\mathrm{sh}}$ is the cohomology of the
modular cyclic deformation complex as a graded-commutative ring,
\[
\cA^{\mathrm{sh}}:=H_\bullet(\Defcyc^{\mathrm{mod}}(\cA))=\bigoplus_{r\ge 2,g\ge 0}\cA^{\mathrm{sh}}_{r,g}.
\]
The single MC element $\Theta_\cA$ decomposes by arity and genus.
Named shadows are projections:
\[
\kappa(\cA)=\cA^{\mathrm{sh}}_{2,0},\quad\mathfrak C(\cA)=\cA^{\mathrm{sh}}_{3,0},\quad\mathfrak Q(\cA)=\cA^{\mathrm{sh}}_{4,0},\quad\operatorname{Sh}_r(\cA)=\cA^{\mathrm{sh}}_{r,0}\text{ for }r\ge 5.
\]
The shadow algebra organises the entire hierarchy; it is a homotopy
invariant, computed by the Weil homomorphism of the Chern--Weil
dictionary (\S5.4).

The all-arity master equation in the shadow algebra:
$\nabla_H(\operatorname{Sh}_r)+o^{(r)}=0$ for all $r\ge 3$.  The
obstruction class $o^{(r)}$ lies in $\cA^{\mathrm{sh}}_{r,0}$ and
measures the failure of $\Theta^{\le r-1}_\cA$ to extend to arity~$r$.

\subsection*{9.7.\enspace BV/BRST identification at genus~$0$}

At genus~$0$, the bar differential of a chiral algebra on $\mathbb C$
is the BRST differential of the associated holomorphic twist,
$(d_{\barB})_{g=0}=Q_{\mathrm{BRST}}$.  The bar complex
$\barB^{(0)}_{\mathbb C}(\cA)$ is the BRST complex of the 2d chiral
algebra, and bar-cobar inversion
$\Omega(\barB(\cA))\xrightarrow{\sim}\cA$ (Theorem~B) says the BRST
cohomology processed through the cobar functor reconstructs the
original algebra.

At genus~$1$, the identification acquires curvature:
$d^2_{\mathrm{fib}}=\kappa(\cA)\cdot\omega_1$ is the one-loop anomaly,
and the curved bar complex is the curved BRST complex of the theory
on an elliptic curve.  The complementarity
$Q_1(\cA)\oplus Q_1(\cA^!)=H^*(\overline{\cM}_{1,1},\cZ_\cA)$ is the
genus-$1$ anomaly cancellation condition.

At genus $g\ge 2$, the identification with Costello--Gwilliam
requires renormalisation: higher-loop integrals over
$\overline{\cM}_{g,n}$ need regularisation beyond the algebraic
bar-cobar framework, and the correct receptacle is the coderived
category.

For the Heisenberg algebra, the identification is proved at all genera
by four independent paths: zeta-regularisation, Selberg zeta function,
Grothendieck--Riemann--Roch, direct bar computation.  The shadow
obstruction tower equals the $L_\infty$-formality obstruction tower
at all arities.

\medskip
\noindent\textbf{Genus-two shell activation as depth diagnostic.}\enspace
The genus-two shell profile independently confirms the depth
classification:

\begin{center}
\renewcommand{\arraystretch}{1.15}
\begin{tabular}{lccc}
& $\Theta^{(2)}_{\mathrm{loop}^2}$ & $\Theta^{(2)}_{\mathrm{sep}\circ\mathrm{loop}}$ & $\Theta^{(2)}_{\mathrm{pf}}$\\
\hline
\textbf{G} (Heisenberg, lattice) & $\checkmark$ & --- & ---\\
\textbf{L} (affine $\widehat{\mathfrak g}_k$) & $\checkmark$ & $\checkmark$ & $\checkmark$\\
\textbf{C} ($\beta\gamma$) & $\checkmark$ & --- & ---\\
\textbf{M} (Virasoro, $\mathcal W_N$) & $\checkmark$ & $\checkmark$ & $\checkmark$
\end{tabular}
\end{center}

\noindent
G and C activate only the loop$\circ$loop shell; L and M activate all
three including planted-forest.  The diagnostic distinguishing L from M
is the critical discriminant $\Delta=8\kappa S_4$: zero for L
($S_4=0$), nonzero for M.

\medskip
\noindent\textbf{Structural refinements.}\enspace
The three bar complexes on a chiral algebra (ordered tensor $T^c$,
symmetric $\operatorname{Sym}^c$, Chevalley--Eilenberg
$\operatorname{Lie}^c$) sit in a chain
$\operatorname{Lie}^c\hookrightarrow\operatorname{Sym}^c\hookrightarrow T^c$
with distinct coalgebra structures at each level
(Theorem~\ref{thm:three-bar-complexes}).  The averaging map
$\mathrm{av}\colon T^c\to\operatorname{Sym}^c$ from ordered to
symmetric is surjective but not split: the kernel at arity~$3$
contains the Drinfeld associator $\Phi_{\mathrm{KZ}}(\cA)$, and the
fibre of splittings is a torsor for $\mathrm{GRT}_1$.  The $E_1$
ordered bar complex is therefore the natural primitive object; the
modular/symmetric story is its $\Sigma_n$-coinvariant image
(Theorem~\ref{thm:e1-primacy}).  The genus-$3$ cross-channel
correction for $\mathcal W_3$ has been computed by summing over all
$42$ stable graphs of $\overline{\cM}_{3,0}$
(Proposition~\ref{prop:w3-genus3-cross-channel}): the cross-channel
tower accelerates with genus and dominates the scalar part
$\kappa\cdot\lambda_g^{\mathrm{FP}}$ by genus~$4$
\textsc{(all-weight, with cross-channel correction)}.

On the open/closed side, the Koszul dual cooperad
$\mathsf{SC}^{\mathrm{ch,top},!}$ decomposes into three sectors:
closed ($\dim=(n-1)!$, Lie cooperad), open ($\dim=m!$, Ass cooperad),
and mixed ($\dim=(k-1)!\binom{k+m}{m}$).  The mixed sector
encodes the bulk-to-boundary module structure by which the chiral
derived centre acts on the boundary algebra.

The BV/BRST${}={}$bar identification at higher genus
(Conjecture~\ref{conj:master-bv-brst}): the chain-level obstruction
(harmonic propagator correction) is resolved for classes G, L, and C,
for G and L by the Heisenberg sewing theorem and its extension to
affine algebras, for C by contact decoupling.  For class M (Virasoro,
$\mathcal W_N$) the identification holds on bar cohomology but fails
at the chain level: the harmonic propagator produces a nonvanishing
quartic obstruction, so a coderived reformulation is needed.

% ====================================================================
% TRACK B: sections 10--13 (open/closed, PVA quantization, holography,
% completion and frontier)
% ====================================================================

\section*{10.\quad The open/closed world (Volume II)}
% ====================================================================

The bar complex lives naturally in three dimensions.  On
$\bC_z \times \bR_t$, the holomorphic direction produces the
closed-string bar differential of Sections~1--9, extracting OPE
residues from collisions in the holomorphic plane.  The
topological direction produces the open-string coproduct,
splitting an ordered tensor sequence at a cut point of the real
line.  The two are unified by the \emph{holomorphic-topological
Swiss-cheese operad} $\SCchtop$: closed colour
$\mathrm{ch} = \FM_k(\bC)$ governs holomorphic factorization,
open colour $\mathrm{top} = E_1(m)$ governs topological
factorization, and the open/closed mixed sector encodes the
action of bulk on boundary.  Without this operadic habitat, the
bar coproduct has no definition and the line-operator category
has no foundation.

The primitive object of 3d HT quantum field theory on
$\bC_z \times \bR_t$ is not the bar complex but the open/closed
factorization dg-category $\cC$ on the bordified curve
$\widetilde X_D$.  Objects are boundary conditions, morphisms
are open-string states, and a vacuum $b$ produces a boundary
algebra $A_b = \End_\cC(b)$ only up to Morita equivalence.  The
governing operad is $\SCchtop$; it is homotopy-Koszul (proved
via Kontsevich formality), so the bar-cobar adjunction on
$\SCchtop$-algebras is a Quillen equivalence.

The bar complex of Volume~I is the coalgebraic shadow of $\cC$:
its differential encodes the closed colour, and its coproduct
the open colour.  Three objects must never be conflated:
\begin{enumerate}[label=\textup{(\roman*)},nosep]
\item the \emph{bar complex} $\barB(A_b)$ classifies twisting
morphisms, universal couplings between $A_b$ and $A_b^!$;
\item the \emph{chiral derived centre}
$\cZ^{\der}_{\ch}(\cC) =
R\Hom_{\Fun(\cC,\cC)}(\Id, \Id)$ is the universal bulk,
computed by chiral Hochschild cochains in any chart;
\item the Koszul dual $A_b^!$ governs the line-operator
category $A_b^!\text{-mod}$.
\end{enumerate}
The primitive hierarchy
\[
\cC_{\op}
\;\xrightarrow{\End(b)}\;
A_b
\;\xrightarrow{\barB}\;
\barB(A_b)
\;\xrightarrow{\Theta}\;
\Theta_\cA
\;\xrightarrow{+\mu^{M_j}}\;
\Theta^{\oc}_\cA
\]
records successive lossy functors: $\cC_{\op}$ is Morita
invariant; $A_b$ depends on $b$; $\barB$ forgets the bulk;
$\Theta_\cA$ is the closed-sector MC element;
$\Theta^{\oc}_\cA$ restores open/closed data.  Modularity is
determined by the open sector: a cyclic trace on $\End(b)$
seeds the Calabi--Yau structure, the annulus identification
$\int_{S^1} \cC \simeq \HH_\bullet(\cC)$ is the first modular
shadow, and clutching along boundary circles raises genus.

Volume II proceeds in four stages.  Stage~1 (local one-colour
theorem: Swiss-cheese universality of the derived centre),
Stage~2 (open primitive: $\cC$ with Morita invariance), and
Stage~3 (globalisation on tangential log curves $(X, D, \tau)$)
are proved (Stage~3 locally; global extension modest).  In Stage~4 (modularisation), the open/closed MC
element $\Theta^{\oc}_\cA$ is constructed, the modular MC
equation holds at all genera, and the projection principle
recovers both the closed-sector genus tower and the
open-sector module structures.  The annulus trace is proved;
the independent higher-genus derivation and chain-level
cooperad compatibilities remain open.

\subsection*{10.1.\enspace The Swiss-cheese operad
$\SCchtop$}

Two colours $\{\mathrm{ch}, \mathrm{top}\}$, with operation
spaces
\begin{align*}
\SC(\mathrm{ch}^k; \mathrm{ch}) &= \FM_k(\bC),\\
\SC(\mathrm{ch}^k, \mathrm{top}^m; \mathrm{top})
&= \FM_k(\bC) \times E_1(m),\\
\SC(\ldots, \mathrm{top}, \ldots; \mathrm{ch}) &= \varnothing.
\end{align*}
Closed-to-closed is the Fulton--MacPherson compactification in
$\bC$; the mixed direction combines holomorphic configurations
acting on the fibre above a topological interval; open-to-closed
is forbidden.  Operadic composition is the boundary stratification
of FM spaces in the holomorphic direction and nested interval
inclusion in the topological direction.

A bar element of tensor degree $k$ lives on the product
$\FM_k(\bC) \times \Conf_k(\bR)$.  The bar differential extracts
OPE residues along collision divisors of the first factor: when
points $z_i$, $z_j$ collide, the residue of
$\eta_{ij} = d\log(z_i - z_j)$ yields the mode $a_{(n)}b$.  The
coproduct splits the ordered sequence at a cut point of the
second factor: for $t_1 < \cdots < t_k$ with $t_p < c < t_{p+1}$,
$(a_1\otimes\cdots\otimes a_k)$ decomposes as
$(a_1\otimes\cdots\otimes a_p)\otimes(a_{p+1}\otimes\cdots\otimes a_k)$.
The bar complex with both structures is a
$\SCchtop$-algebra.

\subsection*{10.2.\enspace Homotopy-Koszulity}

The classical Swiss-cheese operad $\SC$ is Koszul
(Livernet--Voronov, Ginzburg--Kapranov), with quadratic
presentation and distributive-lattice criterion on the poset
of nested disks.  Kontsevich formality supplies a quasi-iso
$\Phi\colon \SCchtop \xrightarrow{\sim} \SC$ from
configuration-space integrals on FM spaces with propagators,
extended by including $\mathbb H$ and $\bR$.  The bar-cobar
adjunction preserves quasi-isos, and two-out-of-three in the
operadic model category gives
$\Omega\mathbf B(\SCchtop) \xrightarrow{\sim} \SCchtop$.  The
bar-cobar adjunction on $\SCchtop$-algebras is therefore a
Quillen equivalence, with three consequences:
\begin{enumerate}[label=(\roman*),nosep]
\item every $\SCchtop$-algebra admits a cofibrant bar-cobar
resolution;
\item the homotopy category of $\SCchtop$-algebras is
equivalent to that of coalgebras over the Koszul dual
cooperad;
\item the line-operator category $A_b^!\text{-mod}$ is
unconditionally defined without finiteness hypotheses on $\cA$.
\end{enumerate}

\subsection*{10.3.\enspace The $A_\infty$ chiral algebra structure}

The full Swiss-cheese MC element $\alpha_T \in \MC(\fg^{\SC}_T)$
lives in the convolution algebra
$\fg^{\SC}_T = \Hom_\Sigma(\mathbf B(\SCchtop),
\End_{(B_\partial, \cC_{\mathrm{line}})})$.  Its closed-colour
projection $(m_k)_{k \ge 1}$ encodes the $A_\infty$ chiral
algebra on $\cA$.  Each $m_k\colon \cA^{\otimes k} \to
\cA(\!(\lambda_1)\!)\cdots(\!(\lambda_{k-1})\!)$ is a
configuration-space integral on
$\FM_k(\bC) \times \Conf_k(\bR)$, with the holomorphic factor
producing OPE singularities and the topological factor
producing the ordering.  Summing over trees $T$ with $k$ leaves,
$m_T(a_1,\ldots,a_k) = \int_{\FM_k(\bC)}\omega_T \prod_{v}
a_{(n_v)} b_v$.  The Maurer--Cartan equation is the spectral
Stasheff identity: Stokes' theorem on boundary strata of
$\FM_k(\bC)$ produces the $A_\infty$ relations.

\subsection*{10.4.\enspace PVA descent (Theorem G)}

On cohomology, the operations simplify.  The regular part of
$m_2$ is a commutative associative product
$\mu(a,b) = \Reg_{z_1 \to z_2} m_2(a(z_1), b(z_2))$;
commutativity follows from the $\bZ/2$ involution of
$\FM_2(\bC)$, associativity from the three FM$_3$ boundary
strata.  The singular part defines the lambda-bracket
$\{a_\lambda b\} = \Res_{z_1 \to z_2} e^{\lambda(z_1 - z_2)}
m_2(a(z_1), b(z_2))$; the exponential extracts all pole orders.
All $m_{k \ge 3} = 0$ on cohomology because $\Conf_k(\bR)$ is
contractible for $k \ge 3$.

The Arnold relation
$\eta_{12}\wedge\eta_{23} + \eta_{23}\wedge\eta_{31}
+ \eta_{31}\wedge\eta_{12} = 0$ on $\FM_3(\bC)$ is equivalent
to the three codim-1 boundary divisors of $\overline M_{0,4}$
summing to zero in homology; its residue against
$a\otimes b\otimes c$ produces the Jacobi identity for
$\{-_\lambda -\}$.  The three-face Stokes identity applied to
$\omega = \eta_{12}\wedge\eta_{23}\cdot a(z_1)b(z_2)c(z_3)$
produces the Leibniz rule $\{a_\lambda(b\cdot c)\} =
\{a_\lambda b\}c + b\{a_\lambda c\}$.  The exchange cylinder
on $\FM_2(\bC)\times E_1(2)$ gives skew-symmetry
$\{a_\lambda b\} = -\{b_{-\lambda - \partial} a\}$.  Thus
$(\mu, \{-_\lambda -\})$ is a $(-1)$-shifted Poisson vertex
algebra, each axiom in one line from the geometry of FM.

(Vol~I convention: we state the PVA in OPE modes; Vol~II uses
lambda-brackets throughout, and the two conventions are
translated by the standard
$\{a_\lambda b\} = \sum_{n \ge 0}\lambda^n a_{(n)}b/n!$.)

\subsection*{10.5.\enspace The Heisenberg algebra in Swiss-cheese coordinates}

For $\cH_k$: $m_2(J, J) = k \cdot \mathbf 1$,
$m_{k \ge 3} = 0$, $\{J_\lambda J\} = k\lambda$.  All higher
operations vanish because the OPE has no simple pole.  The
spectral $R$-matrix is the scalar $R(z) = \exp(k/z)$; YBE
holds trivially.  The PVA is the free bosonic PVA, the
universal enveloping PVA of the rank-one abelian Lie conformal
algebra.

\subsection*{10.6.\enspace The Koszul triangle}

From $\cC$ on $\widetilde X_D$ three independent objects arise:
\begin{enumerate}[label=\textup{(\roman*)},nosep]
\item Boundary $A_b = \End_\cC(b)$: a chart.
\item Bulk $\cZ^{\der}_{\ch}(\cC) =
R\Hom_{\Fun(\cC,\cC)}(\Id, \Id)$: the universal bulk, Morita
invariant, computed by chiral Hochschild cochains.
\item Lines
$A_b^!{}_\infty = D_{\Ran}(\barB(A_b))$: the homotopy Koszul
dual, whose cohomology on the Koszul locus recovers the strict
dual $A_b^! = (H^*(\barB(A_b)))^\vee$.  The module category
$A_b^!\text{-mod}$ with spectral braiding $R(z)$ is the
braided tensor category of line operators.
\end{enumerate}
The triangle of identifications
\[
\cZ^{\der}_{\ch}(\cC)
\;\simeq\;
\HH^\bullet(A_b)
\;\simeq\;
\mathrm{ChirHoch}^\bullet(A_b^!),
\qquad
A_b^!\text{-mod} = \{\text{line operators}\}
\]
relates three distinct functors: endomorphisms (bulk),
Ran-space Verdier duality (lines), and identity-bimodule
endomorphisms (centre).  The bar complex classifies twisting
morphisms, not the bulk.

\medskip\noindent\emph{Spectral braiding.}  Stokes' theorem on
$\FM_3(\bC)$ applied to $\alpha\star\alpha\star\alpha$ decomposes
the boundary into three pairwise collision divisors, one for
each side of the Yang--Baxter equation
$R_{12}R_{13}R_{23} = R_{23}R_{13}R_{12}$.

\medskip\noindent\emph{Bulk $=$ derived centre.}  Let
$\cB_\partial$ be the boundary factorization algebra on $\bR$.
Its derived centre
$Z_{\der}(\cB_\partial) = R\Hom_{\cB_\partial\text{-bimod}}
(\cB_\partial, \cB_\partial)$
computes endomorphisms of the identity bimodule.  The
$E_1$-structure gives
$Z_{\der}(\cB_\partial) \simeq \HH^\bullet(\cB_\partial)$ as
$E_2$-algebra.  The Swiss-cheese structure identifies the
holomorphic bulk with this $E_2$-algebra.

\subsection*{10.7.\enspace Curved Swiss-cheese at $g \ge 1$}

At genus $g \ge 1$, the bar complex carries curvature:
\[
d_{\mathrm{fib}}^2 \;=\; \kappa(\cA)\cdot\omega_g,
\]
with $\omega_g$ the Arakelov $(1,1)$-form representing the first
Chern class of the relative dualising sheaf.  The coproduct
remains strictly coassociative: the interval $\bR$ has no
topology, so the curvature lives entirely in the holomorphic
direction.  The total corrected differential
$D_g = d_{\mathrm{fib}} + \nabla^{\GM}_g$ restores flatness by
coupling the ordering to the Hodge bundle:
$d_{\mathrm{fib}}^2 + [d_{\mathrm{fib}}, \nabla^{\GM}]
+ (\nabla^{\GM})^2 = \kappa\omega_g - \kappa\omega_g = 0$.

The derived category must be replaced by the coderived category
$\mathsf D^{\co}(\cA^{(g)})$: ordinary $d^2 = 0$ fails at $g \ge
1$, every object is acyclic in the derived sense, and the
correct homotopy theory comes from totalising the curvature
bicomplex.  Passage from derived to coderived is the categorified
form of passage from $d^2 = 0$ to curved MC
$d\alpha + \alpha^2 + \kappa\omega = 0$.

\subsection*{10.8.\enspace Lagrangian self-intersection}

Volume~II identifies the geometric substrate of the bar
construction: the bar complex is the structure sheaf of the
Lagrangian self-intersection
$\mathfrak S = \cL \times_\cM \cL$ in a $(-2)$-shifted
symplectic stack.  Every theorem is a face of that geometry:
Theorem~A is the groupoid comodule-module adjunction for
$\mathfrak S$; Theorem~B (bar-cobar inversion) reconstructs
the Lagrangian from its clean self-intersection; Theorem~C
(complementarity) is two Lagrangians whose intersection
carries a $(-1)$-shifted symplectic structure; Theorem~D is
the Kodaira--Spencer class; Theorem~H (Hochschild polynomial
growth) is the HKR theorem for the Lagrangian embedding.

\subsection*{10.9.\enspace Factorization-primary hierarchy}

$\SCchtop$ is a \emph{local model}: the formal completion of a
factorisation structure at collision points.  The global truth
is a factorisable $\cD$-module on $\Ran(\Sigma_g)$ over
$\overline\cM_g$.  Six layers:
\begin{enumerate}[label=(\roman*),nosep]
\item Geometric substrate: smooth and stable nodal curves,
$\Ran$ space, $\overline\cM_{g,n}$.
\item Factorizable $\cD$-modules (Beilinson--Drinfeld).
\item HT product geometry $\Sigma_g \times \bR$.
\item Bar as factorization coalgebra on $\Ran(X)$.
\item Koszul duality as factorization duality; Feynman
involutivity $\mathrm{FT}^2 \simeq \id$.
\item Three models as three gauges: flat associated graded
($D_0^2 = 0$, operadic/local), corrected holomorphic
($D^{(g)2} = 0$, Fay trisecant, still local), curved geometric
($d_{\mathrm{fib}}^2 = \kappa\cdot\omega_g$, Arakelov, global).
\end{enumerate}
The curvature measures what the local extraction loses: it is
the monodromy of the $\cD$-module connection around $B$-cycles
of $\Sigma_g$.  The modular bar complex $\barB_{\mod}(\cA)$ is
flat because it couples all genera simultaneously via $D_1$;
the genus-fixed bar complex is curved because restriction
decouples this connection.

The relative Feynman transform
$\mathrm{FT}_{\mathsf{Com}_{\mod}/\SCchtop}$ mediates between
factorisation (global) and operadic (local) viewpoints.  The
modular bar coalgebra is an algebra over this transform;
homotopy-involutivity plus factorisation Koszul duality gives
the flat-derived / curved-coderived comparison at all genera.

\subsection*{10.10.\enspace Four-stage architecture}

\emph{Stage~1.}  $\SCchtop$ homotopy-Koszul; bar-cobar is
Quillen-equivalent; chiral derived centre
$\simeq C^\bullet_{\ch}(A_b, A_b)$ is the universal bulk.
\textsc{Proved.}

\emph{Stage~2.}  Primitive object is the open factorisation
dg-category $\cC$; $A_b$ is a chart; bulk is Morita invariant.
\textsc{Proved.}

\emph{Stage~3.}  Tangential log curves $(X, D, \tau)$ provide
bordered FM compactification; local-global bridge lifts
operadic structure to factorisable $\cD$-module on $\Ran(X)$.
\textsc{Proved locally; modest globally.}

\emph{Stage~4.}  Trace and clutching on the open sector produce
the full open/closed MC element $\Theta^{\oc}$.  Annulus trace
is proved; full chain-level clutching is programme.

Five negative principles: (N1) bar $\neq$ bulk; (N2) boundary
algebra is a chart; (N3) tangential log geometry required;
(N4) modularity is trace plus clutching on the open sector;
(N5) dg-shifted Yangian is proved for affine lineage only,
extension to non-standard families is programme.


% ====================================================================
\section*{11.\quad Modular PVA quantization}
% ====================================================================

The open/closed world supplies a genus-zero datum: the
Swiss-cheese MC element $\alpha_T$, the PVA
$(\mu, \{-_\lambda -\})$ on cohomology, the KZ/BPZ shadow
connections at arity $n$.  All are tree-level.  How does this
structure extend across the genus filtration?  The answer is
a genus-by-genus obstruction theory whose input is the genus-$0$
PVA and whose output is the full modular MC element.  The
nonseparating degeneration operator $D_1$ is the primitive
genus-raising differential; it squares to zero, anticommutes
with the tree-level differential $D_0$, and defines a spectral
sequence whose obstructions are the quantisation anomalies.

\subsection*{11.1.\enspace One-edge principle}

Every nontrivial modular obstruction theory must contain a
primitive genus-raising operator; otherwise any genus-$0$ MC
element lifts tautologically to all genera.  The one-edge
expansion
\[
D_{\exp} \;=\; D_{\sep} + D_{\nsep},
\qquad
D_0 \;=\; D_{\mathrm{int}} + D_{\sep},
\qquad
D_1 \;=\; D_{\nsep}
\]
separates genus-preserving from genus-raising degenerations.
$(D_0 + D_1)^2 = 0$ stratifies by genus shift:
$D_0^2 = 0$, $D_0 D_1 + D_1 D_0 = 0$, $D_1^2 = 0$.  Trees
govern chirality; loops govern modularity.  The three
identities are codim-2 cancellation stratified by genus.

\subsection*{11.2.\enspace Modular bar datum}

A modular bar datum on a reduced complete graded collection
$F \mapsto C(F)$ consists of (i)~an internal differential
$d\colon C(F) \to C(F)$ with $d^2 = 0$;
(ii)~for every one-edge expansion
$\epsilon \in \mathrm{OE}(F)$ (separating or nonseparating), a
degree-$+1$ map
\[
\Delta_\epsilon\colon
sC(F)
\;\longrightarrow\;
\orline(\epsilon) \otimes \bigotimes_{u \in V(\epsilon)}
sC(\Fl_\epsilon(u));
\]
(iii)~the codim-$2$ cancellation
$\sum_{\epsilon_1, \epsilon_2} \sgn\,
\Delta_{\epsilon_2}\circ\Delta_{\epsilon_1} = 0$
for every two-edge graph; and
(iv)~$d\circ\Delta_\epsilon + \Delta_\epsilon\circ d = 0$.
The total $D = d + \sum_\epsilon \Delta_\epsilon$ satisfies
$D^2 = 0$ by three mechanisms: $d^2 = 0$, anticommutation,
and the codim-$2$ axiom.  This is the algebraic incarnation
of the cell decomposition of $\overline\cM_{g,n}$ into
stable-graph strata.

\subsection*{11.3.\enspace Complete modular bar coalgebra}

For a connected stable graph $\Gamma$ of genus $g$, set
\[
C[\Gamma]
\;=\;
\orline(\Gamma) \otimes
\bigotimes_{v \in V(\Gamma)} sC(\Fl(v)),
\qquad
\orline(\Gamma) = \bigotimes_{e} \orline(e).
\]
The complete modular bar coalgebra is
\[
\mathbf B_{\mod}(C)
\;=\;
\prod_{[\Gamma] \in \mathsf{StGraph}^{\conn}}
(C[\Gamma])_{\Aut(\Gamma)},
\]
completed over total genus $g(\Gamma) = b_1 + \sum_v g(v)$.
Edge contraction is the reduced coproduct; the coalgebra is
conilpotent.

\subsection*{11.4.\enspace Strict modular deformation algebra}

$L_{\mod}(C) := \Coder(\mathbf B_{\mod}(C))[-1]$ is complete
filtered by genus, with convolution
$L_{\mod}(C) \cong \Hom(\mathbf B_{\mod}(C), sC)[-1]$ under
commutator bracket.  The algebra is pro-nilpotent in
$F^1$, and the associated graded inherits a dg Lie structure
at each genus.

\subsection*{11.5.\enspace Genus-$1$ obstruction theorem}

For a genus-$0$ MC element
$\Theta_0 \in \MC(\gr_F^0 L)$, the genus-$1$ obstruction is
\[
\Ob_1(\Theta_0)
\;=\;
[D_1 \Theta_0]
\;\in\;
H^2((\gr_F^1 L)^{\Theta_0}, d_0^{\Theta_0}).
\]
$D_1\Theta_0$ is a $d_0^{\Theta_0}$-cocycle by
$D_0 D_1 + D_1 D_0 = 0$ and the MC equation for $\Theta_0$.
The class vanishes iff $\Theta_0$ lifts to
$\Theta_0 + \hbar\Theta_1$ satisfying the modular MC equation
modulo $\hbar^2$.

In PVA coordinates, the nonseparating one-edge operator acts
as the reduced odd Laplacian: for cyclic $(A, m_k, \langle-,-\rangle)$
with basis $\{e^i\}$ and dual basis $\{e_i\}$,
\[
\ell_1^{(1)}(a)
\;=\;
\hbar\Delta_{\mathrm{cyc}}(a)
\;=\;
\hbar\sum_i m_2(e^i, m_2(a, e_i)).
\]
The full $\Ob_1$ is the cohomology class of
$\Delta_{\mathrm{cyc}}(\Theta_0)$.

\subsection*{11.6.\enspace $\cW_3$ computation}

A filtered quadratic PVA is generated by $u^1, \ldots, u^N$ with
$\lambda$-bracket $\Pi^{ab}(\partial)$ of total generator degree
$\le 2$.  Classical $\cW_3$ fits: generators $T$ (weight $2$),
$W$ (weight $3$), with the composite
$\Lambda = :TT: - (3/10)\partial^2 T$ controlling the nonlinear
channel; the $WW$ OPE has the $c = -22/5$ pole in
$16/(22 + 5c)$ where $\Lambda$ becomes null.

The triangular $W$-normal form vanishing theorem: in the local,
translation-invariant, filtration-preserving sector,
$\Ob_1(\Theta_0) = 0$ for classical $\cW_3$.  The relevant
cohomology is
$H^1((\gr^1 L)^{\Theta_0}, d_0^{\Theta_0})
\cong \bC\cdot\partial_c P_c$,
one-dimensional and spanned by the central-parameter direction.
Every genus-$1$ lift is gauge equivalent to a renormalisation
$c \mapsto c + \hbar c_1$; the genus-$1$ quantum $\cW_3$ is the
unique quantisation in the direction $\partial_c$.

\subsection*{11.7.\enspace Curved extension via $S$-tails}

In the inhomogeneous quadratic setting (Gui--Li--Zeng), the
correct input is a twisted pair $(B, B^\circ, S)$: a
CDG-algebra $B$, a CDG-coalgebra $B^\circ$, and curvature $S$.
An $S$-tail expansion attaches a univalent $S$-decorated vertex
to a corolla, followed by contraction of the new internal edge.
Adjoining $S$-tail expansions to one-edge expansions with the
same codim-$2$ compatibilities gives a curved modular bar
object.  The curved MC equation is
$\mu((S + \alpha) \boxtimes (S + \alpha)) = 0$; at genus $0$
it is the curved $A_\infty$ MC equation with $m_0 = S$, at
genus $1$ it is $\Delta_{\mathrm{cyc}}(S + \alpha) = 0$, and
at genus $g$ it is the BV quantum master equation
$\hbar\Delta(e^{(S+\alpha)/\hbar}) = 0$.

\subsection*{11.8.\enspace Genus-$0$ Gui--Li--Zeng comparison}

At tree level, the modular MC equation reduces to the
classical twisting morphism equation.  For a Koszul pair
$(\cA, \cA^!)$ with $\cA^!$ a CDG algebra with curvature $S$,
\[
\MC_0(L_{\mod})
\;\cong\;
\Hom_{\mathsf{dg}}(\cA, B^\circ)
\;\hookrightarrow\;
\MC(\cA^! \widehat\otimes \cB).
\]
The modular theory is the genus extension of the classical
Gui--Li--Zeng picture.


% ====================================================================
\section*{12.\quad Holographic modular Koszul duality}
% ====================================================================

The open/closed world (Section~10) supplies the 3d geometric
framework; the modular PVA quantisation programme
(Section~11) supplies the obstruction theory.  The holographic
modular Koszul datum packages both into a single algebraic
object: a six-component structure that records the boundary
algebra, its Koszul dual, the line-operator category, the
spectral $R$-matrix, the universal MC element, and the modular
shadow connection.  Every ingredient constructed so far (the
bar complex, the shadow tower, the Swiss-cheese structure, the
PVA descent, the genus expansion) is a projection of this
datum.  The datum is the smallest package from which the full
correspondence reconstructs.

\subsection*{12.1.\enspace The holographic modular Koszul datum}

Every 3d HT system $T$ is controlled by
\[
\cH(T)
\;=\;
(\cA,\; \cA^!,\; \cC,\; r(z),\; \Theta_\cA,\; \nabla^{\hol}),
\]
where:
\begin{enumerate}[label=\textup{(\roman*)}]
\item $\cA$ is the boundary chiral algebra on a curve $X$,
restricted from the holomorphic boundary
$\{0\} \times X \subset \bR_+ \times X$.
\item $\cA^!$ is the strict chiral Koszul dual,
$\cA^! = (\cA^{\mathrm i})^\vee$ with $\cA^{\mathrm i} = H^*(\barB(\cA))$.
Holographically, $\cA^!$ is a dg-shifted Yangian: the bar
complex $\barB(A_b)$ is a $\SCchtop$-algebra, and $\cA^!$ is
the $E_1$-projection of the $\SCchtop$-dual.
\item $\cC \simeq \cA^!\text{-mod}$ is the braided tensor
category of line operators: objects are topological defect
lines, morphisms junction operators, braiding from crossing
in the holomorphic plane.
\item $r(z) \in \cA^! \otimes \cA^!(\!(z^{-1})\!)$ is the
classical $r$-matrix satisfying CYBE, the binary genus-$0$
shadow $r(z) = \Res^{\coll}_{0,2}(\Theta_\cA)$.
\item $\Theta_\cA \in \MC(\gAmod)$ is the universal MC element
from the bar-intrinsic construction $\Theta_\cA := D_\cA - d_0$,
always MC because $D_\cA^2 = 0$.
\item $\nabla^{\hol}_{g,n} = d - \mathrm{Sh}_{g,n}(\Theta_\cA)$
is the modular shadow connection on
$\overline\cM_{g,n}$, flat by the MC equation projected to
$(g, n)$.
\end{enumerate}

\subsection*{12.2.\enspace Collision filtration}

The modular convolution $\gAmod$ carries a collision filtration
$F^0_{\coll} \supset F^1_{\coll} \supset \cdots$ indexed by
minimum pairwise distance on $\FM_n(\bC)$.  Depth $0$ is
free propagation; depth $1$ is first collision; depth $k$
is $k$-fold nested collision.  The spectral sequence has
$E_1^{p,*} = H^*(\gr^p_{\coll}\gAmod, d_0)$ with $d_1$
Arnold cancellations; it converges to $H^*(\gAmod)$, and
collapse is governed by shadow depth: class $\mathbf G$
collapses at $E_2$; class $\mathbf M$ does not collapse at
any finite page.

\subsection*{12.3.\enspace Four proved recovery theorems}

\medskip\noindent\emph{(R1) Collision residue is twisting.}
At depth $1$, arity $2$:
$\Res^{\coll}_{0,2}(\Theta_\cA) = \pi_\cA$, the universal
twisting morphism producing the dg-shifted Yangian $r_\cA(z)$.
The MC equation becomes $\partial\pi + \pi\star\pi = 0$.

\medskip\noindent\emph{(R2) CYBE from Arnold.}  The three
boundary divisors of $\overline M_{0,4}$ correspond to channels
$(12|34)$, $(13|24)$, $(14|23)$.  Evaluating the Arnold relation
on the Casimir $\Omega\otimes\Omega \in (\cA^!)^{\otimes 3}$
produces the three terms of the CYBE
$\sum [r_{ij}, r_{ik}] = 0$.

\medskip\noindent\emph{(R3) Affine shadow / KZ comparison.}
For $\cA = \widehat\fg_k$, on the affine comparison surface the
genus-$0$ arity-$n$ shadow one-form is the KZ connection form
\[
\mathrm{Sh}_{0,n}(\Theta_{\widehat\fg_k})
\;=\;
\sum_{i < j} r_{ij} \, d\log(z_i - z_j)
\;=\;
\frac{1}{h^\vee + k}\sum_{i < j}
\frac{\Omega_{ij}\,dz_{ij}}{z_{ij}},
\]
where $r_{ij} = \Omega_{ij}/(h^\vee + k)$ is the
propagator-weighted Casimir.  Flatness is MC at genus $0$,
arity $n$.  For Virasoro, the Virasoro conformal Ward extraction
plus higher collision residues yields the BPZ equations.  The
affine Drinfeld--Kohno theorem compares monodromy with the
quantum-group representation category.

\medskip\noindent\emph{(R4) Quartic as $L_\infty$ obstruction.}
$[\mathfrak R^{\mod}_4(\cA)] = [o_4(\cA)]
\in H^2(\cA^{\mathrm{sh}}_{4,0})$ is the obstruction to extending
$r(z)$ from arity $2$ to arity $4$.  Vanishes for classes
$\mathbf G, \mathbf L, \mathbf C$; nonzero for class
$\mathbf M$ (Virasoro: $Q^{\contact} = 10/[c(5c+22)]$).

\subsection*{12.4.\enspace Five conjectural targets}

\begin{enumerate}[label=(\roman*)]
\item \emph{Boundary-defect realisation.}  $\cA_\partial(T)$
and $\cA^!_\partial(T)$ are the two algebras of $\cH(T)$ for
every 3d HT system, with $\cC \simeq \cA^!\text{-mod}$.
\item \emph{Yangian-shadow.}  In full derived generality, the
collision residue recovers the dg-shifted Yangian $r(z)$ at
the chain level, including higher $A_\infty$ homotopies.
\item \emph{Sphere reconstruction.}  The Gaiotto--Zeng (2026)
commuting differentials on the non-planar sphere algebra are
the scalar shadow $\mathrm{Sh}_{0,n}(\Theta_\cA)$; sphere
amplitudes are flat sections of $\nabla^{\hol}_{0,n}$ on
established comparison surfaces.
\item \emph{Quartic resonance obstruction.}
$\mathfrak R^{\mod}_4$ is the first global nonlinear invariant,
measuring the failure of the large-$N$ expansion to truncate
at tree level.
\item \emph{Singular-fibre descent.}
$\nabla^{\hol}|_{\partial\overline\cM} \sim \nabla^{\clutch}$
with residues from the clutching law of $\Theta_\cA$:
$\Theta_\cA|_{\sigma_\Gamma} = \prod_v \Theta_v \cdot \prod_e P_e$.
\end{enumerate}

\subsection*{12.5.\enspace Coloured modular deformation object}

A 3d HT theory produces three observation algebras:
$\Obs^{\hol}$ (bulk), $\Obs^\partial$ (boundary), $\Obs^\ell$
(line).  The coloured Swiss-cheese FM space
$\FM^{\log}_{\bullet, \SC}$ encodes the collision geometry of
all three.  The three-colour modular deformation object is
\[
\fg^{\mathrm{HT},\log}_{\mod}
\;=\;
\hom^{\alpha_{\mathrm{HT}}}\!\bigl(
C_\bullet(\FM^{\log}_{\bullet, \SC}),
\End_{\Ch_\infty}(\Obs^{\hol}_\infty, \Obs^\partial_\infty,
\Obs^\ell_\infty)\bigr).
\]
The three colours interact through $\SCchtop$: bulk-to-boundary,
boundary-to-line, no line-to-bulk.

The modular effective action is the graph sum
\[
S^{\mod}_{\mathrm{HT}}
\;=\;
\sum_{\Gamma \in \mathsf{StGraph}_{\SC}}
\frac{\hbar^{b_1(\Gamma)}}{|\Aut(\Gamma)|}
W^{\log}_\Gamma O_\Gamma,
\]
and the BV quantum master equation
$(d_{\BV} + \hbar\Delta_{\mathrm{odd}})\exp(S^{\mod}/\hbar) = 0$
is the single equation governing $T$.  At genus $0$: classical
BV $\{S_0, S_0\} = 0$; genus $1$: one-loop anomaly cancellation;
genus $g$: $g$-loop Ward identity.

\subsection*{12.6.\enspace Khan--Zeng bridge}

Khan and Zeng construct a 3d HT Poisson sigma model on
$\bR_+ \times \bC$ attached to a freely generated PVA:
$\fg$-valued $(0,1)$-form $A_{\bar z}$, $\fg^*$-valued scalar
$\phi$, $\fg$-valued $1$-form $\eta$ on $\bR_+$, with action
\[
S_{\KZ}
\;=\;
\int_{\bR_+\times\bC}
(\langle\phi, \bar\partial A\rangle
+ \langle\phi, \tfrac12\{A_\lambda A\}|_{\lambda=0}\rangle
+ \langle\eta, d_t\phi\rangle)\,dt\wedge dz\wedge d\bar z.
\]
Gauge invariance is the PVA Jacobi identity; a Virasoro
element promotes HT to topological; the boundary carries a
$(-1)$-shifted Poisson bracket and generates the genus-$0$
quantisation problem.  The modular extension couples the
bulk theory to stable curves at $g \ge 1$ and controls the
full genus tower.  The pipeline:
\[
\text{classical coisson}
\;\xrightarrow{\text{KZ quantisation}}\;
\text{genus-$0$ quantum Koszul}
\;\xrightarrow{\text{modular lift}}\;
\cH(T).
\]

The bridge is completed in Volume~II by the doubling theorem:
the half-space BV problem for any logarithmic $\SCchtop$-algebra
embeds into the doubled whole-space via method of images; the
reflected weight identity is the $\sigma$-invariant restriction
of the standard Stokes identity on the doubled FM compactification.
The affine Poisson--vertex sigma model satisfies the physics
bridge hypotheses, and Drinfeld--Sokolov reduction preserves
this structure.  The entire standard landscape enters through
these two gates.

The genus expansion of the modular bar differential,
$D = \sum_{g \ge 0} \hbar^g D^{(g)}$, is the Feynman transform
of the modular operad; the canonical MC element
$\delta_\cA = D - D^{(0)}$ computes positive-genus corrections.
Stable graphs, automorphism weights, genus bookkeeping: the
mathematical shadow of the 't~Hooft expansion.  The
identification $\hbar \leftrightarrow 1/N$ for a specific
holographic model remains the principal open bridge.

\subsection*{12.7.\enspace Heisenberg holographic datum}

Base case: shadow class $\mathbf G$, shadow depth $2$, trivial
line-operator category, scalar $R$-matrix.  For $\cA = \cH_k$:
\begin{align*}
\cA &= \cH_k, & \cA^! &= \Sym^{\ch}(V^*),
  & \cC &= \mathsf{Vect},\\
r(z) &= k/z, & R(z) &= \exp(k/z),
  & \Theta_{\cH_k} &= k\cdot\eta\otimes\Lambda.
\end{align*}
The shadow connection $\nabla^{\hol}_{g,n} = d$ is trivial
(scalar MC element).  All five conjectural targets are
trivially verified: free boundary, curved symmetric dual
(with $\cH_k^! \neq \cH_{-k}$, $m_0 = -k\omega$), trivial
lines, vanishing quartic, no moduli dependence.

\subsection*{12.8.\enspace Affine $\widehat{\mathfrak{sl}}_2$ holographic datum}

First nontrivial example: class $\mathbf L$, shadow depth $3$,
Yang's $r(z) = k\,\Omega/z$, line-operator category
$Y(\mathfrak{sl}_2)\text{-mod}$, KZ shadow connection,
quantum-group representation theory from genus-$0$ monodromy.
\[
\cH(\widehat{\mathfrak{sl}}_2{}_k)
\;=\;
\Bigl(
\widehat{\mathfrak{sl}}_2{}_k,\;
\widehat{\mathfrak{sl}}_2{}_{-k-4},\;
\cO_q^{\mathrm{sh}},\;
\tfrac{k\,\Omega}{z},\;
\Theta_k,\;
d - \kappa_k\,\omega_g
\Bigr),
\]
with $q = e^{i\pi/(k+2)}$ and $\kappa_k = 3(k+2)/4$.  At
$k = -2$: $\kappa = 0$, the Feigin--Frenkel centre is the
algebra of $\mathfrak{sl}_2$-opers.  The KZ monodromy is
compared with the $U_q(\mathfrak{sl}_2)$ $R$-matrix by affine
Drinfeld--Kohno; on the reduced evaluation comparison surface,
$\cC_{\mathrm{line}}^{\red}|_{\mathrm{eval}}$ is identified
with the corresponding quantum-group side, and the coloured
Jones polynomial is recovered from monodromy of the bar complex
(Steinberg convolution of the chiral self-intersection on
$\FM_k(\bC)$).

This is 3d Chern--Simons at level $k$.  The holographic datum
is the $k = 0$-vanishing $r$-matrix plus the full KZ package.

\subsection*{12.9.\enspace Virasoro holographic datum}

First example of shadow class $\mathbf M$.  The $r$-matrix
$r_c(z) = (c/2)/z^3 + 2T/z$ has both a third-order pole
(gravitational) and a first-order pole (stress-tensor
exchange); the quartic contact invariant
$Q^{\contact} = 10/[c(5c+22)] \neq 0$.  The datum packages
3d gravity on $\bR_+ \times \bC$:
\begin{align*}
\cH(\Vir_c)
&\;=\;
(\Vir_c,\; \Vir_{26-c},\; \cC_c,\; r_c(z),\; \Theta_c,\;
\nabla^{\hol}_c),\\
r_c(z) &\;=\; \frac{c/2}{z^3} + \frac{2T}{z},
\qquad
\kappa_c \;=\; \frac{c}{2},\\
\nabla^{\hol}_{1,1}
&\;=\; d \;-\; \kappa_c\,\omega_1
\;-\; \frac{10}{c(5c+22)}\,\delta_4 \;-\; \cdots.
\end{align*}
Koszul dual $\Vir_{26-c}$ (self-dual at $c = 13$, not $c = 26$);
shadow connection with curvature $\kappa_c\omega_g$ at leading
order.  The transferred coproduct $\Delta_z^{\Vir}$ is strictly
primitive at all arities
(Principle~\ref{princ:gravitational-primitivity}): ghost-number
obstruction intrinsic to DS reduction kills every HPL tree
correction.  All gravitational dynamics resides in the
collision residue and the infinite $A_\infty$ product tower;
the coproduct channel is identically zero.  Gravity does not
produce particles, only braids them.

\subsection*{12.10.\enspace The M2 brane example}

First genuinely interacting holographic example: boundary
algebra $\cW_{1+\infty}$ (infinitely many generators), Koszul
dual dg-shifted Yangian for $\mathfrak{psl}(2|2)$, line-operator
category a quantum affine superalgebra.  All three genus-$2$
shells active; the planted-forest shell receives its first
nontrivial contribution from Mok's log-FM strata.

Physical setup: 3d $\cN = 4$ ADHM gauge theory, $N$ M2 branes
probing $\bC^2/\bZ_k$.  At large $N$, the boundary VOA is
$\cW_{1+\infty}[c]$ at $c \sim N$, with strong generators of
weights $1, 2, 3, \ldots$ and rational OPE coefficients in $N$.

The Koszul dual is
$\cA^!_{M2,\infty} \simeq Y^{\dg}_\hbar(\mathfrak{psl}(2|2))$.
The classical $r$-matrix $k\,\Omega_{\mathfrak{psl}(2|2)}/z$
vanishes at $k = 0$ and satisfies CYBE via the Arnold relation
on the $\mathfrak{psl}(2|2)$ Casimir.  The line-operator
category is
$\cC_{M2} \simeq \Rep(U_q(\widehat{\mathfrak{psl}}(2|2)))$
with $q = e^{2\pi i/(k+2)}$.

The datum:
\begin{align*}
\cA &= \cA_{M2,\infty},
  & \cA^! &= Y^{\dg}_\hbar(\mathfrak{psl}(2|2)),\\
\cC &= \Rep(U_q(\widehat{\mathfrak{psl}}(2|2))),
  & r(z) &= k\,\Omega_{\mathfrak{psl}(2|2)}/z,\\
\Theta_{M2} &= \sum_\Gamma
  \tfrac{W^{\log\FM}_\Gamma}{|\Aut(\Gamma)|}\otimes\Phi^{M2}_\Gamma,
  & \nabla^{\hol}_{g,n} &= d - \mathrm{Sh}_{g,n}(\Theta_{M2}).
\end{align*}
Shadow tower: cubic nonzero (nonabelian current subalgebra);
quartic nonzero ($16/(22 + 5c(N))$); quintic
$\{\mathfrak C_{M2}, Q^{\contact}_{M2}\}_H \neq 0$; tower
infinite.  At genus $2$, all three shells active; the
planted-forest shell is the first contribution from Mok's
log-FM strata to the holographic genus expansion.  The quartic
contact invariant is the first obstruction to planarity: it
distinguishes the full string-theoretic amplitude from its
planar limit.  The modular Koszul datum $\cH(M2)$ packages the
complete holographic correspondence into a single MC problem.

\bigskip\noindent\emph{Summary.}\enspace Sections~1--2
constructed the bar complex from $\eta = d\log(z_1 - z_2)$ and
the Arnold relation.  Sections~3--4 identified these as
projections of $\Theta_\cA \in \MC(\gAmod)$.  Sections~5--6
extracted the characteristic invariants from $\Theta_\cA$ and
classified algebras by shadow depth.  Section~7 ran the
machine on every standard family, reading both the geometric
face and the $E_1$ algebraic face for each family.  Section~8
lifted the shadow obstruction tower into number theory.
Section~9 proved scalar saturation and the Koszulness
meta-theorem.  Sections~10--11 transported the machine to
three dimensions via the Swiss-cheese operad and the modular
PVA quantisation obstruction theory.  Section~12 packaged
the full holographic correspondence into a six-component
modular Koszul datum controlled by a single MC equation.


% ====================================================================
\section*{13.\quad Completion and the frontier}
% ====================================================================

The holographic datum of Section~12 packages a chiral algebra,
its Koszul dual, the line-operator category, the spectral
$R$-matrix, the universal MC element, and the shadow connection
into a single algebraic object.  The frontier consists of
structural questions the packaging forces: the completed
bar-cobar problem by infinite-type algebras, the analytic
sewing problem by $g \ge 1$ curvature, the factorization
envelope problem by the passage from Lie conformal input to
modular output, and the non-principal duality problem by
Drinfeld--Sokolov reduction at arbitrary nilpotent orbits.

\subsection*{13.1.\enspace The Maurer--Cartan frontier}

\noindent\textsc{Status: master conjectures MC1 through MC4 are
proved; MC5 is partially proved, with the analytic HS-sewing package
established at all genera and the genuswise BV/BRST/bar identification
conjectural.}  MC3 holds for all simple types on the
evaluation-generated core via multiplicity-free $\ell$-weights
(Theorem~\ref{thm:categorical-cg-all-types},
Corollary~\ref{cor:mc3-all-types}).  The residual problem
DK-4/5 (extension beyond evaluation modules) is downstream.

\medskip\noindent\textbf{MC1} (PBW concentration).  For every
standard family (Kac--Moody, Virasoro, principal $\cW_N$,
$\beta\gamma$, lattice) and every genus $g \ge 1$, the PBW
spectral sequence on the genus-$g$ bar complex has
$E_\infty^{p,q}(g) = 0$ for $q \neq 0$.  Proof: genus-$0$
Koszulness gives concentration of $E_1^{p,q}(g{=}0)$; the
genus-$g$ enrichment from the Hodge bundle is
$d_0^{\mathrm{PBW}}$-exact because it arises from the regular
Eisenstein series (exact under Poincar\'e residue).
Concentration is uniform in $k$.

\medskip\noindent\textbf{MC2} (universal $\Theta_\cA$).  The
bar-intrinsic construction $\Theta_\cA := D_\cA - d_0
\in \gAmod$ is automatically MC because $D_\cA^2 = 0$.  Scalar
saturation universality: for every one-parameter algebraic
family $\cA_k$ with rational OPE coefficients,
$\Theta_{\cA_k} = \kappa(\cA_k)\cdot\eta_{\mathrm{univ}}
\otimes\Lambda_{\mathrm{univ}}$.  Proof via Whitehead
reduction ($E_2^{p,q} = H^p(\fg; \Sym^q V) = 0$ at $p = 2$ by
Whitehead's lemma for reductive $\fg$) and algebraic
semicontinuity (the non-vanishing locus is Zariski-closed,
proper, and empty at the free-field point).  Ambient
$D^2 = 0$ on the planted-forest complex follows from Mok's
normal-crossings result.

\medskip\noindent\textbf{MC3} (thick generation).  On the
type-$A$ Yangian module surface, the four-package MC3 problem
reduces to a single compact/completed extension packet via
(a) chromatic filtration of $K_0$ by $q$-character support,
(b) prefundamental Clebsch--Gordan closure
$V_n \otimes L^- = \bigoplus_{j=0}^n L^-(n-2j)$ at character
level, (c) Francis--Gaitsgory pro-nilpotent completion, and
(d) the remaining DK compact/completion packet.  The CG
decomposition is extended to all simple types via the
Chari--Moura multiplicity-free $\ell$-weight property,
replacing the minuscule hypothesis; the remaining post-CG
completion packets beyond type~$A$ are the downstream problem.

\medskip\noindent\textbf{MC4} (completed bar-cobar).  The
strong completion-tower theorem: the strong filtration axiom
$\mu_r(F^{i_1},\ldots,F^{i_r}) \subset F^{i_1+\cdots+i_r}$
yields an automatic arity cutoff, so for fixed $N$ only
finitely many arities contribute to $F^N$, making continuity
and Mittag-Leffler automatic.  The completion closure
$\CompCl(\cF_{\mathrm{ft}})$ of complete filtered objects with
finite-type associated graded carries a quasi-inverse
bar-cobar homotopy equivalence.  Key results: coefficient
stability, MC twisting closure, completed twisting
representability, and the uniform PBW bridge (MC1 implies
MC4 for all standard families).

MC4 splits into two sub-problems, both proved:
\begin{itemize}
\item \emph{MC4${}^+$} (positive towers: $\cW_{1+\infty}$,
affine Yangians, RTT algebras): solved by weight stabilisation.
\item \emph{MC4${}^0$} (resonant towers: Virasoro,
non-quadratic $\cW_N$): reduced to a finite resonance problem
by the resonance-filtered bar-cobar theorem.  The resonance
rank $\rho(\cA)$ measures the depth; every positive-energy
chiral algebra has $\rho < \infty$
(Theorem~\ref{thm:platonic-completion}).  The Virasoro case:
$\Vir_{26-c}$ is the depth-zero resonance shadow; the genuine
$\cW_\infty$-type dual requires the full resonance-filtered
completion.
\end{itemize}

\medskip\noindent\textbf{MC5} (genus expansion).  Partially proved.
The analytic sewing package is established at all genera, together
with the algebraic genus tower.  The full genuswise BV/BRST/bar
identification remains conjectural.  At genus~$0$, the algebraic
BRST/bar comparison is proved
(Theorem~\ref{thm:algebraic-string-dictionary}); the tree-level
amplitude pairing is conditional on
Corollary~\ref{cor:string-amplitude-genus0}; bar-cobar inversion
$\Omega(\barB(\cA)) \xrightarrow{\sim} \cA$ is Theorem~B.  At
genus~$1$, the curved bar complex with
$d_{\mathrm{fib}}^2 = \kappa\cdot\omega_1$ supplies the
one-loop package, and the complementarity splitting
$Q_1(\cA) \oplus Q_1(\cA^!) = H^*(\overline\cM_{1,n}, Z(\cA))$
is the genus-$1$ modular-invariance constraint.

At $g \ge 2$ the resolution proceeds in five steps:
\begin{enumerate}[label=(\roman*),nosep]
\item \emph{Inductive genus determination}
(Theorem~\ref{thm:inductive-genus-determination}):
$\Theta^{(g)}_\cA$ is determined by $\Theta^{(<g)}_\cA$ via
the graph-sum recursion on $\overline\cM_{g,n}$.
\item \emph{2D convergence}: no UV renormalisation on curves;
the Feynman integrals
$W_\Gamma = \int_{\sigma_\Gamma}\omega_\Gamma$ converge
absolutely on each stratum without counterterm subtraction.
\item \emph{Analytic-algebraic comparison}
$D^{\mathrm{an}} = D^{\mathrm{alg}}$ on the algebraic core.
\item \emph{General HS-sewing}
(Theorem~\ref{thm:general-hs-sewing}): polynomial OPE growth
plus subexponential sector growth implies the
Hilbert--Schmidt sewing condition, yielding trace-class closed
amplitudes at every genus.
\item \emph{Heisenberg one-particle sewing} identifies the
Heisenberg sewing amplitude with a Fredholm determinant,
supplying the decisive first-case verification.
\end{enumerate}

At multi-weight $g \ge 2$ (ALL-WEIGHT), the genus expansion
admits an explicit cross-channel correction:
$F_g(\cA) = \kappa(\cA)\cdot\lambda_g^{\FP} + \delta F_g^{\cross}$,
proved for $\cW_3$ by graph-by-graph computation over all seven
stable graphs at $g = 2$ and all forty-two stable graphs at
$g = 3$.  The scalar formula of the uniform-weight expansion
\emph{fails} for multi-weight algebras at $g \ge 2$
(resolving Open Problem~\ref{op:multi-generator-universality}
negatively).  A universal closed form for the gravitational
cross-channel correction covers principal $\cW_N$:
$\delta F_2^{\grav}(\cW_N, c)
= (N-2)(N+3)/96 + (N-2)(3N^3 + 14N^2 + 22N + 33)/(24c)$,
exact for $\cW_3$ and a lower bound for $N \ge 4$.  The
correct receptacle at $g \ge 1$ remains the coderived category:
curvature forces coderived/contraderived passage.

\medskip\noindent\textbf{Completion programme: $q$-windows and
completion entropy.}  The bar-cobar quasi-isomorphism
$\Omega(\barB(\cA)) \xrightarrow{\sim} \cA$ is proved at each
finite-type stage $\cW_N$.  The passage to the inverse limit
$\cW_\infty = \varprojlim \cW_N$ requires completed bar-cobar;
the strong filtration axiom gives a quotientwise-finite
differential.  The completed bar complex decomposes into finite
windows indexed by reduced weight $\rho(a) = \mathrm{wt}(a) - 1$.

For Virasoro at $q = 2$: the window $K_2(\Vir)$ has basis
$\{sT, sT|sT\}$, and the bar differential on the two-dimensional
complex is $d(sT|sT) = s(T\cdot T)$.  The product $T\cdot T$
decomposes into the quasi-primary $\Lambda$ (weight $4$,
reduced weight $3 > 2$, projects to zero), $\partial^2 T$
(reduced weight $3 > 2$, projects to zero), and the central
term $(c/2)\cdot\mathbf 1$.  Only the central term survives:
\[
d(sT|sT)|_{K_2} \;=\; (c/2)\cdot s\mathbf 1.
\]
The modular characteristic $\kappa(\Vir_c) = c/2$ is the sole
surviving coefficient of the bar differential on the smallest
nontrivial window.

Window dimensions grow exponentially.  Define the completion
entropy $h_K(\cA) := \log(\rho_\cA^{-1})$, with $\rho_\cA$ the
radius of convergence of the primitive generating series from
the vacuum character.  The entropy measures kinematic
completion complexity and forms a monotone ladder
\[
h_K(\cW_2) \approx 0.567
< h_K(\cW_3) \approx 0.772
< h_K(\cW_4) \approx 0.835
< \cdots
< h_K(\cW_\infty) \approx 0.872,
\]
governed at the limit by the MacMahon plane-partition
generating function
$\chi_{\cW_\infty}(q) = \prod_{n \ge 2}(1 - q^n)^{-(n-1)}$.

\subsection*{13.2.\enspace The analytic frontier}

Curvature forces this direction.  At $g \ge 1$ the bar complex
satisfies $d_{\mathrm{fib}}^2 = \kappa\cdot\omega_g \neq 0$;
the ordinary derived category is empty, and the correct
receptacle is the coderived category, requiring
Hilbert-space-valued structures.

The bare VOA is a dense algebraic skeleton; the object for
partition functions is the sewing envelope $\cA^{\mathrm{sew}}$,
the Hausdorff completion in the topology of all-amplitude
seminorms $\|a\|_{\mathrm{amp}} := \sup_{v \in \cV}
\langle v, \pi(a) v\rangle$.  The HS-sewing condition:
\[
\sum_{a,b,c \ge 0} q^{a+b+c}\|m^c_{a,b}\|^2_{\mathrm{HS}}
\;<\;\infty
\]
for $|q| < 1$ implies closed amplitudes
$\Tr(\pi(a)\cdot q^{L_0})$ are trace-class and the sewing
operation converges absolutely.  An analytic Koszul pair is
one for which the analytic bar-cobar adjunction restricts to a
Quillen equivalence on sewing-complete categories.

Four targets (two proved, two conjectural):
\begin{enumerate}[label=(\alph*),nosep]
\item \emph{Heisenberg sewing} (decisive first case, proved):
sewing envelope of $\cH_k$ is
$\Sym\,A_2(D) \cong \mathrm{ind\text{-}Hilb}(\cH_k)$
(Moriwaki 2026).
\item \emph{Lattice sewing} (proved): sewing envelope of
$V_\Lambda$ is vertex-algebraic completion by lattice theta
functions.
\item \emph{Analytic realisation criterion} (conjectural):
HS-sewing on $\cA$ plus dual HS on $\cA^!$ should suffice.
\item \emph{Boundary bar duality} (conjectural):
$B^{\mathrm{an}}(\cA_\partial) \simeq
B^{\mathrm{top}}(\mathrm{Fact}^{\mathrm{bulk}}_T)$.
\end{enumerate}

\subsection*{13.3.\enspace Factorization envelopes and the modular Koszul datum}

The bar construction starts from a chiral algebra, but the
natural input in most geometric constructions is a Lie conformal
algebra.  The factorisation envelope is the functor producing a
chiral algebra from Lie conformal data; the programme asks
whether the entire modular Koszul package can be constructed
functorially at the Lie conformal level, without passing
through the vertex algebra.  The target is a universal modular
factorization envelope $U^{\mod}_X(L)$, left adjoint to
extracting primitive currents.

A cyclically admissible Lie conformal algebra $L$ (conformal
grading, compatible decreasing filtration, bounded poles,
invariant bilinear pairing) produces the modular Koszul datum
\[
\Pi_X(L)
\;=\;
\bigl(
\mathrm{Fact}_X(L),\;
\bar B_X(L),\;
\Theta_L,\;
\cL_L,\;
(V^{\mathrm{br}}_L, T^{\mathrm{br}}_L),\;
R_4^{\mod}(L)
\bigr),
\]
with each component functorial in $L$.  The target universal
envelope $U^{\mod}_X(L)$ satisfies the adjunction
\[
\Hom_{\mathsf{Fact}^{\mod}}(U^{\mod}_X(L), \cF)
\;\cong\;
\Hom_{\mathsf{LCA}^{\cyc}}(L, \mathrm{Prim}^{\mod}(\cF)),
\]
and $\Theta_L = \Theta^{\mathrm{env}}_{\le\infty}(L)$ is the
universal MC class.  Direct sums: $\kappa$, $T^{\mathrm{br}}$,
$\Delta$, $R_4^{\mod}$ decompose additively.  BRST reduction:
$\Theta_{\cW_N} = H^0_{Q_{\DS}}(\Theta_{V_k(\fg)\otimes F_{\gh}})$,
the shadow tower of the $W$-algebra is the DS reduction of
the full current-plus-ghost tower, not of the bare affine tower.

\subsection*{13.4.\enspace Non-principal $W$-algebra duality}

Drinfeld--Sokolov reduction connects affine landscape to
$W$-algebra landscape.  For principal $f$, the reduction commutes
with bar-cobar duality by Feigin--Frenkel, and the holographic
datum descends.  For non-principal $f$, this commutation is
unproved.  The theory demands the extension: non-principal
$\cW$-algebras include the physically central cases
($\cW_k(\mathfrak{sl}_N, f_{\mathrm{sub}})$ controls $\cN = 2$
superconformal theories) and the mathematically central ones
(Feigin--Tipunin at admissible levels).

The hook-type nilpotent orbit in type $A$ is the first
non-principal corridor, conditional on DS-bar compatibility
(Theorem~\ref{thm:hook-transport-corridor}).  For a partition
$\lambda = (n - p, 1^p)$ in $\mathfrak{sl}_N$,
\[
(\cW_k(\mathfrak{sl}_N, f_\lambda))^!
\;\simeq\;
\cW_{k^\vee_\lambda}(\mathfrak{sl}_N, f_{\lambda^t}),
\]
where $\lambda^t$ is the transpose and $k^\vee_\lambda$ is a
rational function of $k$.  Proof: for hook-type partitions, DS
commutes with bar-cobar by Fehily and
Creutzig--Linshaw--Nakatsuka--Sato (2023--2025).

The transport propagation lemma: hook-type seeds combined with
edge-compatibility propagate duality results along edges of
the nilpotent Hasse diagram in type~$A$.  The type-$A$
transport-to-transpose conjecture: for every partition
$\lambda$ of $N$, the duality formula holds.  The general
obstruction is proving that DS reduction commutes with
bar-cobar at the chain level, i.e.\ that the natural map
$\mathrm{DS}(\barB(\widehat\fg_k)) \to
\barB(\mathrm{DS}(\widehat\fg_k))$ is a quasi-isomorphism.

\subsection*{13.5.\enspace The boxed machine}

\[
\boxed{
\text{homotopy chiral input }\cA
\;+\;
\text{log-FM cutting cooperad }\cC^!_{\ch}
\;+\;
\text{modular convolution }L_\infty
\;\;\Longrightarrow\;\;
\Theta_\cA,\;
\cT^{\mod}_{\Theta_\cA},\;
\kappa,\;
\Delta_\cA,\;
\mathfrak R^{\mod}_4,\;\ldots
}
\]

The modular convolution $L_\infty$-algebra $\gAmod$ is the
single organising structure.  It carries a natural modular
$L_\infty$ structure from the Feynman transform of the modular
operad; $\Conv_{\mathrm{str}}(\cC^!_{\ch}, \cP^{\ch})$ is its
strict model.  The universal MC element
$\Theta_\cA \in \MC(\gAmod)$ is proved: it is the bar-intrinsic
extraction $\Theta_\cA := D_\cA - d_0$, automatically MC because
$D_\cA^2 = 0$.

The shadow obstruction tower consists of finite-order
projections:
\begin{align*}
\Theta_\cA^{\le 2}&:\qquad\kappa(\cA)
  &&\text{(modular characteristic, scalar curvature)},\\
\Theta_\cA^{\le 3}&:\qquad\mathfrak C(\cA)
  &&\text{(cubic shadow, Lie structure)},\\
\Theta_\cA^{\le 4}&:\qquad\mathfrak Q(\cA)
  &&\text{(quartic resonance class, contact)},\\
\Theta_\cA^{\le r}&:\qquad\mathrm{Sh}_r(\cA)
  &&\text{(arity-$r$ shadow)},\\
\Theta_\cA&:\qquad\varprojlim_r\Theta_\cA^{\le r}
  &&\text{(universal MC element)}.
\end{align*}
Theorems~A--D and Theorem~H are proved projections at the
scalar level $\kappa$.  The holographic datum
$\cH(T) = (\cA, \cA^!, \cC, r(z), \Theta_\cA, \nabla^{\hol})$
packages $\Theta_\cA$ and its projections into the data of a
3d HT field theory.

\begin{principle}[The four-test interface]
\label{princ:four-test-interface}
\index{four-test interface|textbf}%
The modular Koszul machine has a complete algebraic-geometric
interface with $\overline\cM_{g,n}$ consisting of four
independent proved tests:
\begin{enumerate}[label=\textup{(\roman*)},nosep]
\item $D_\cA^2 = 0$: the bar differential squares to zero at
all genera and arities (Theorems~\ref{thm:bar-modular-operad},
\ref{thm:quantum-diff-squares-zero}).
\item $\mathrm{obs}_g(\cA) = \kappa(\cA)\cdot\lambda_g$ for
uniform-weight algebras at all genera (UNIFORM-WEIGHT), and
unconditionally at genus~$1$ for all families; for multi-weight
algebras at $g \ge 2$ (ALL-WEIGHT), with cross-channel
correction $\delta F_g^{\cross}$
(Theorem~\ref{thm:modular-characteristic},
Theorem~\ref{thm:multi-weight-genus-expansion}).
\item $Q_g(\cA) + Q_g(\cA^!) = H^*(\overline\cM_g, Z(\cA))$:
complementarity assembles the Koszul pair into a Lagrangian
decomposition
(Theorem~\ref{thm:quantum-complementarity-main}).
\item Sewing: genus-$g$ amplitudes converge from lower-genus
data (Theorem~\ref{thm:general-hs-sewing}).
\end{enumerate}
Each test captures a distinct aspect: (i)~is purely algebraic
and requires no Koszulness; (ii)~identifies the tautological
class but does not involve $\cA^!$; (iii)~requires the Koszul
pair but is genus-by-genus; (iv)~is analytic and independent
of the Hodge-class identification.  Any three can hold with
the fourth failing.  Three regions lie outside: multi-weight
at $g \ge 2$ (ALL-WEIGHT, resolved with cross-channel
correction, Open
Problem~\ref{op:multi-generator-universality}), BV/BRST${}={}$bar
at higher genus (Conjecture~\ref{conj:master-bv-brst}), and
non-perturbative completion.
\end{principle}

\medskip\noindent\textbf{The holographic frontier.}  The
holographic datum $\cH(T)$ is a boundary object.  Its lift to a
genuine 3d bulk theory is controlled by the relative
holographic deformation complex
$\fh_T^\Theta := \fib(\fg_T^{\SC} \to \gAmod)_{\alpha_T}$:
lifts of $\Theta_\cA$ are MC elements of the twisted fibre.
The graded obstruction tower
$[\fo_r(\beta)] \in H^2(\gr^F_r \fh_T^\Theta)$
turns the frontier into cohomology calculus.  At depth~$4$,
the relative quartic obstruction
$\mathfrak R_4^{\mathrm{rel}}(T)$ is the first invariant
detecting when a boundary datum cannot arise from a planar
holographic realisation.  The Loop--Connes principle
$\beta_{\der}^*(\Delta_{\mathrm{open}}) = \Lambda_P$:
closed genus creation is the image of the open BV operator
under the derived-centre map; loops in the bulk are traces on
the boundary.

\medskip\noindent\textbf{Cyclic trace, ribbon structure, and $1/N$.}
A bare modular graph sum does not define a literal 't~Hooft
expansion (Lemma~\ref{lem:bare-graph-no-thooft}).  Cyclic trace
structure on the open $E_1$-chiral algebra resolves this:
cyclicity enables ribbon thickening
(Theorem~\ref{thm:cyclicity-ribbon}), and a matrix realisation
gives the exact weight $N^{\chi(\Sigma_\Gamma)}$
(Theorem~\ref{thm:exact-n-chi-weighting}).  The three levels of
language (plain modular $E_1$ genus expansion, cyclic traced
modular $E_1$ 't~Hooft expansion, ribbonised modular
Swiss-cheese holographic 't~Hooft language) are made precise
in the $E_1$ modular Koszul chapter
(Chapter~\ref{chap:e1-modular-koszul}).

\subsection*{13.6.\enspace Admissible levels and logarithmic CFT}

At admissible levels $k = -h^\vee + p/q$ (coprime $p \ge h^\vee$,
$q \ge 1$), the Zhu algebra $A(\cV_k(\fg))$ is semisimple but
the vertex algebra is not simple: there is a maximal proper
ideal $J_k$ with $L_k(\fg) = \cV_k(\fg)/J_k$.  The bar complex
of $\cV_k$ acquires a periodic CDG structure indexed by $q$-th
roots of unity: the curvature $m_0 = \kappa\cdot\omega_g$ factors
through $\Phi_q(\zeta)$, and the weight filtration has a finite
plateau at level $N = q$.  This is the algebraic door to
logarithmic CFT on the degenerate admissible lane.  On the
non-degenerate admissible lane the simple quotient is rational;
on the degenerate lane the simple quotient has a logarithmic
module category.  Weightwise finite-page character degeneration
holds on the non-degenerate lane; identification of the
resulting bar cohomology with logarithmic extensions is not
proved here.

\subsection*{13.7.\enspace Relative Feynman transform}

The modular bar coalgebra $\barB_{\mod}(\cA)$ is an algebra
over the relative Feynman transform
$\mathrm{FT}_{\mathsf{Com}_{\mod}/\SCchtop}$, the algebraic
skeleton shared by the factorisation approach (\S10.8) and the
operadic approach.  Factorisation on
$\Ran(\Sigma_g)\times\Ran(\bR)$ over $\overline\cM_g$ produces
it via families of factorisable $\cD$-modules; the operadic
structure on $\SCchtop_{\mod}$ produces it via formal
completion.

\emph{Recognition.}  The tree-level part $D_0$ is the
$\SCchtop$-coalgebra structure; the genus-raising part $D_1$
is the $\mathsf{Com}_{\mod}$-modular structure.

\emph{Homotopy-involutivity.}
$\mathrm{FT}_{\mathsf{Com}_{\mod}/\SCchtop}^{\;2} \simeq \id$.
Genus-$0$ involutivity from operadic homotopy-Koszulity;
closed-colour modular involutivity from the Feynman involution
(Theorem~\ref{thm:feynman-involution}); assembly via the
genus-filtration spectral sequence.

\emph{Three routes unified.}  Route~B (factorisation) is global
truth; Route~A (operadic) is local approximation; Route~C
(this section) is the categorical image.  The factorisation
structure projects to the relative Feynman transform by
forgetting $\cD$-module data; the operadic structure injects
by formal completion.

\subsection*{13.8.\enspace Computational infrastructure}

Every formula has been verified by an independent symbolic
computation layer comprising tens of thousands of tests across
several hundred modules, covering the shadow obstruction tower
through arity $32$, the standard landscape for all seven
families, the genus-$2$ and genus-$3$ planted-forest formulas,
the DS-shadow cascade (arities $2$--$8$ for $N = 2, \ldots, 5$),
the $\cW_3$ multi-generator shadow metric, the B\"ocherer bridge
at genus~$2$, the Pixton shadow bridge, the entanglement
programme, the non-simply-laced bar computations ($B_2$, $C_2$,
$G_2$, $F_4$), minimal model bar and fusion rules, the dressed
propagator engine (scalar complementarity, $R$-matrix
$R_1 = 1/12$), and the constrained Epstein zeta.  The
computation is adversarial: tests encode cross-family
consistency checks that cannot be fooled by hardcoded values.

Volume~III extends the engine to Calabi--Yau geometry: forty-two
CY/BKM compute engines with $5700+$ tests cover the
$K3 \times E$ programme (shadow-Siegel gap, BKM root
multiplicities, moonshine multiplier, Schottky barrier), toric
CY$_3$ CoHA integrality, and the automorphic correction as
shadow obstruction tower identification.

Volume~II develops the 3d theory in three new chapters: the
factorisation Swiss-cheese algebra (global truth on
$\Ran(\Sigma_g)\times\Ran(\bR)$), the modular Swiss-cheese
operad (local approximation), and the relative Feynman
transform (algebraic skeleton).  Together with the 3d gravity
chapter, the affine half-space BV theorem, and eleven
additional computational engines, these constitute the
Volume~II programme.

\subsection*{13.9.\enspace The single open problem}

Everything proved in this monograph is a finite-order projection
of the universal MC element
$\Theta_\cA \in \MC(\gAmod)$.  Everything open is a question
about its ambient:

\begin{center}
\fbox{\parbox{0.85\textwidth}{\centering
\emph{Characterise the modular convolution algebra $\gAmod$
as a geometric object.}
}}
\end{center}

\noindent
Precisely: the modular convolution algebra is the endomorphism
algebra of the Lagrangian embedding
$\cL_\cA \hookrightarrow \cM_{\mathrm{vac}}(\cA)$ in the
$(-2)$-shifted symplectic stack of vacua.  Volume~I computes
the Taylor jets of this embedding (the shadow obstruction
tower).  Volume~II identifies the local composition law (the
Swiss-cheese operad).  What remains is the global geometry of
the stack itself: the derived symplectic structure, the moduli
of Lagrangians, and the higher-categorical descent that would
make the construction independent of choices.

The approximately $150$ open conjectures in these volumes are
projections of this single question:
\begin{itemize}[nosep]
\item the $\cD$-module purity conjecture asks whether the
Lagrangian is pure (characterisation~(xii) of Koszulness);
\item the non-principal $W$-duality programme asks how the
Lagrangian transforms under DS reduction at arbitrary nilpotent
orbits;
\item the analytic completion programme asks when the algebraic
Lagrangian extends to a convergent one (the sewing envelope);
\item the factorization envelope programme asks for the universal
construction of the Lagrangian from Lie conformal data alone;
\item the holographic completion question asks whether the
genus-$0$ Lagrangian extends over $\overline\cM_{g,n}$;
\item the Calabi--Yau quantum group programme asks whether the
automorphic correction of a generalised Borcherds--Kac--Moody
superalgebra is the shadow obstruction tower of an
$E_2$-chiral algebra constructed from a CY category (proved
for $d = 2$; conditional on chain-level $\mathbb S^3$-framing
for $d = 3$; Volume~III).
\end{itemize}
Each is a different facet of the geometry of $\gAmod$, and
none requires going beyond the framework already established.

\bigskip

\noindent
The entire structure unfolds from a single datum: the
logarithmic $1$-form $\eta = d\log(z_1 - z_2)$.  At $n$ points
it became a differential on the Fulton--MacPherson
compactification; across all genera it became the total bar
differential $D_\cA$; in the modular convolution algebra it
became the Maurer--Cartan element $\Theta_\cA = D_\cA - d_0$.
One $1$-form, one relation (Arnold), one extraction (residues):
a universal invariant of chiral algebras that is algebraic,
computable, and complete.

\begin{thebibliography}{99}

\bibitem{BD04}
A.~Beilinson and V.~Drinfeld,
\textit{Chiral Algebras},
AMS Colloquium Publications, vol.~51,
American Mathematical Society, Providence, RI, 2004.

\bibitem{BGS96}
A.~Beilinson, V.~Ginzburg, and W.~Soergel,
Koszul duality patterns in representation theory,
\textit{J.~Amer.\ Math.\ Soc.}~\textbf{9} (1996), 473--527.

\bibitem{BS88}
A.~Beilinson and V.~Schechtman,
Determinant bundles and Virasoro algebras,
\textit{Comm.\ Math.\ Phys.}~\textbf{118} (1988), 651--701.

\bibitem{CG17}
K.~Costello and O.~Gwilliam,
\textit{Factorization Algebras in Quantum Field Theory},
vols.~1--2, Cambridge University Press, 2017/2021.

\bibitem{DS85}
V.~Drinfeld and V.~Sokolov,
Lie algebras and equations of Korteweg--de Vries type,
\textit{J.~Soviet Math.}~\textbf{30} (1985), 1975--2036.

\bibitem{FG12}
J.~Francis and D.~Gaitsgory,
Chiral Koszul duality,
\textit{Selecta Math.\ (N.S.)}~\textbf{18} (2012), 27--87.

\bibitem{FP00}
C.~Faber and R.~Pandharipande,
Hodge integrals and moduli spaces of curves,
in \textit{Moduli of Curves and Abelian Varieties},
Aspects Math.~E33, Vieweg, 2000, pp.~109--130.

\bibitem{GeK98}
E.~Getzler and M.~Kapranov,
Modular operads,
\textit{Compositio Math.}~\textbf{110} (1998), 65--126.

\bibitem{FF90}
B.~Feigin and E.~Frenkel,
Quantization of the Drinfeld--Sokolov reduction,
\textit{Phys.\ Lett.\ B}~\textbf{246} (1990), 75--81.

\bibitem{FlajoletSedgewick}
P.~Flajolet and R.~Sedgewick,
\textit{Analytic Combinatorics},
Cambridge University Press, 2009.

\bibitem{KM94}
M.~Kontsevich and Yu.~Manin,
Gromov--Witten classes, quantum cohomology, and enumerative geometry,
\textit{Comm.\ Math.\ Phys.}~\textbf{164} (1994), 525--562.

\bibitem{Kontsevich03}
M.~Kontsevich,
Deformation quantization of Poisson manifolds,
\textit{Lett.\ Math.\ Phys.}~\textbf{66} (2003), 157--216.

\bibitem{KRW03}
V.~Kac, S.-S.~Roan, and M.~Wakimoto,
Quantum reduction for affine superalgebras,
\textit{Comm.\ Math.\ Phys.}~\textbf{241} (2003), 307--342.

\bibitem{LV12}
J.-L.~Loday and B.~Vallette,
\textit{Algebraic Operads},
Grundlehren der Mathematischen Wissenschaften, vol.~346,
Springer, Berlin, 2012.

\bibitem{Mumford83}
D.~Mumford,
Towards an enumerative geometry of the moduli space of curves,
in \textit{Arithmetic and Geometry, Vol.~II},
Progr.\ Math., vol.~36, Birkh\"auser, 1983, pp.~271--328.

\bibitem{Priddy70}
S.~Priddy,
Koszul resolutions,
\textit{Trans.\ Amer.\ Math.\ Soc.}~\textbf{152} (1970), 39--60.

\bibitem{EK00}
P.~Etingof and D.~Kazhdan,
Quantization of Lie bialgebras V: quantum vertex operator algebras,
\textit{Selecta Math.\ (N.S.)}~\textbf{6} (2000), 105--130.

\bibitem{Mok25}
C.-P.~Mok,
\emph{Logarithmic Fulton--MacPherson configuration spaces},
arXiv:2503.17563, 2025.

\bibitem{MS24}
A.~Moreno and P.~Safronov,
Chiral Koszul duality on moduli of curves,
preprint, 2024.

\bibitem{RNW19}
N.~Rozenblyum, S.~Nagy, and D.~Wagner,
Convolution $L_\infty$-algebras and Maurer--Cartan,
preprint, 2019.

\bibitem{LorgatVirR}
R.~Lorgat,
\textit{The Virasoro spectral $R$-matrix: closed form and non-formality},
preprint, 2026.

\bibitem{LorgatSevenFaces}
R.~Lorgat,
\textit{The seven faces of the collision residue},
preprint, 2026.

\end{thebibliography}

\end{document}
